\section*{Lecture 4 - Dynamism}
\addcontentsline{toc}{section}{Lecture 4 - Dynamism}

\clearpage
\SlideHeader{Lecture 4 - Dynamism}{001}{Data Driven Decision Making}
\begin{minipage}[t]{0.405\textwidth}
\SlideImage{1}{../Materials/2026/3DM_04_Dynamism_SS26.pdf}
\begin{adjustbox}{max totalsize={\textwidth}{0.43\textheight},center}
\begin{minipage}{\textwidth}\scriptsize
\NoteSection{日本語訳}\begin{itemize}
\item データ駆動型意思決定
\item 講義 4 – 動的性質
\end{itemize}\NoteSection{試験対策ポイント}\begin{itemize}
\item dynamismは、時間の進行により状態が変わり、意思決定を繰り返す性質である。
\item static planning, dynamic planning, a priori optimization, online reoptimizationを区別する。
\end{itemize}
\end{minipage}
\end{adjustbox}
\end{minipage}\hfill
\begin{minipage}[t]{0.58\textwidth}
\begin{adjustbox}{max totalsize={\textwidth}{0.89\textheight},center}
\begin{minipage}{\textwidth}\scriptsize
\NoteSection{解説}このページでは「Data Driven Decision Making」を理解する。スライド上の表現は短いが、試験では定義、具体例、モデル上の意味、手法の限界まで説明できる必要がある。\par
Dynamism は、意思決定問題が時間とともに進行し、状態が変化することを意味する。静的問題では、すべての入力を最初に与えられ、その情報に基づいて一度だけ計画を作る。動的問題では、新しい注文、交通状況、車両位置、キャンセル、故障などが途中で発生し、計画を更新する必要がある。\par
重要なのは、dynamism は stochasticity と同じではないという点である。stochasticity は未来が不確実であること、dynamism は時間とともに決定と状態が進む構造を持つことである。例えば、すべての未来注文が既知でも、時点ごとに車両位置が変わり決定を繰り返すなら動的である。逆に、将来の需要が確率的でも、最初に一括で決めて変更しないなら静的な確率計画として扱える。\par
a priori optimization は、実行前に全体計画を作り、実行中は大きく変更しない考え方である。計算しやすく、全体最適を考えやすいが、予期しない情報に弱い。online または dynamic planning は、各時点で現在状態を観測し、必要に応じて決定を変える。現実に適応しやすいが、計算時間や一貫性が問題になる。\par
配送の例では、朝に全ルートを決めるのが静的計画である。昼に新規注文が入ったり、車両が遅れたりしたときに、現在の車両位置と残り注文を見て再計画するのが動的計画である。動的計画では、今の選択が将来の柔軟性を奪わないかも考える必要がある。\par
試験答案では、静的/動的の違いを『計画を変更するか』だけでなく、状態、時点、決定、外生情報、遷移の連鎖として説明する。特に動的問題では、状態が更新され、次の決定の入力になるというループを明確に書くと満点答案に近づく。\par\NoteSection{確認問題と解答}\begin{enumerate}\item \textbf{L04-S001: 「Data Driven Decision Making」の概念を、定義・具体例・モデル上の意味の順に説明せよ。}\\定義としては、dynamismは、時間の進行により状態が変わり、意思決定を繰り返す性質である。。具体例では、配送・在庫・フードデリバリーのような現実問題を用い、状態、決定、外生情報、報酬、または情報モデル/意思決定モデルのどこに対応するかを明示する。モデル上の意味として、この概念が目的関数、制約、状態表現、将来価値、または方策選択にどのような影響を与えるかを書く。試験では、用語だけを列挙せず、入力データから意思決定、さらに現実の実装へつながる流れを文章で示すことが重要である。\item \textbf{L04-S001: このスライドの考え方を、講義全体のDDDMプロセスの中に位置づけよ。}\\DDDMプロセスでは、現実世界のデータを集約して情報モデルを作り、意思決定モデルまたは方策で行動を選び、実装後に結果を観測して次の状態へ進む。このスライドの概念は、そのどこか一部だけではなく、データ表現、意思決定、将来不確実性、計算可能性のいずれかに関わる。答案では、まずこの概念がどの段階に属するかを述べ、次に他の段階へどう影響するかを書く。例えば価値関数なら将来価値を現在決定へ戻す役割、FIFOなら旅行時間情報モデルの整合性を保証する役割、CFAなら目的関数を調整して将来に良い行動を誘導する役割である。\end{enumerate}
\end{minipage}
\end{adjustbox}
\end{minipage}


\clearpage
\SlideHeader{Lecture 4 - Dynamism}{002}{Recap and Motivation}
\begin{minipage}[t]{0.405\textwidth}
\SlideImage{2}{../Materials/2026/3DM_04_Dynamism_SS26.pdf}
\begin{adjustbox}{max totalsize={\textwidth}{0.43\textheight},center}
\begin{minipage}{\textwidth}\scriptsize
\NoteSection{日本語訳}\begin{itemize}
\item 前回の復習と今回の位置づけ
\item Uncertainty（この用語は講義スライド上の専門語として扱う）
\item Information modeling（この用語は講義スライド上の専門語として扱う）
\item 決定 modeling
\item Today: Can we react to observed changes in（この用語は講義スライド上の専門語として扱う）
\item the information model?（この用語は講義スライド上の専門語として扱う）
\item web.stanford.edu（この用語は講義スライド上の専門語として扱う）
\end{itemize}\NoteSection{試験対策ポイント}\begin{itemize}
\item dynamismは、時間の進行により状態が変わり、意思決定を繰り返す性質である。
\item static planning, dynamic planning, a priori optimization, online reoptimizationを区別する。
\end{itemize}
\end{minipage}
\end{adjustbox}
\end{minipage}\hfill
\begin{minipage}[t]{0.58\textwidth}
\begin{adjustbox}{max totalsize={\textwidth}{0.89\textheight},center}
\begin{minipage}{\textwidth}\scriptsize
\NoteSection{解説}このページでは「Recap and Motivation」を理解する。スライド上の表現は短いが、試験では定義、具体例、モデル上の意味、手法の限界まで説明できる必要がある。\par
Dynamism は、意思決定問題が時間とともに進行し、状態が変化することを意味する。静的問題では、すべての入力を最初に与えられ、その情報に基づいて一度だけ計画を作る。動的問題では、新しい注文、交通状況、車両位置、キャンセル、故障などが途中で発生し、計画を更新する必要がある。\par
重要なのは、dynamism は stochasticity と同じではないという点である。stochasticity は未来が不確実であること、dynamism は時間とともに決定と状態が進む構造を持つことである。例えば、すべての未来注文が既知でも、時点ごとに車両位置が変わり決定を繰り返すなら動的である。逆に、将来の需要が確率的でも、最初に一括で決めて変更しないなら静的な確率計画として扱える。\par
a priori optimization は、実行前に全体計画を作り、実行中は大きく変更しない考え方である。計算しやすく、全体最適を考えやすいが、予期しない情報に弱い。online または dynamic planning は、各時点で現在状態を観測し、必要に応じて決定を変える。現実に適応しやすいが、計算時間や一貫性が問題になる。\par
配送の例では、朝に全ルートを決めるのが静的計画である。昼に新規注文が入ったり、車両が遅れたりしたときに、現在の車両位置と残り注文を見て再計画するのが動的計画である。動的計画では、今の選択が将来の柔軟性を奪わないかも考える必要がある。\par
試験答案では、静的/動的の違いを『計画を変更するか』だけでなく、状態、時点、決定、外生情報、遷移の連鎖として説明する。特に動的問題では、状態が更新され、次の決定の入力になるというループを明確に書くと満点答案に近づく。\par\NoteSection{確認問題と解答}\begin{enumerate}\item \textbf{L04-S002: 「Recap and Motivation」の概念を、定義・具体例・モデル上の意味の順に説明せよ。}\\定義としては、dynamismは、時間の進行により状態が変わり、意思決定を繰り返す性質である。。具体例では、配送・在庫・フードデリバリーのような現実問題を用い、状態、決定、外生情報、報酬、または情報モデル/意思決定モデルのどこに対応するかを明示する。モデル上の意味として、この概念が目的関数、制約、状態表現、将来価値、または方策選択にどのような影響を与えるかを書く。試験では、用語だけを列挙せず、入力データから意思決定、さらに現実の実装へつながる流れを文章で示すことが重要である。\item \textbf{L04-S002: このスライドの考え方を、講義全体のDDDMプロセスの中に位置づけよ。}\\DDDMプロセスでは、現実世界のデータを集約して情報モデルを作り、意思決定モデルまたは方策で行動を選び、実装後に結果を観測して次の状態へ進む。このスライドの概念は、そのどこか一部だけではなく、データ表現、意思決定、将来不確実性、計算可能性のいずれかに関わる。答案では、まずこの概念がどの段階に属するかを述べ、次に他の段階へどう影響するかを書く。例えば価値関数なら将来価値を現在決定へ戻す役割、FIFOなら旅行時間情報モデルの整合性を保証する役割、CFAなら目的関数を調整して将来に良い行動を誘導する役割である。\end{enumerate}
\end{minipage}
\end{adjustbox}
\end{minipage}


\clearpage
\SlideHeader{Lecture 4 - Dynamism}{003}{Goals Today}
\begin{minipage}[t]{0.405\textwidth}
\SlideImage{3}{../Materials/2026/3DM_04_Dynamism_SS26.pdf}
\begin{adjustbox}{max totalsize={\textwidth}{0.43\textheight},center}
\begin{minipage}{\textwidth}\scriptsize
\NoteSection{日本語訳}\begin{itemize}
\item 目標 Today
\item 動的性質 in information and 決定 model
\item Dynamic vs. Static 計画
\item Dynamic 決定 process
\item Stochastic and Dynamic Vehicle Routing（この用語は講義スライド上の専門語として扱う）
\item First steps of modelling（この用語は講義スライド上の専門語として扱う）
\end{itemize}\NoteSection{試験対策ポイント}\begin{itemize}
\item dynamismは、時間の進行により状態が変わり、意思決定を繰り返す性質である。
\item static planning, dynamic planning, a priori optimization, online reoptimizationを区別する。
\end{itemize}
\end{minipage}
\end{adjustbox}
\end{minipage}\hfill
\begin{minipage}[t]{0.58\textwidth}
\begin{adjustbox}{max totalsize={\textwidth}{0.89\textheight},center}
\begin{minipage}{\textwidth}\scriptsize
\NoteSection{解説}このページでは「Goals Today」を理解する。スライド上の表現は短いが、試験では定義、具体例、モデル上の意味、手法の限界まで説明できる必要がある。\par
Dynamism は、意思決定問題が時間とともに進行し、状態が変化することを意味する。静的問題では、すべての入力を最初に与えられ、その情報に基づいて一度だけ計画を作る。動的問題では、新しい注文、交通状況、車両位置、キャンセル、故障などが途中で発生し、計画を更新する必要がある。\par
重要なのは、dynamism は stochasticity と同じではないという点である。stochasticity は未来が不確実であること、dynamism は時間とともに決定と状態が進む構造を持つことである。例えば、すべての未来注文が既知でも、時点ごとに車両位置が変わり決定を繰り返すなら動的である。逆に、将来の需要が確率的でも、最初に一括で決めて変更しないなら静的な確率計画として扱える。\par
a priori optimization は、実行前に全体計画を作り、実行中は大きく変更しない考え方である。計算しやすく、全体最適を考えやすいが、予期しない情報に弱い。online または dynamic planning は、各時点で現在状態を観測し、必要に応じて決定を変える。現実に適応しやすいが、計算時間や一貫性が問題になる。\par
配送の例では、朝に全ルートを決めるのが静的計画である。昼に新規注文が入ったり、車両が遅れたりしたときに、現在の車両位置と残り注文を見て再計画するのが動的計画である。動的計画では、今の選択が将来の柔軟性を奪わないかも考える必要がある。\par
試験答案では、静的/動的の違いを『計画を変更するか』だけでなく、状態、時点、決定、外生情報、遷移の連鎖として説明する。特に動的問題では、状態が更新され、次の決定の入力になるというループを明確に書くと満点答案に近づく。\par\NoteSection{確認問題と解答}\begin{enumerate}\item \textbf{L04-S003: 「Goals Today」の概念を、定義・具体例・モデル上の意味の順に説明せよ。}\\定義としては、dynamismは、時間の進行により状態が変わり、意思決定を繰り返す性質である。。具体例では、配送・在庫・フードデリバリーのような現実問題を用い、状態、決定、外生情報、報酬、または情報モデル/意思決定モデルのどこに対応するかを明示する。モデル上の意味として、この概念が目的関数、制約、状態表現、将来価値、または方策選択にどのような影響を与えるかを書く。試験では、用語だけを列挙せず、入力データから意思決定、さらに現実の実装へつながる流れを文章で示すことが重要である。\item \textbf{L04-S003: このスライドの考え方を、講義全体のDDDMプロセスの中に位置づけよ。}\\DDDMプロセスでは、現実世界のデータを集約して情報モデルを作り、意思決定モデルまたは方策で行動を選び、実装後に結果を観測して次の状態へ進む。このスライドの概念は、そのどこか一部だけではなく、データ表現、意思決定、将来不確実性、計算可能性のいずれかに関わる。答案では、まずこの概念がどの段階に属するかを述べ、次に他の段階へどう影響するかを書く。例えば価値関数なら将来価値を現在決定へ戻す役割、FIFOなら旅行時間情報モデルの整合性を保証する役割、CFAなら目的関数を調整して将来に良い行動を誘導する役割である。\end{enumerate}
\end{minipage}
\end{adjustbox}
\end{minipage}


\clearpage
\SlideHeader{Lecture 4 - Dynamism}{004}{Agenda}
\begin{minipage}[t]{0.405\textwidth}
\SlideImage{4}{../Materials/2026/3DM_04_Dynamism_SS26.pdf}
\begin{adjustbox}{max totalsize={\textwidth}{0.43\textheight},center}
\begin{minipage}{\textwidth}\scriptsize
\NoteSection{日本語訳}\begin{itemize}
\item アジェンダ
\item 動的性質
\item ケーススタディ: Urban Delivery and Traffic Management
\item Preparation of Modelling（この用語は講義スライド上の専門語として扱う）
\end{itemize}\NoteSection{試験対策ポイント}\begin{itemize}
\item dynamismは、時間の進行により状態が変わり、意思決定を繰り返す性質である。
\item static planning, dynamic planning, a priori optimization, online reoptimizationを区別する。
\end{itemize}
\end{minipage}
\end{adjustbox}
\end{minipage}\hfill
\begin{minipage}[t]{0.58\textwidth}
\begin{adjustbox}{max totalsize={\textwidth}{0.89\textheight},center}
\begin{minipage}{\textwidth}\scriptsize
\NoteSection{解説}このページでは「Agenda」を理解する。スライド上の表現は短いが、試験では定義、具体例、モデル上の意味、手法の限界まで説明できる必要がある。\par
Dynamism は、意思決定問題が時間とともに進行し、状態が変化することを意味する。静的問題では、すべての入力を最初に与えられ、その情報に基づいて一度だけ計画を作る。動的問題では、新しい注文、交通状況、車両位置、キャンセル、故障などが途中で発生し、計画を更新する必要がある。\par
重要なのは、dynamism は stochasticity と同じではないという点である。stochasticity は未来が不確実であること、dynamism は時間とともに決定と状態が進む構造を持つことである。例えば、すべての未来注文が既知でも、時点ごとに車両位置が変わり決定を繰り返すなら動的である。逆に、将来の需要が確率的でも、最初に一括で決めて変更しないなら静的な確率計画として扱える。\par
a priori optimization は、実行前に全体計画を作り、実行中は大きく変更しない考え方である。計算しやすく、全体最適を考えやすいが、予期しない情報に弱い。online または dynamic planning は、各時点で現在状態を観測し、必要に応じて決定を変える。現実に適応しやすいが、計算時間や一貫性が問題になる。\par
配送の例では、朝に全ルートを決めるのが静的計画である。昼に新規注文が入ったり、車両が遅れたりしたときに、現在の車両位置と残り注文を見て再計画するのが動的計画である。動的計画では、今の選択が将来の柔軟性を奪わないかも考える必要がある。\par
試験答案では、静的/動的の違いを『計画を変更するか』だけでなく、状態、時点、決定、外生情報、遷移の連鎖として説明する。特に動的問題では、状態が更新され、次の決定の入力になるというループを明確に書くと満点答案に近づく。\par\NoteSection{確認問題と解答}\begin{enumerate}\item \textbf{L04-S004: 「Agenda」の概念を、定義・具体例・モデル上の意味の順に説明せよ。}\\定義としては、dynamismは、時間の進行により状態が変わり、意思決定を繰り返す性質である。。具体例では、配送・在庫・フードデリバリーのような現実問題を用い、状態、決定、外生情報、報酬、または情報モデル/意思決定モデルのどこに対応するかを明示する。モデル上の意味として、この概念が目的関数、制約、状態表現、将来価値、または方策選択にどのような影響を与えるかを書く。試験では、用語だけを列挙せず、入力データから意思決定、さらに現実の実装へつながる流れを文章で示すことが重要である。\item \textbf{L04-S004: このスライドの考え方を、講義全体のDDDMプロセスの中に位置づけよ。}\\DDDMプロセスでは、現実世界のデータを集約して情報モデルを作り、意思決定モデルまたは方策で行動を選び、実装後に結果を観測して次の状態へ進む。このスライドの概念は、そのどこか一部だけではなく、データ表現、意思決定、将来不確実性、計算可能性のいずれかに関わる。答案では、まずこの概念がどの段階に属するかを述べ、次に他の段階へどう影響するかを書く。例えば価値関数なら将来価値を現在決定へ戻す役割、FIFOなら旅行時間情報モデルの整合性を保証する役割、CFAなら目的関数を調整して将来に良い行動を誘導する役割である。\end{enumerate}
\end{minipage}
\end{adjustbox}
\end{minipage}


\clearpage
\SlideHeader{Lecture 4 - Dynamism}{005}{Dynamism}
\begin{minipage}[t]{0.405\textwidth}
\SlideImage{5}{../Materials/2026/3DM_04_Dynamism_SS26.pdf}
\begin{adjustbox}{max totalsize={\textwidth}{0.43\textheight},center}
\begin{minipage}{\textwidth}\scriptsize
\NoteSection{日本語訳}\begin{itemize}
\item 動的性質
\item Dynamic problems require not only one but several subsequent 決定s.
\item The 決定s impact the future and are therefore interconnected.
\item There is new information revealed after each 決定 was made.
\item The solution for a 動的 problem is a strategy (方策)
\item It defines how to react to newly revealed information.（この用語は講義スライド上の専門語として扱う）
\item 報酬 or cost accumulate over time and 決定s.
\end{itemize}\NoteSection{試験対策ポイント}\begin{itemize}
\item dynamismは、時間の進行により状態が変わり、意思決定を繰り返す性質である。
\item static planning, dynamic planning, a priori optimization, online reoptimizationを区別する。
\end{itemize}
\end{minipage}
\end{adjustbox}
\end{minipage}\hfill
\begin{minipage}[t]{0.58\textwidth}
\begin{adjustbox}{max totalsize={\textwidth}{0.89\textheight},center}
\begin{minipage}{\textwidth}\scriptsize
\NoteSection{解説}このページでは「Dynamism」を理解する。スライド上の表現は短いが、試験では定義、具体例、モデル上の意味、手法の限界まで説明できる必要がある。\par
Dynamism は、意思決定問題が時間とともに進行し、状態が変化することを意味する。静的問題では、すべての入力を最初に与えられ、その情報に基づいて一度だけ計画を作る。動的問題では、新しい注文、交通状況、車両位置、キャンセル、故障などが途中で発生し、計画を更新する必要がある。\par
重要なのは、dynamism は stochasticity と同じではないという点である。stochasticity は未来が不確実であること、dynamism は時間とともに決定と状態が進む構造を持つことである。例えば、すべての未来注文が既知でも、時点ごとに車両位置が変わり決定を繰り返すなら動的である。逆に、将来の需要が確率的でも、最初に一括で決めて変更しないなら静的な確率計画として扱える。\par
a priori optimization は、実行前に全体計画を作り、実行中は大きく変更しない考え方である。計算しやすく、全体最適を考えやすいが、予期しない情報に弱い。online または dynamic planning は、各時点で現在状態を観測し、必要に応じて決定を変える。現実に適応しやすいが、計算時間や一貫性が問題になる。\par
配送の例では、朝に全ルートを決めるのが静的計画である。昼に新規注文が入ったり、車両が遅れたりしたときに、現在の車両位置と残り注文を見て再計画するのが動的計画である。動的計画では、今の選択が将来の柔軟性を奪わないかも考える必要がある。\par
試験答案では、静的/動的の違いを『計画を変更するか』だけでなく、状態、時点、決定、外生情報、遷移の連鎖として説明する。特に動的問題では、状態が更新され、次の決定の入力になるというループを明確に書くと満点答案に近づく。\par\NoteSection{確認問題と解答}\begin{enumerate}\item \textbf{L04-S005: 「Dynamism」の概念を、定義・具体例・モデル上の意味の順に説明せよ。}\\定義としては、dynamismは、時間の進行により状態が変わり、意思決定を繰り返す性質である。。具体例では、配送・在庫・フードデリバリーのような現実問題を用い、状態、決定、外生情報、報酬、または情報モデル/意思決定モデルのどこに対応するかを明示する。モデル上の意味として、この概念が目的関数、制約、状態表現、将来価値、または方策選択にどのような影響を与えるかを書く。試験では、用語だけを列挙せず、入力データから意思決定、さらに現実の実装へつながる流れを文章で示すことが重要である。\item \textbf{L04-S005: このスライドの考え方を、講義全体のDDDMプロセスの中に位置づけよ。}\\DDDMプロセスでは、現実世界のデータを集約して情報モデルを作り、意思決定モデルまたは方策で行動を選び、実装後に結果を観測して次の状態へ進む。このスライドの概念は、そのどこか一部だけではなく、データ表現、意思決定、将来不確実性、計算可能性のいずれかに関わる。答案では、まずこの概念がどの段階に属するかを述べ、次に他の段階へどう影響するかを書く。例えば価値関数なら将来価値を現在決定へ戻す役割、FIFOなら旅行時間情報モデルの整合性を保証する役割、CFAなら目的関数を調整して将来に良い行動を誘導する役割である。\end{enumerate}
\end{minipage}
\end{adjustbox}
\end{minipage}


\clearpage
\SlideHeader{Lecture 4 - Dynamism}{006}{Static and Dynamic Planning}
\begin{minipage}[t]{0.405\textwidth}
\SlideImage{6}{../Materials/2026/3DM_04_Dynamism_SS26.pdf}
\begin{adjustbox}{max totalsize={\textwidth}{0.43\textheight},center}
\begin{minipage}{\textwidth}\scriptsize
\NoteSection{日本語訳}\begin{itemize}
\item 静的計画と動的計画
\item 例: uncertain travel times
\item Static planning:（この用語は講義スライド上の専門語として扱う）
\item Replanning not possible（この用語は講義スライド上の専門語として扱う）
\item 実装 of a (reliable) a-priori plan www.braunschweiger-
\item 例: Bus
\item Dynamic planning:（この用語は講義スライド上の専門語として扱う）
\item High impact of 不確実性 requires replanning
\item Adaptations of plans（この用語は講義スライド上の専門語として扱う）
\end{itemize}\NoteSection{試験対策ポイント}\begin{itemize}
\item dynamismは、時間の進行により状態が変わり、意思決定を繰り返す性質である。
\item static planning, dynamic planning, a priori optimization, online reoptimizationを区別する。
\end{itemize}
\end{minipage}
\end{adjustbox}
\end{minipage}\hfill
\begin{minipage}[t]{0.58\textwidth}
\begin{adjustbox}{max totalsize={\textwidth}{0.89\textheight},center}
\begin{minipage}{\textwidth}\scriptsize
\NoteSection{解説}このページでは「Static and Dynamic Planning」を理解する。スライド上の表現は短いが、試験では定義、具体例、モデル上の意味、手法の限界まで説明できる必要がある。\par
Dynamism は、意思決定問題が時間とともに進行し、状態が変化することを意味する。静的問題では、すべての入力を最初に与えられ、その情報に基づいて一度だけ計画を作る。動的問題では、新しい注文、交通状況、車両位置、キャンセル、故障などが途中で発生し、計画を更新する必要がある。\par
重要なのは、dynamism は stochasticity と同じではないという点である。stochasticity は未来が不確実であること、dynamism は時間とともに決定と状態が進む構造を持つことである。例えば、すべての未来注文が既知でも、時点ごとに車両位置が変わり決定を繰り返すなら動的である。逆に、将来の需要が確率的でも、最初に一括で決めて変更しないなら静的な確率計画として扱える。\par
a priori optimization は、実行前に全体計画を作り、実行中は大きく変更しない考え方である。計算しやすく、全体最適を考えやすいが、予期しない情報に弱い。online または dynamic planning は、各時点で現在状態を観測し、必要に応じて決定を変える。現実に適応しやすいが、計算時間や一貫性が問題になる。\par
配送の例では、朝に全ルートを決めるのが静的計画である。昼に新規注文が入ったり、車両が遅れたりしたときに、現在の車両位置と残り注文を見て再計画するのが動的計画である。動的計画では、今の選択が将来の柔軟性を奪わないかも考える必要がある。\par
試験答案では、静的/動的の違いを『計画を変更するか』だけでなく、状態、時点、決定、外生情報、遷移の連鎖として説明する。特に動的問題では、状態が更新され、次の決定の入力になるというループを明確に書くと満点答案に近づく。\par\NoteSection{確認問題と解答}\begin{enumerate}\item \textbf{L04-S006: 「Static and Dynamic Planning」の概念を、定義・具体例・モデル上の意味の順に説明せよ。}\\定義としては、dynamismは、時間の進行により状態が変わり、意思決定を繰り返す性質である。。具体例では、配送・在庫・フードデリバリーのような現実問題を用い、状態、決定、外生情報、報酬、または情報モデル/意思決定モデルのどこに対応するかを明示する。モデル上の意味として、この概念が目的関数、制約、状態表現、将来価値、または方策選択にどのような影響を与えるかを書く。試験では、用語だけを列挙せず、入力データから意思決定、さらに現実の実装へつながる流れを文章で示すことが重要である。\item \textbf{L04-S006: このスライドの考え方を、講義全体のDDDMプロセスの中に位置づけよ。}\\DDDMプロセスでは、現実世界のデータを集約して情報モデルを作り、意思決定モデルまたは方策で行動を選び、実装後に結果を観測して次の状態へ進む。このスライドの概念は、そのどこか一部だけではなく、データ表現、意思決定、将来不確実性、計算可能性のいずれかに関わる。答案では、まずこの概念がどの段階に属するかを述べ、次に他の段階へどう影響するかを書く。例えば価値関数なら将来価値を現在決定へ戻す役割、FIFOなら旅行時間情報モデルの整合性を保証する役割、CFAなら目的関数を調整して将来に良い行動を誘導する役割である。\end{enumerate}
\end{minipage}
\end{adjustbox}
\end{minipage}


\clearpage
\SlideHeader{Lecture 4 - Dynamism}{007}{Examples for Dynamism}
\begin{minipage}[t]{0.405\textwidth}
\SlideImage{7}{../Materials/2026/3DM_04_Dynamism_SS26.pdf}
\begin{adjustbox}{max totalsize={\textwidth}{0.43\textheight},center}
\begin{minipage}{\textwidth}\scriptsize
\NoteSection{日本語訳}\begin{itemize}
\item 例s for Dynamism
\item Urban Delivery Line Production Asset Allocation（この用語は講義スライド上の専門語として扱う）
\item https://www.venitec.de/de/leistungen-de/fertigungsplanung-（この用語は講義スライド上の専門語として扱う）
\item https://www.handelsblatt.com/unternehmen/handel- https://www.forbes.com/sites/randywarren/2020/08/06/maxi（この用語は講義スライド上の専門語として扱う）
\item fertigungsoptimierung/fliessfertigung（この用語は講義スライド上の専門語として扱う）
\item konsumgueter/onlinehandel-verkehrskollaps-per- mizing-your-portfolio-benefits-via-asset-allocation/（この用語は講義スライド上の専門語として扱う）
\item mausklick/20604804.html?ticket=ST-190133-（この用語は講義スライド上の専門語として扱う）
\item fsLx4SuFxff4v6EHdGk0-ap5（この用語は講義スライド上の専門語として扱う）
\item Courier Service Inventory Management Games（この用語は講義スライド上の専門語として扱う）
\end{itemize}\NoteSection{試験対策ポイント}\begin{itemize}
\item dynamismは、時間の進行により状態が変わり、意思決定を繰り返す性質である。
\item static planning, dynamic planning, a priori optimization, online reoptimizationを区別する。
\end{itemize}
\end{minipage}
\end{adjustbox}
\end{minipage}\hfill
\begin{minipage}[t]{0.58\textwidth}
\begin{adjustbox}{max totalsize={\textwidth}{0.89\textheight},center}
\begin{minipage}{\textwidth}\scriptsize
\NoteSection{解説}このページでは「Examples for Dynamism」を理解する。スライド上の表現は短いが、試験では定義、具体例、モデル上の意味、手法の限界まで説明できる必要がある。\par
Dynamism は、意思決定問題が時間とともに進行し、状態が変化することを意味する。静的問題では、すべての入力を最初に与えられ、その情報に基づいて一度だけ計画を作る。動的問題では、新しい注文、交通状況、車両位置、キャンセル、故障などが途中で発生し、計画を更新する必要がある。\par
重要なのは、dynamism は stochasticity と同じではないという点である。stochasticity は未来が不確実であること、dynamism は時間とともに決定と状態が進む構造を持つことである。例えば、すべての未来注文が既知でも、時点ごとに車両位置が変わり決定を繰り返すなら動的である。逆に、将来の需要が確率的でも、最初に一括で決めて変更しないなら静的な確率計画として扱える。\par
a priori optimization は、実行前に全体計画を作り、実行中は大きく変更しない考え方である。計算しやすく、全体最適を考えやすいが、予期しない情報に弱い。online または dynamic planning は、各時点で現在状態を観測し、必要に応じて決定を変える。現実に適応しやすいが、計算時間や一貫性が問題になる。\par
配送の例では、朝に全ルートを決めるのが静的計画である。昼に新規注文が入ったり、車両が遅れたりしたときに、現在の車両位置と残り注文を見て再計画するのが動的計画である。動的計画では、今の選択が将来の柔軟性を奪わないかも考える必要がある。\par
試験答案では、静的/動的の違いを『計画を変更するか』だけでなく、状態、時点、決定、外生情報、遷移の連鎖として説明する。特に動的問題では、状態が更新され、次の決定の入力になるというループを明確に書くと満点答案に近づく。\par\NoteSection{確認問題と解答}\begin{enumerate}\item \textbf{L04-S007: 「Examples for Dynamism」の概念を、定義・具体例・モデル上の意味の順に説明せよ。}\\定義としては、dynamismは、時間の進行により状態が変わり、意思決定を繰り返す性質である。。具体例では、配送・在庫・フードデリバリーのような現実問題を用い、状態、決定、外生情報、報酬、または情報モデル/意思決定モデルのどこに対応するかを明示する。モデル上の意味として、この概念が目的関数、制約、状態表現、将来価値、または方策選択にどのような影響を与えるかを書く。試験では、用語だけを列挙せず、入力データから意思決定、さらに現実の実装へつながる流れを文章で示すことが重要である。\item \textbf{L04-S007: このスライドの考え方を、講義全体のDDDMプロセスの中に位置づけよ。}\\DDDMプロセスでは、現実世界のデータを集約して情報モデルを作り、意思決定モデルまたは方策で行動を選び、実装後に結果を観測して次の状態へ進む。このスライドの概念は、そのどこか一部だけではなく、データ表現、意思決定、将来不確実性、計算可能性のいずれかに関わる。答案では、まずこの概念がどの段階に属するかを述べ、次に他の段階へどう影響するかを書く。例えば価値関数なら将来価値を現在決定へ戻す役割、FIFOなら旅行時間情報モデルの整合性を保証する役割、CFAなら目的関数を調整して将来に良い行動を誘導する役割である。\end{enumerate}
\end{minipage}
\end{adjustbox}
\end{minipage}


\clearpage
\SlideHeader{Lecture 4 - Dynamism}{008}{Urban Mobility and Transportation:}
\begin{minipage}[t]{0.405\textwidth}
\SlideImage{8}{../Materials/2026/3DM_04_Dynamism_SS26.pdf}
\begin{adjustbox}{max totalsize={\textwidth}{0.43\textheight},center}
\begin{minipage}{\textwidth}\scriptsize
\NoteSection{日本語訳}\begin{itemize}
\item Urban モビリティ and 交通・輸送:
\item Savelsbergh, van Woensel (交通・輸送 Science, 2016):
\item “Highly 動的 and volatile 決定 making environments with access to
\item massive amounts of near real-time data, which will be seen more and more in（この用語は講義スライド上の専門語として扱う）
\item the future, may necessitate new views, paradigms, and models for 決定
\item support.”（この用語は講義スライド上の専門語として扱う）
\item Speranza (European Journal of Operational Research, 2016):（この用語は講義スライド上の専門語として扱う）
\item New challenge of “appropriately modeling 動的 events and simultaneously
\item incorporating information about the 不確実性 of future events.”
\end{itemize}\NoteSection{試験対策ポイント}\begin{itemize}
\item dynamismは、時間の進行により状態が変わり、意思決定を繰り返す性質である。
\item static planning, dynamic planning, a priori optimization, online reoptimizationを区別する。
\end{itemize}
\end{minipage}
\end{adjustbox}
\end{minipage}\hfill
\begin{minipage}[t]{0.58\textwidth}
\begin{adjustbox}{max totalsize={\textwidth}{0.89\textheight},center}
\begin{minipage}{\textwidth}\scriptsize
\NoteSection{解説}このページでは「Urban Mobility and Transportation:」を理解する。スライド上の表現は短いが、試験では定義、具体例、モデル上の意味、手法の限界まで説明できる必要がある。\par
Dynamism は、意思決定問題が時間とともに進行し、状態が変化することを意味する。静的問題では、すべての入力を最初に与えられ、その情報に基づいて一度だけ計画を作る。動的問題では、新しい注文、交通状況、車両位置、キャンセル、故障などが途中で発生し、計画を更新する必要がある。\par
重要なのは、dynamism は stochasticity と同じではないという点である。stochasticity は未来が不確実であること、dynamism は時間とともに決定と状態が進む構造を持つことである。例えば、すべての未来注文が既知でも、時点ごとに車両位置が変わり決定を繰り返すなら動的である。逆に、将来の需要が確率的でも、最初に一括で決めて変更しないなら静的な確率計画として扱える。\par
a priori optimization は、実行前に全体計画を作り、実行中は大きく変更しない考え方である。計算しやすく、全体最適を考えやすいが、予期しない情報に弱い。online または dynamic planning は、各時点で現在状態を観測し、必要に応じて決定を変える。現実に適応しやすいが、計算時間や一貫性が問題になる。\par
配送の例では、朝に全ルートを決めるのが静的計画である。昼に新規注文が入ったり、車両が遅れたりしたときに、現在の車両位置と残り注文を見て再計画するのが動的計画である。動的計画では、今の選択が将来の柔軟性を奪わないかも考える必要がある。\par
試験答案では、静的/動的の違いを『計画を変更するか』だけでなく、状態、時点、決定、外生情報、遷移の連鎖として説明する。特に動的問題では、状態が更新され、次の決定の入力になるというループを明確に書くと満点答案に近づく。\par\NoteSection{確認問題と解答}\begin{enumerate}\item \textbf{L04-S008: 「Urban Mobility and Transportation:」の概念を、定義・具体例・モデル上の意味の順に説明せよ。}\\定義としては、dynamismは、時間の進行により状態が変わり、意思決定を繰り返す性質である。。具体例では、配送・在庫・フードデリバリーのような現実問題を用い、状態、決定、外生情報、報酬、または情報モデル/意思決定モデルのどこに対応するかを明示する。モデル上の意味として、この概念が目的関数、制約、状態表現、将来価値、または方策選択にどのような影響を与えるかを書く。試験では、用語だけを列挙せず、入力データから意思決定、さらに現実の実装へつながる流れを文章で示すことが重要である。\item \textbf{L04-S008: このスライドの考え方を、講義全体のDDDMプロセスの中に位置づけよ。}\\DDDMプロセスでは、現実世界のデータを集約して情報モデルを作り、意思決定モデルまたは方策で行動を選び、実装後に結果を観測して次の状態へ進む。このスライドの概念は、そのどこか一部だけではなく、データ表現、意思決定、将来不確実性、計算可能性のいずれかに関わる。答案では、まずこの概念がどの段階に属するかを述べ、次に他の段階へどう影響するかを書く。例えば価値関数なら将来価値を現在決定へ戻す役割、FIFOなら旅行時間情報モデルの整合性を保証する役割、CFAなら目的関数を調整して将来に良い行動を誘導する役割である。\end{enumerate}
\end{minipage}
\end{adjustbox}
\end{minipage}


\clearpage
\SlideHeader{Lecture 4 - Dynamism}{009}{Components of a Dynamic Decision Problem}
\begin{minipage}[t]{0.405\textwidth}
\SlideImage{9}{../Materials/2026/3DM_04_Dynamism_SS26.pdf}
\begin{adjustbox}{max totalsize={\textwidth}{0.43\textheight},center}
\begin{minipage}{\textwidth}\scriptsize
\NoteSection{日本語訳}\begin{itemize}
\item Components of a Dynamic 決定 Problem
\end{itemize}\NoteSection{試験対策ポイント}\begin{itemize}
\item dynamismは、時間の進行により状態が変わり、意思決定を繰り返す性質である。
\item static planning, dynamic planning, a priori optimization, online reoptimizationを区別する。
\end{itemize}
\end{minipage}
\end{adjustbox}
\end{minipage}\hfill
\begin{minipage}[t]{0.58\textwidth}
\begin{adjustbox}{max totalsize={\textwidth}{0.89\textheight},center}
\begin{minipage}{\textwidth}\scriptsize
\NoteSection{解説}このページでは「Components of a Dynamic Decision Problem」を理解する。スライド上の表現は短いが、試験では定義、具体例、モデル上の意味、手法の限界まで説明できる必要がある。\par
Dynamism は、意思決定問題が時間とともに進行し、状態が変化することを意味する。静的問題では、すべての入力を最初に与えられ、その情報に基づいて一度だけ計画を作る。動的問題では、新しい注文、交通状況、車両位置、キャンセル、故障などが途中で発生し、計画を更新する必要がある。\par
重要なのは、dynamism は stochasticity と同じではないという点である。stochasticity は未来が不確実であること、dynamism は時間とともに決定と状態が進む構造を持つことである。例えば、すべての未来注文が既知でも、時点ごとに車両位置が変わり決定を繰り返すなら動的である。逆に、将来の需要が確率的でも、最初に一括で決めて変更しないなら静的な確率計画として扱える。\par
a priori optimization は、実行前に全体計画を作り、実行中は大きく変更しない考え方である。計算しやすく、全体最適を考えやすいが、予期しない情報に弱い。online または dynamic planning は、各時点で現在状態を観測し、必要に応じて決定を変える。現実に適応しやすいが、計算時間や一貫性が問題になる。\par
配送の例では、朝に全ルートを決めるのが静的計画である。昼に新規注文が入ったり、車両が遅れたりしたときに、現在の車両位置と残り注文を見て再計画するのが動的計画である。動的計画では、今の選択が将来の柔軟性を奪わないかも考える必要がある。\par
試験答案では、静的/動的の違いを『計画を変更するか』だけでなく、状態、時点、決定、外生情報、遷移の連鎖として説明する。特に動的問題では、状態が更新され、次の決定の入力になるというループを明確に書くと満点答案に近づく。\par\NoteSection{確認問題と解答}\begin{enumerate}\item \textbf{L04-S009: 「Components of a Dynamic Decision Problem」の概念を、定義・具体例・モデル上の意味の順に説明せよ。}\\定義としては、dynamismは、時間の進行により状態が変わり、意思決定を繰り返す性質である。。具体例では、配送・在庫・フードデリバリーのような現実問題を用い、状態、決定、外生情報、報酬、または情報モデル/意思決定モデルのどこに対応するかを明示する。モデル上の意味として、この概念が目的関数、制約、状態表現、将来価値、または方策選択にどのような影響を与えるかを書く。試験では、用語だけを列挙せず、入力データから意思決定、さらに現実の実装へつながる流れを文章で示すことが重要である。\item \textbf{L04-S009: このスライドの考え方を、講義全体のDDDMプロセスの中に位置づけよ。}\\DDDMプロセスでは、現実世界のデータを集約して情報モデルを作り、意思決定モデルまたは方策で行動を選び、実装後に結果を観測して次の状態へ進む。このスライドの概念は、そのどこか一部だけではなく、データ表現、意思決定、将来不確実性、計算可能性のいずれかに関わる。答案では、まずこの概念がどの段階に属するかを述べ、次に他の段階へどう影響するかを書く。例えば価値関数なら将来価値を現在決定へ戻す役割、FIFOなら旅行時間情報モデルの整合性を保証する役割、CFAなら目的関数を調整して将来に良い行動を誘導する役割である。\end{enumerate}
\end{minipage}
\end{adjustbox}
\end{minipage}


\clearpage
\SlideHeader{Lecture 4 - Dynamism}{010}{Ingredients of a Dynamic Decision System}
\begin{minipage}[t]{0.405\textwidth}
\SlideImage{10}{../Materials/2026/3DM_04_Dynamism_SS26.pdf}
\begin{adjustbox}{max totalsize={\textwidth}{0.43\textheight},center}
\begin{minipage}{\textwidth}\scriptsize
\NoteSection{日本語訳}\begin{itemize}
\item Ingredients of a Dynamic 決定 System
\item Exogeneous: Development that cannot（この用語は講義スライド上の専門語として扱う）
\item be controlled by 決定 maker
\item (information model)（この用語は講義スライド上の専門語として扱う）
\item 決定: Endogenous 決定-making
\item process, provides 報酬 or costs
\item (決定 model)
\item Both exogeneous and endogenous（この用語は講義スライド上の専門語として扱う）
\item processes interact via 決定 状態s
\end{itemize}\NoteSection{試験対策ポイント}\begin{itemize}
\item dynamismは、時間の進行により状態が変わり、意思決定を繰り返す性質である。
\item static planning, dynamic planning, a priori optimization, online reoptimizationを区別する。
\end{itemize}
\end{minipage}
\end{adjustbox}
\end{minipage}\hfill
\begin{minipage}[t]{0.58\textwidth}
\begin{adjustbox}{max totalsize={\textwidth}{0.89\textheight},center}
\begin{minipage}{\textwidth}\scriptsize
\NoteSection{解説}このページでは「Ingredients of a Dynamic Decision System」を理解する。スライド上の表現は短いが、試験では定義、具体例、モデル上の意味、手法の限界まで説明できる必要がある。\par
Dynamism は、意思決定問題が時間とともに進行し、状態が変化することを意味する。静的問題では、すべての入力を最初に与えられ、その情報に基づいて一度だけ計画を作る。動的問題では、新しい注文、交通状況、車両位置、キャンセル、故障などが途中で発生し、計画を更新する必要がある。\par
重要なのは、dynamism は stochasticity と同じではないという点である。stochasticity は未来が不確実であること、dynamism は時間とともに決定と状態が進む構造を持つことである。例えば、すべての未来注文が既知でも、時点ごとに車両位置が変わり決定を繰り返すなら動的である。逆に、将来の需要が確率的でも、最初に一括で決めて変更しないなら静的な確率計画として扱える。\par
a priori optimization は、実行前に全体計画を作り、実行中は大きく変更しない考え方である。計算しやすく、全体最適を考えやすいが、予期しない情報に弱い。online または dynamic planning は、各時点で現在状態を観測し、必要に応じて決定を変える。現実に適応しやすいが、計算時間や一貫性が問題になる。\par
配送の例では、朝に全ルートを決めるのが静的計画である。昼に新規注文が入ったり、車両が遅れたりしたときに、現在の車両位置と残り注文を見て再計画するのが動的計画である。動的計画では、今の選択が将来の柔軟性を奪わないかも考える必要がある。\par
試験答案では、静的/動的の違いを『計画を変更するか』だけでなく、状態、時点、決定、外生情報、遷移の連鎖として説明する。特に動的問題では、状態が更新され、次の決定の入力になるというループを明確に書くと満点答案に近づく。\par\NoteSection{確認問題と解答}\begin{enumerate}\item \textbf{L04-S010: 「Ingredients of a Dynamic Decision System」の概念を、定義・具体例・モデル上の意味の順に説明せよ。}\\定義としては、dynamismは、時間の進行により状態が変わり、意思決定を繰り返す性質である。。具体例では、配送・在庫・フードデリバリーのような現実問題を用い、状態、決定、外生情報、報酬、または情報モデル/意思決定モデルのどこに対応するかを明示する。モデル上の意味として、この概念が目的関数、制約、状態表現、将来価値、または方策選択にどのような影響を与えるかを書く。試験では、用語だけを列挙せず、入力データから意思決定、さらに現実の実装へつながる流れを文章で示すことが重要である。\item \textbf{L04-S010: このスライドの考え方を、講義全体のDDDMプロセスの中に位置づけよ。}\\DDDMプロセスでは、現実世界のデータを集約して情報モデルを作り、意思決定モデルまたは方策で行動を選び、実装後に結果を観測して次の状態へ進む。このスライドの概念は、そのどこか一部だけではなく、データ表現、意思決定、将来不確実性、計算可能性のいずれかに関わる。答案では、まずこの概念がどの段階に属するかを述べ、次に他の段階へどう影響するかを書く。例えば価値関数なら将来価値を現在決定へ戻す役割、FIFOなら旅行時間情報モデルの整合性を保証する役割、CFAなら目的関数を調整して将来に良い行動を誘導する役割である。\end{enumerate}
\end{minipage}
\end{adjustbox}
\end{minipage}


\clearpage
\SlideHeader{Lecture 4 - Dynamism}{011}{Urban Delivery}
\begin{minipage}[t]{0.405\textwidth}
\SlideImage{11}{../Materials/2026/3DM_04_Dynamism_SS26.pdf}
\begin{adjustbox}{max totalsize={\textwidth}{0.43\textheight},center}
\begin{minipage}{\textwidth}\scriptsize
\NoteSection{日本語訳}\begin{itemize}
\item Urban Delivery（この用語は講義スライド上の専門語として扱う）
\item Vehicle routing problem（この用語は講義スライド上の専門語として扱う）
\item All customers must be served（この用語は講義スライド上の専門語として扱う）
\item One depot, one vehicle（この用語は講義スライド上の専門語として扱う）
\item Travel times vary over the day（この用語は講義スライド上の専門語として扱う）
\item D
\item Exogenous: Varying 旅行時間s
\item 状態: Customers to serve, Position of vehicle 1 3
\item 決定: Whom to visit next
\end{itemize}\NoteSection{試験対策ポイント}\begin{itemize}
\item dynamismは、時間の進行により状態が変わり、意思決定を繰り返す性質である。
\item static planning, dynamic planning, a priori optimization, online reoptimizationを区別する。
\end{itemize}
\end{minipage}
\end{adjustbox}
\end{minipage}\hfill
\begin{minipage}[t]{0.58\textwidth}
\begin{adjustbox}{max totalsize={\textwidth}{0.89\textheight},center}
\begin{minipage}{\textwidth}\scriptsize
\NoteSection{解説}このページでは「Urban Delivery」を理解する。スライド上の表現は短いが、試験では定義、具体例、モデル上の意味、手法の限界まで説明できる必要がある。\par
Dynamism は、意思決定問題が時間とともに進行し、状態が変化することを意味する。静的問題では、すべての入力を最初に与えられ、その情報に基づいて一度だけ計画を作る。動的問題では、新しい注文、交通状況、車両位置、キャンセル、故障などが途中で発生し、計画を更新する必要がある。\par
重要なのは、dynamism は stochasticity と同じではないという点である。stochasticity は未来が不確実であること、dynamism は時間とともに決定と状態が進む構造を持つことである。例えば、すべての未来注文が既知でも、時点ごとに車両位置が変わり決定を繰り返すなら動的である。逆に、将来の需要が確率的でも、最初に一括で決めて変更しないなら静的な確率計画として扱える。\par
a priori optimization は、実行前に全体計画を作り、実行中は大きく変更しない考え方である。計算しやすく、全体最適を考えやすいが、予期しない情報に弱い。online または dynamic planning は、各時点で現在状態を観測し、必要に応じて決定を変える。現実に適応しやすいが、計算時間や一貫性が問題になる。\par
配送の例では、朝に全ルートを決めるのが静的計画である。昼に新規注文が入ったり、車両が遅れたりしたときに、現在の車両位置と残り注文を見て再計画するのが動的計画である。動的計画では、今の選択が将来の柔軟性を奪わないかも考える必要がある。\par
試験答案では、静的/動的の違いを『計画を変更するか』だけでなく、状態、時点、決定、外生情報、遷移の連鎖として説明する。特に動的問題では、状態が更新され、次の決定の入力になるというループを明確に書くと満点答案に近づく。\par\NoteSection{確認問題と解答}\begin{enumerate}\item \textbf{L04-S011: 「Urban Delivery」の概念を、定義・具体例・モデル上の意味の順に説明せよ。}\\定義としては、dynamismは、時間の進行により状態が変わり、意思決定を繰り返す性質である。。具体例では、配送・在庫・フードデリバリーのような現実問題を用い、状態、決定、外生情報、報酬、または情報モデル/意思決定モデルのどこに対応するかを明示する。モデル上の意味として、この概念が目的関数、制約、状態表現、将来価値、または方策選択にどのような影響を与えるかを書く。試験では、用語だけを列挙せず、入力データから意思決定、さらに現実の実装へつながる流れを文章で示すことが重要である。\item \textbf{L04-S011: このスライドの考え方を、講義全体のDDDMプロセスの中に位置づけよ。}\\DDDMプロセスでは、現実世界のデータを集約して情報モデルを作り、意思決定モデルまたは方策で行動を選び、実装後に結果を観測して次の状態へ進む。このスライドの概念は、そのどこか一部だけではなく、データ表現、意思決定、将来不確実性、計算可能性のいずれかに関わる。答案では、まずこの概念がどの段階に属するかを述べ、次に他の段階へどう影響するかを書く。例えば価値関数なら将来価値を現在決定へ戻す役割、FIFOなら旅行時間情報モデルの整合性を保証する役割、CFAなら目的関数を調整して将来に良い行動を誘導する役割である。\end{enumerate}
\end{minipage}
\end{adjustbox}
\end{minipage}


\clearpage
\SlideHeader{Lecture 4 - Dynamism}{012}{Example: Urban Delivery}
\begin{minipage}[t]{0.405\textwidth}
\SlideImage{12}{../Materials/2026/3DM_04_Dynamism_SS26.pdf}
\begin{adjustbox}{max totalsize={\textwidth}{0.43\textheight},center}
\begin{minipage}{\textwidth}\scriptsize
\NoteSection{日本語訳}\begin{itemize}
\item 例: Urban Delivery
\item Route has been optimized prior to execution（この用語は講義スライド上の専門語として扱う）
\item D
\item Execution starts with proceeding to customer 1（この用語は講義スライド上の専門語として扱う）
\end{itemize}\NoteSection{試験対策ポイント}\begin{itemize}
\item dynamismは、時間の進行により状態が変わり、意思決定を繰り返す性質である。
\item static planning, dynamic planning, a priori optimization, online reoptimizationを区別する。
\end{itemize}
\end{minipage}
\end{adjustbox}
\end{minipage}\hfill
\begin{minipage}[t]{0.58\textwidth}
\begin{adjustbox}{max totalsize={\textwidth}{0.89\textheight},center}
\begin{minipage}{\textwidth}\scriptsize
\NoteSection{解説}このページでは「Example: Urban Delivery」を理解する。スライド上の表現は短いが、試験では定義、具体例、モデル上の意味、手法の限界まで説明できる必要がある。\par
Dynamism は、意思決定問題が時間とともに進行し、状態が変化することを意味する。静的問題では、すべての入力を最初に与えられ、その情報に基づいて一度だけ計画を作る。動的問題では、新しい注文、交通状況、車両位置、キャンセル、故障などが途中で発生し、計画を更新する必要がある。\par
重要なのは、dynamism は stochasticity と同じではないという点である。stochasticity は未来が不確実であること、dynamism は時間とともに決定と状態が進む構造を持つことである。例えば、すべての未来注文が既知でも、時点ごとに車両位置が変わり決定を繰り返すなら動的である。逆に、将来の需要が確率的でも、最初に一括で決めて変更しないなら静的な確率計画として扱える。\par
a priori optimization は、実行前に全体計画を作り、実行中は大きく変更しない考え方である。計算しやすく、全体最適を考えやすいが、予期しない情報に弱い。online または dynamic planning は、各時点で現在状態を観測し、必要に応じて決定を変える。現実に適応しやすいが、計算時間や一貫性が問題になる。\par
配送の例では、朝に全ルートを決めるのが静的計画である。昼に新規注文が入ったり、車両が遅れたりしたときに、現在の車両位置と残り注文を見て再計画するのが動的計画である。動的計画では、今の選択が将来の柔軟性を奪わないかも考える必要がある。\par
試験答案では、静的/動的の違いを『計画を変更するか』だけでなく、状態、時点、決定、外生情報、遷移の連鎖として説明する。特に動的問題では、状態が更新され、次の決定の入力になるというループを明確に書くと満点答案に近づく。\par\NoteSection{確認問題と解答}\begin{enumerate}\item \textbf{L04-S012: 「Example: Urban Delivery」の概念を、定義・具体例・モデル上の意味の順に説明せよ。}\\定義としては、dynamismは、時間の進行により状態が変わり、意思決定を繰り返す性質である。。具体例では、配送・在庫・フードデリバリーのような現実問題を用い、状態、決定、外生情報、報酬、または情報モデル/意思決定モデルのどこに対応するかを明示する。モデル上の意味として、この概念が目的関数、制約、状態表現、将来価値、または方策選択にどのような影響を与えるかを書く。試験では、用語だけを列挙せず、入力データから意思決定、さらに現実の実装へつながる流れを文章で示すことが重要である。\item \textbf{L04-S012: このスライドの考え方を、講義全体のDDDMプロセスの中に位置づけよ。}\\DDDMプロセスでは、現実世界のデータを集約して情報モデルを作り、意思決定モデルまたは方策で行動を選び、実装後に結果を観測して次の状態へ進む。このスライドの概念は、そのどこか一部だけではなく、データ表現、意思決定、将来不確実性、計算可能性のいずれかに関わる。答案では、まずこの概念がどの段階に属するかを述べ、次に他の段階へどう影響するかを書く。例えば価値関数なら将来価値を現在決定へ戻す役割、FIFOなら旅行時間情報モデルの整合性を保証する役割、CFAなら目的関数を調整して将来に良い行動を誘導する役割である。\end{enumerate}
\end{minipage}
\end{adjustbox}
\end{minipage}


\clearpage
\SlideHeader{Lecture 4 - Dynamism}{013}{Example: Urban Delivery}
\begin{minipage}[t]{0.405\textwidth}
\SlideImage{13}{../Materials/2026/3DM_04_Dynamism_SS26.pdf}
\begin{adjustbox}{max totalsize={\textwidth}{0.43\textheight},center}
\begin{minipage}{\textwidth}\scriptsize
\NoteSection{日本語訳}\begin{itemize}
\item 例: Urban Delivery
\item Having arrived at customer 1, a congestion alert comes in（この用語は講義スライド上の専門語として扱う）
\item D
\end{itemize}\NoteSection{試験対策ポイント}\begin{itemize}
\item dynamismは、時間の進行により状態が変わり、意思決定を繰り返す性質である。
\item static planning, dynamic planning, a priori optimization, online reoptimizationを区別する。
\end{itemize}
\end{minipage}
\end{adjustbox}
\end{minipage}\hfill
\begin{minipage}[t]{0.58\textwidth}
\begin{adjustbox}{max totalsize={\textwidth}{0.89\textheight},center}
\begin{minipage}{\textwidth}\scriptsize
\NoteSection{解説}このページでは「Example: Urban Delivery」を理解する。スライド上の表現は短いが、試験では定義、具体例、モデル上の意味、手法の限界まで説明できる必要がある。\par
Dynamism は、意思決定問題が時間とともに進行し、状態が変化することを意味する。静的問題では、すべての入力を最初に与えられ、その情報に基づいて一度だけ計画を作る。動的問題では、新しい注文、交通状況、車両位置、キャンセル、故障などが途中で発生し、計画を更新する必要がある。\par
重要なのは、dynamism は stochasticity と同じではないという点である。stochasticity は未来が不確実であること、dynamism は時間とともに決定と状態が進む構造を持つことである。例えば、すべての未来注文が既知でも、時点ごとに車両位置が変わり決定を繰り返すなら動的である。逆に、将来の需要が確率的でも、最初に一括で決めて変更しないなら静的な確率計画として扱える。\par
a priori optimization は、実行前に全体計画を作り、実行中は大きく変更しない考え方である。計算しやすく、全体最適を考えやすいが、予期しない情報に弱い。online または dynamic planning は、各時点で現在状態を観測し、必要に応じて決定を変える。現実に適応しやすいが、計算時間や一貫性が問題になる。\par
配送の例では、朝に全ルートを決めるのが静的計画である。昼に新規注文が入ったり、車両が遅れたりしたときに、現在の車両位置と残り注文を見て再計画するのが動的計画である。動的計画では、今の選択が将来の柔軟性を奪わないかも考える必要がある。\par
試験答案では、静的/動的の違いを『計画を変更するか』だけでなく、状態、時点、決定、外生情報、遷移の連鎖として説明する。特に動的問題では、状態が更新され、次の決定の入力になるというループを明確に書くと満点答案に近づく。\par\NoteSection{確認問題と解答}\begin{enumerate}\item \textbf{L04-S013: 「Example: Urban Delivery」の概念を、定義・具体例・モデル上の意味の順に説明せよ。}\\定義としては、dynamismは、時間の進行により状態が変わり、意思決定を繰り返す性質である。。具体例では、配送・在庫・フードデリバリーのような現実問題を用い、状態、決定、外生情報、報酬、または情報モデル/意思決定モデルのどこに対応するかを明示する。モデル上の意味として、この概念が目的関数、制約、状態表現、将来価値、または方策選択にどのような影響を与えるかを書く。試験では、用語だけを列挙せず、入力データから意思決定、さらに現実の実装へつながる流れを文章で示すことが重要である。\item \textbf{L04-S013: このスライドの考え方を、講義全体のDDDMプロセスの中に位置づけよ。}\\DDDMプロセスでは、現実世界のデータを集約して情報モデルを作り、意思決定モデルまたは方策で行動を選び、実装後に結果を観測して次の状態へ進む。このスライドの概念は、そのどこか一部だけではなく、データ表現、意思決定、将来不確実性、計算可能性のいずれかに関わる。答案では、まずこの概念がどの段階に属するかを述べ、次に他の段階へどう影響するかを書く。例えば価値関数なら将来価値を現在決定へ戻す役割、FIFOなら旅行時間情報モデルの整合性を保証する役割、CFAなら目的関数を調整して将来に良い行動を誘導する役割である。\end{enumerate}
\end{minipage}
\end{adjustbox}
\end{minipage}


\clearpage
\SlideHeader{Lecture 4 - Dynamism}{014}{Example: Urban Delivery}
\begin{minipage}[t]{0.405\textwidth}
\SlideImage{14}{../Materials/2026/3DM_04_Dynamism_SS26.pdf}
\begin{adjustbox}{max totalsize={\textwidth}{0.43\textheight},center}
\begin{minipage}{\textwidth}\scriptsize
\NoteSection{日本語訳}\begin{itemize}
\item 例: Urban Delivery
\item One option on the way to customer 2 is a costly detour（この用語は講義スライド上の専門語として扱う）
\item D
\end{itemize}\NoteSection{試験対策ポイント}\begin{itemize}
\item dynamismは、時間の進行により状態が変わり、意思決定を繰り返す性質である。
\item static planning, dynamic planning, a priori optimization, online reoptimizationを区別する。
\end{itemize}
\end{minipage}
\end{adjustbox}
\end{minipage}\hfill
\begin{minipage}[t]{0.58\textwidth}
\begin{adjustbox}{max totalsize={\textwidth}{0.89\textheight},center}
\begin{minipage}{\textwidth}\scriptsize
\NoteSection{解説}このページでは「Example: Urban Delivery」を理解する。スライド上の表現は短いが、試験では定義、具体例、モデル上の意味、手法の限界まで説明できる必要がある。\par
Dynamism は、意思決定問題が時間とともに進行し、状態が変化することを意味する。静的問題では、すべての入力を最初に与えられ、その情報に基づいて一度だけ計画を作る。動的問題では、新しい注文、交通状況、車両位置、キャンセル、故障などが途中で発生し、計画を更新する必要がある。\par
重要なのは、dynamism は stochasticity と同じではないという点である。stochasticity は未来が不確実であること、dynamism は時間とともに決定と状態が進む構造を持つことである。例えば、すべての未来注文が既知でも、時点ごとに車両位置が変わり決定を繰り返すなら動的である。逆に、将来の需要が確率的でも、最初に一括で決めて変更しないなら静的な確率計画として扱える。\par
a priori optimization は、実行前に全体計画を作り、実行中は大きく変更しない考え方である。計算しやすく、全体最適を考えやすいが、予期しない情報に弱い。online または dynamic planning は、各時点で現在状態を観測し、必要に応じて決定を変える。現実に適応しやすいが、計算時間や一貫性が問題になる。\par
配送の例では、朝に全ルートを決めるのが静的計画である。昼に新規注文が入ったり、車両が遅れたりしたときに、現在の車両位置と残り注文を見て再計画するのが動的計画である。動的計画では、今の選択が将来の柔軟性を奪わないかも考える必要がある。\par
試験答案では、静的/動的の違いを『計画を変更するか』だけでなく、状態、時点、決定、外生情報、遷移の連鎖として説明する。特に動的問題では、状態が更新され、次の決定の入力になるというループを明確に書くと満点答案に近づく。\par\NoteSection{確認問題と解答}\begin{enumerate}\item \textbf{L04-S014: 「Example: Urban Delivery」の概念を、定義・具体例・モデル上の意味の順に説明せよ。}\\定義としては、dynamismは、時間の進行により状態が変わり、意思決定を繰り返す性質である。。具体例では、配送・在庫・フードデリバリーのような現実問題を用い、状態、決定、外生情報、報酬、または情報モデル/意思決定モデルのどこに対応するかを明示する。モデル上の意味として、この概念が目的関数、制約、状態表現、将来価値、または方策選択にどのような影響を与えるかを書く。試験では、用語だけを列挙せず、入力データから意思決定、さらに現実の実装へつながる流れを文章で示すことが重要である。\item \textbf{L04-S014: このスライドの考え方を、講義全体のDDDMプロセスの中に位置づけよ。}\\DDDMプロセスでは、現実世界のデータを集約して情報モデルを作り、意思決定モデルまたは方策で行動を選び、実装後に結果を観測して次の状態へ進む。このスライドの概念は、そのどこか一部だけではなく、データ表現、意思決定、将来不確実性、計算可能性のいずれかに関わる。答案では、まずこの概念がどの段階に属するかを述べ、次に他の段階へどう影響するかを書く。例えば価値関数なら将来価値を現在決定へ戻す役割、FIFOなら旅行時間情報モデルの整合性を保証する役割、CFAなら目的関数を調整して将来に良い行動を誘導する役割である。\end{enumerate}
\end{minipage}
\end{adjustbox}
\end{minipage}


\clearpage
\SlideHeader{Lecture 4 - Dynamism}{015}{Example: Urban Delivery}
\begin{minipage}[t]{0.405\textwidth}
\SlideImage{15}{../Materials/2026/3DM_04_Dynamism_SS26.pdf}
\begin{adjustbox}{max totalsize={\textwidth}{0.43\textheight},center}
\begin{minipage}{\textwidth}\scriptsize
\NoteSection{日本語訳}\begin{itemize}
\item 例: Urban Delivery
\item Another option is to switch the sequence of customer visits（この用語は講義スライド上の専門語として扱う）
\item D
\end{itemize}\NoteSection{試験対策ポイント}\begin{itemize}
\item dynamismは、時間の進行により状態が変わり、意思決定を繰り返す性質である。
\item static planning, dynamic planning, a priori optimization, online reoptimizationを区別する。
\end{itemize}
\end{minipage}
\end{adjustbox}
\end{minipage}\hfill
\begin{minipage}[t]{0.58\textwidth}
\begin{adjustbox}{max totalsize={\textwidth}{0.89\textheight},center}
\begin{minipage}{\textwidth}\scriptsize
\NoteSection{解説}このページでは「Example: Urban Delivery」を理解する。スライド上の表現は短いが、試験では定義、具体例、モデル上の意味、手法の限界まで説明できる必要がある。\par
Dynamism は、意思決定問題が時間とともに進行し、状態が変化することを意味する。静的問題では、すべての入力を最初に与えられ、その情報に基づいて一度だけ計画を作る。動的問題では、新しい注文、交通状況、車両位置、キャンセル、故障などが途中で発生し、計画を更新する必要がある。\par
重要なのは、dynamism は stochasticity と同じではないという点である。stochasticity は未来が不確実であること、dynamism は時間とともに決定と状態が進む構造を持つことである。例えば、すべての未来注文が既知でも、時点ごとに車両位置が変わり決定を繰り返すなら動的である。逆に、将来の需要が確率的でも、最初に一括で決めて変更しないなら静的な確率計画として扱える。\par
a priori optimization は、実行前に全体計画を作り、実行中は大きく変更しない考え方である。計算しやすく、全体最適を考えやすいが、予期しない情報に弱い。online または dynamic planning は、各時点で現在状態を観測し、必要に応じて決定を変える。現実に適応しやすいが、計算時間や一貫性が問題になる。\par
配送の例では、朝に全ルートを決めるのが静的計画である。昼に新規注文が入ったり、車両が遅れたりしたときに、現在の車両位置と残り注文を見て再計画するのが動的計画である。動的計画では、今の選択が将来の柔軟性を奪わないかも考える必要がある。\par
試験答案では、静的/動的の違いを『計画を変更するか』だけでなく、状態、時点、決定、外生情報、遷移の連鎖として説明する。特に動的問題では、状態が更新され、次の決定の入力になるというループを明確に書くと満点答案に近づく。\par\NoteSection{確認問題と解答}\begin{enumerate}\item \textbf{L04-S015: 「Example: Urban Delivery」の概念を、定義・具体例・モデル上の意味の順に説明せよ。}\\定義としては、dynamismは、時間の進行により状態が変わり、意思決定を繰り返す性質である。。具体例では、配送・在庫・フードデリバリーのような現実問題を用い、状態、決定、外生情報、報酬、または情報モデル/意思決定モデルのどこに対応するかを明示する。モデル上の意味として、この概念が目的関数、制約、状態表現、将来価値、または方策選択にどのような影響を与えるかを書く。試験では、用語だけを列挙せず、入力データから意思決定、さらに現実の実装へつながる流れを文章で示すことが重要である。\item \textbf{L04-S015: このスライドの考え方を、講義全体のDDDMプロセスの中に位置づけよ。}\\DDDMプロセスでは、現実世界のデータを集約して情報モデルを作り、意思決定モデルまたは方策で行動を選び、実装後に結果を観測して次の状態へ進む。このスライドの概念は、そのどこか一部だけではなく、データ表現、意思決定、将来不確実性、計算可能性のいずれかに関わる。答案では、まずこの概念がどの段階に属するかを述べ、次に他の段階へどう影響するかを書く。例えば価値関数なら将来価値を現在決定へ戻す役割、FIFOなら旅行時間情報モデルの整合性を保証する役割、CFAなら目的関数を調整して将来に良い行動を誘導する役割である。\end{enumerate}
\end{minipage}
\end{adjustbox}
\end{minipage}


\clearpage
\SlideHeader{Lecture 4 - Dynamism}{016}{Dynamic Decision Process}
\begin{minipage}[t]{0.405\textwidth}
\SlideImage{16}{../Materials/2026/3DM_04_Dynamism_SS26.pdf}
\begin{adjustbox}{max totalsize={\textwidth}{0.43\textheight},center}
\begin{minipage}{\textwidth}\scriptsize
\NoteSection{日本語訳}\begin{itemize}
\item 動的意思決定プロセス
\item 旅行時間 Realizations
\item Instance（この用語は講義スライド上の専門語として扱う）
\item Open Traveling Salesperson Problem（この用語は講義スライド上の専門語として扱う）
\end{itemize}\NoteSection{試験対策ポイント}\begin{itemize}
\item dynamismは、時間の進行により状態が変わり、意思決定を繰り返す性質である。
\item static planning, dynamic planning, a priori optimization, online reoptimizationを区別する。
\end{itemize}
\end{minipage}
\end{adjustbox}
\end{minipage}\hfill
\begin{minipage}[t]{0.58\textwidth}
\begin{adjustbox}{max totalsize={\textwidth}{0.89\textheight},center}
\begin{minipage}{\textwidth}\scriptsize
\NoteSection{解説}このページでは「Dynamic Decision Process」を理解する。スライド上の表現は短いが、試験では定義、具体例、モデル上の意味、手法の限界まで説明できる必要がある。\par
Dynamism は、意思決定問題が時間とともに進行し、状態が変化することを意味する。静的問題では、すべての入力を最初に与えられ、その情報に基づいて一度だけ計画を作る。動的問題では、新しい注文、交通状況、車両位置、キャンセル、故障などが途中で発生し、計画を更新する必要がある。\par
重要なのは、dynamism は stochasticity と同じではないという点である。stochasticity は未来が不確実であること、dynamism は時間とともに決定と状態が進む構造を持つことである。例えば、すべての未来注文が既知でも、時点ごとに車両位置が変わり決定を繰り返すなら動的である。逆に、将来の需要が確率的でも、最初に一括で決めて変更しないなら静的な確率計画として扱える。\par
a priori optimization は、実行前に全体計画を作り、実行中は大きく変更しない考え方である。計算しやすく、全体最適を考えやすいが、予期しない情報に弱い。online または dynamic planning は、各時点で現在状態を観測し、必要に応じて決定を変える。現実に適応しやすいが、計算時間や一貫性が問題になる。\par
配送の例では、朝に全ルートを決めるのが静的計画である。昼に新規注文が入ったり、車両が遅れたりしたときに、現在の車両位置と残り注文を見て再計画するのが動的計画である。動的計画では、今の選択が将来の柔軟性を奪わないかも考える必要がある。\par
試験答案では、静的/動的の違いを『計画を変更するか』だけでなく、状態、時点、決定、外生情報、遷移の連鎖として説明する。特に動的問題では、状態が更新され、次の決定の入力になるというループを明確に書くと満点答案に近づく。\par\NoteSection{確認問題と解答}\begin{enumerate}\item \textbf{L04-S016: 「Dynamic Decision Process」の概念を、定義・具体例・モデル上の意味の順に説明せよ。}\\定義としては、dynamismは、時間の進行により状態が変わり、意思決定を繰り返す性質である。。具体例では、配送・在庫・フードデリバリーのような現実問題を用い、状態、決定、外生情報、報酬、または情報モデル/意思決定モデルのどこに対応するかを明示する。モデル上の意味として、この概念が目的関数、制約、状態表現、将来価値、または方策選択にどのような影響を与えるかを書く。試験では、用語だけを列挙せず、入力データから意思決定、さらに現実の実装へつながる流れを文章で示すことが重要である。\item \textbf{L04-S016: このスライドの考え方を、講義全体のDDDMプロセスの中に位置づけよ。}\\DDDMプロセスでは、現実世界のデータを集約して情報モデルを作り、意思決定モデルまたは方策で行動を選び、実装後に結果を観測して次の状態へ進む。このスライドの概念は、そのどこか一部だけではなく、データ表現、意思決定、将来不確実性、計算可能性のいずれかに関わる。答案では、まずこの概念がどの段階に属するかを述べ、次に他の段階へどう影響するかを書く。例えば価値関数なら将来価値を現在決定へ戻す役割、FIFOなら旅行時間情報モデルの整合性を保証する役割、CFAなら目的関数を調整して将来に良い行動を誘導する役割である。\end{enumerate}
\end{minipage}
\end{adjustbox}
\end{minipage}


\clearpage
\SlideHeader{Lecture 4 - Dynamism}{017}{Dynamic Decision Process}
\begin{minipage}[t]{0.405\textwidth}
\SlideImage{17}{../Materials/2026/3DM_04_Dynamism_SS26.pdf}
\begin{adjustbox}{max totalsize={\textwidth}{0.43\textheight},center}
\begin{minipage}{\textwidth}\scriptsize
\NoteSection{日本語訳}\begin{itemize}
\item 動的意思決定プロセス
\item Change in 旅行時間 matrix
\item Position（この用語は講義スライド上の専門語として扱う）
\item Customers left to serve（この用語は講義スライド上の専門語として扱う）
\item Travel time matrix（この用語は講義スライド上の専門語として扱う）
\item Next customer to visit（この用語は講義スライド上の専門語として扱う）
\item (new tentative route)（この用語は講義スライド上の専門語として扱う）
\item Travel time to next（この用語は講義スライド上の専門語として扱う）
\item customer（この用語は講義スライド上の専門語として扱う）
\end{itemize}\NoteSection{試験対策ポイント}\begin{itemize}
\item dynamismは、時間の進行により状態が変わり、意思決定を繰り返す性質である。
\item static planning, dynamic planning, a priori optimization, online reoptimizationを区別する。
\end{itemize}
\end{minipage}
\end{adjustbox}
\end{minipage}\hfill
\begin{minipage}[t]{0.58\textwidth}
\begin{adjustbox}{max totalsize={\textwidth}{0.89\textheight},center}
\begin{minipage}{\textwidth}\scriptsize
\NoteSection{解説}このページでは「Dynamic Decision Process」を理解する。スライド上の表現は短いが、試験では定義、具体例、モデル上の意味、手法の限界まで説明できる必要がある。\par
Dynamism は、意思決定問題が時間とともに進行し、状態が変化することを意味する。静的問題では、すべての入力を最初に与えられ、その情報に基づいて一度だけ計画を作る。動的問題では、新しい注文、交通状況、車両位置、キャンセル、故障などが途中で発生し、計画を更新する必要がある。\par
重要なのは、dynamism は stochasticity と同じではないという点である。stochasticity は未来が不確実であること、dynamism は時間とともに決定と状態が進む構造を持つことである。例えば、すべての未来注文が既知でも、時点ごとに車両位置が変わり決定を繰り返すなら動的である。逆に、将来の需要が確率的でも、最初に一括で決めて変更しないなら静的な確率計画として扱える。\par
a priori optimization は、実行前に全体計画を作り、実行中は大きく変更しない考え方である。計算しやすく、全体最適を考えやすいが、予期しない情報に弱い。online または dynamic planning は、各時点で現在状態を観測し、必要に応じて決定を変える。現実に適応しやすいが、計算時間や一貫性が問題になる。\par
配送の例では、朝に全ルートを決めるのが静的計画である。昼に新規注文が入ったり、車両が遅れたりしたときに、現在の車両位置と残り注文を見て再計画するのが動的計画である。動的計画では、今の選択が将来の柔軟性を奪わないかも考える必要がある。\par
試験答案では、静的/動的の違いを『計画を変更するか』だけでなく、状態、時点、決定、外生情報、遷移の連鎖として説明する。特に動的問題では、状態が更新され、次の決定の入力になるというループを明確に書くと満点答案に近づく。\par\NoteSection{確認問題と解答}\begin{enumerate}\item \textbf{L04-S017: 「Dynamic Decision Process」の概念を、定義・具体例・モデル上の意味の順に説明せよ。}\\定義としては、dynamismは、時間の進行により状態が変わり、意思決定を繰り返す性質である。。具体例では、配送・在庫・フードデリバリーのような現実問題を用い、状態、決定、外生情報、報酬、または情報モデル/意思決定モデルのどこに対応するかを明示する。モデル上の意味として、この概念が目的関数、制約、状態表現、将来価値、または方策選択にどのような影響を与えるかを書く。試験では、用語だけを列挙せず、入力データから意思決定、さらに現実の実装へつながる流れを文章で示すことが重要である。\item \textbf{L04-S017: このスライドの考え方を、講義全体のDDDMプロセスの中に位置づけよ。}\\DDDMプロセスでは、現実世界のデータを集約して情報モデルを作り、意思決定モデルまたは方策で行動を選び、実装後に結果を観測して次の状態へ進む。このスライドの概念は、そのどこか一部だけではなく、データ表現、意思決定、将来不確実性、計算可能性のいずれかに関わる。答案では、まずこの概念がどの段階に属するかを述べ、次に他の段階へどう影響するかを書く。例えば価値関数なら将来価値を現在決定へ戻す役割、FIFOなら旅行時間情報モデルの整合性を保証する役割、CFAなら目的関数を調整して将来に良い行動を誘導する役割である。\end{enumerate}
\end{minipage}
\end{adjustbox}
\end{minipage}


\clearpage
\SlideHeader{Lecture 4 - Dynamism}{018}{Dynamic Decision Process}
\begin{minipage}[t]{0.405\textwidth}
\SlideImage{18}{../Materials/2026/3DM_04_Dynamism_SS26.pdf}
\begin{adjustbox}{max totalsize={\textwidth}{0.43\textheight},center}
\begin{minipage}{\textwidth}\scriptsize
\NoteSection{日本語訳}\begin{itemize}
\item 動的意思決定プロセス
\item Position（この用語は講義スライド上の専門語として扱う）
\item Customers left to serve（この用語は講義スライド上の専門語として扱う）
\item Travel time matrix（この用語は講義スライド上の専門語として扱う）
\item Next customer to visit（この用語は講義スライド上の専門語として扱う）
\item (new tentative route)（この用語は講義スライド上の専門語として扱う）
\item Travel time to next（この用語は講義スライド上の専門語として扱う）
\item customer（この用語は講義スライド上の専門語として扱う）
\end{itemize}\NoteSection{試験対策ポイント}\begin{itemize}
\item dynamismは、時間の進行により状態が変わり、意思決定を繰り返す性質である。
\item static planning, dynamic planning, a priori optimization, online reoptimizationを区別する。
\end{itemize}
\end{minipage}
\end{adjustbox}
\end{minipage}\hfill
\begin{minipage}[t]{0.58\textwidth}
\begin{adjustbox}{max totalsize={\textwidth}{0.89\textheight},center}
\begin{minipage}{\textwidth}\scriptsize
\NoteSection{解説}このページでは「Dynamic Decision Process」を理解する。スライド上の表現は短いが、試験では定義、具体例、モデル上の意味、手法の限界まで説明できる必要がある。\par
Dynamism は、意思決定問題が時間とともに進行し、状態が変化することを意味する。静的問題では、すべての入力を最初に与えられ、その情報に基づいて一度だけ計画を作る。動的問題では、新しい注文、交通状況、車両位置、キャンセル、故障などが途中で発生し、計画を更新する必要がある。\par
重要なのは、dynamism は stochasticity と同じではないという点である。stochasticity は未来が不確実であること、dynamism は時間とともに決定と状態が進む構造を持つことである。例えば、すべての未来注文が既知でも、時点ごとに車両位置が変わり決定を繰り返すなら動的である。逆に、将来の需要が確率的でも、最初に一括で決めて変更しないなら静的な確率計画として扱える。\par
a priori optimization は、実行前に全体計画を作り、実行中は大きく変更しない考え方である。計算しやすく、全体最適を考えやすいが、予期しない情報に弱い。online または dynamic planning は、各時点で現在状態を観測し、必要に応じて決定を変える。現実に適応しやすいが、計算時間や一貫性が問題になる。\par
配送の例では、朝に全ルートを決めるのが静的計画である。昼に新規注文が入ったり、車両が遅れたりしたときに、現在の車両位置と残り注文を見て再計画するのが動的計画である。動的計画では、今の選択が将来の柔軟性を奪わないかも考える必要がある。\par
試験答案では、静的/動的の違いを『計画を変更するか』だけでなく、状態、時点、決定、外生情報、遷移の連鎖として説明する。特に動的問題では、状態が更新され、次の決定の入力になるというループを明確に書くと満点答案に近づく。\par\NoteSection{確認問題と解答}\begin{enumerate}\item \textbf{L04-S018: 「Dynamic Decision Process」の概念を、定義・具体例・モデル上の意味の順に説明せよ。}\\定義としては、dynamismは、時間の進行により状態が変わり、意思決定を繰り返す性質である。。具体例では、配送・在庫・フードデリバリーのような現実問題を用い、状態、決定、外生情報、報酬、または情報モデル/意思決定モデルのどこに対応するかを明示する。モデル上の意味として、この概念が目的関数、制約、状態表現、将来価値、または方策選択にどのような影響を与えるかを書く。試験では、用語だけを列挙せず、入力データから意思決定、さらに現実の実装へつながる流れを文章で示すことが重要である。\item \textbf{L04-S018: このスライドの考え方を、講義全体のDDDMプロセスの中に位置づけよ。}\\DDDMプロセスでは、現実世界のデータを集約して情報モデルを作り、意思決定モデルまたは方策で行動を選び、実装後に結果を観測して次の状態へ進む。このスライドの概念は、そのどこか一部だけではなく、データ表現、意思決定、将来不確実性、計算可能性のいずれかに関わる。答案では、まずこの概念がどの段階に属するかを述べ、次に他の段階へどう影響するかを書く。例えば価値関数なら将来価値を現在決定へ戻す役割、FIFOなら旅行時間情報モデルの整合性を保証する役割、CFAなら目的関数を調整して将来に良い行動を誘導する役割である。\end{enumerate}
\end{minipage}
\end{adjustbox}
\end{minipage}


\clearpage
\SlideHeader{Lecture 4 - Dynamism}{019}{Courier Services}
\begin{minipage}[t]{0.405\textwidth}
\SlideImage{19}{../Materials/2026/3DM_04_Dynamism_SS26.pdf}
\begin{adjustbox}{max totalsize={\textwidth}{0.43\textheight},center}
\begin{minipage}{\textwidth}\scriptsize
\NoteSection{日本語訳}\begin{itemize}
\item Courier Services（この用語は講義スライド上の専門語として扱う）
\item Pickup and delivery problem（この用語は講義スライド上の専門語として扱う）
\item Customer request service（この用語は講義スライド上の専門語として扱う）
\item Requests are accepted/rejected（この用語は講義スライド上の専門語として扱う）
\item 報酬: Total no. of customer served
\item Exogenous: Customer requests（この用語は講義スライド上の専門語として扱う）
\item 状態: Customer accepted, position of vehicle
\item 決定: Which request to accept, whom to visit next
\item 報酬: Newly accepted customers
\end{itemize}\NoteSection{試験対策ポイント}\begin{itemize}
\item dynamismは、時間の進行により状態が変わり、意思決定を繰り返す性質である。
\item static planning, dynamic planning, a priori optimization, online reoptimizationを区別する。
\end{itemize}
\end{minipage}
\end{adjustbox}
\end{minipage}\hfill
\begin{minipage}[t]{0.58\textwidth}
\begin{adjustbox}{max totalsize={\textwidth}{0.89\textheight},center}
\begin{minipage}{\textwidth}\scriptsize
\NoteSection{解説}このページでは「Courier Services」を理解する。スライド上の表現は短いが、試験では定義、具体例、モデル上の意味、手法の限界まで説明できる必要がある。\par
Dynamism は、意思決定問題が時間とともに進行し、状態が変化することを意味する。静的問題では、すべての入力を最初に与えられ、その情報に基づいて一度だけ計画を作る。動的問題では、新しい注文、交通状況、車両位置、キャンセル、故障などが途中で発生し、計画を更新する必要がある。\par
重要なのは、dynamism は stochasticity と同じではないという点である。stochasticity は未来が不確実であること、dynamism は時間とともに決定と状態が進む構造を持つことである。例えば、すべての未来注文が既知でも、時点ごとに車両位置が変わり決定を繰り返すなら動的である。逆に、将来の需要が確率的でも、最初に一括で決めて変更しないなら静的な確率計画として扱える。\par
a priori optimization は、実行前に全体計画を作り、実行中は大きく変更しない考え方である。計算しやすく、全体最適を考えやすいが、予期しない情報に弱い。online または dynamic planning は、各時点で現在状態を観測し、必要に応じて決定を変える。現実に適応しやすいが、計算時間や一貫性が問題になる。\par
配送の例では、朝に全ルートを決めるのが静的計画である。昼に新規注文が入ったり、車両が遅れたりしたときに、現在の車両位置と残り注文を見て再計画するのが動的計画である。動的計画では、今の選択が将来の柔軟性を奪わないかも考える必要がある。\par
試験答案では、静的/動的の違いを『計画を変更するか』だけでなく、状態、時点、決定、外生情報、遷移の連鎖として説明する。特に動的問題では、状態が更新され、次の決定の入力になるというループを明確に書くと満点答案に近づく。\par\NoteSection{確認問題と解答}\begin{enumerate}\item \textbf{L04-S019: 「Courier Services」の概念を、定義・具体例・モデル上の意味の順に説明せよ。}\\定義としては、dynamismは、時間の進行により状態が変わり、意思決定を繰り返す性質である。。具体例では、配送・在庫・フードデリバリーのような現実問題を用い、状態、決定、外生情報、報酬、または情報モデル/意思決定モデルのどこに対応するかを明示する。モデル上の意味として、この概念が目的関数、制約、状態表現、将来価値、または方策選択にどのような影響を与えるかを書く。試験では、用語だけを列挙せず、入力データから意思決定、さらに現実の実装へつながる流れを文章で示すことが重要である。\item \textbf{L04-S019: このスライドの考え方を、講義全体のDDDMプロセスの中に位置づけよ。}\\DDDMプロセスでは、現実世界のデータを集約して情報モデルを作り、意思決定モデルまたは方策で行動を選び、実装後に結果を観測して次の状態へ進む。このスライドの概念は、そのどこか一部だけではなく、データ表現、意思決定、将来不確実性、計算可能性のいずれかに関わる。答案では、まずこの概念がどの段階に属するかを述べ、次に他の段階へどう影響するかを書く。例えば価値関数なら将来価値を現在決定へ戻す役割、FIFOなら旅行時間情報モデルの整合性を保証する役割、CFAなら目的関数を調整して将来に良い行動を誘導する役割である。\end{enumerate}
\end{minipage}
\end{adjustbox}
\end{minipage}


\clearpage
\SlideHeader{Lecture 4 - Dynamism}{020}{Line Production}
\begin{minipage}[t]{0.405\textwidth}
\SlideImage{20}{../Materials/2026/3DM_04_Dynamism_SS26.pdf}
\begin{adjustbox}{max totalsize={\textwidth}{0.43\textheight},center}
\begin{minipage}{\textwidth}\scriptsize
\NoteSection{日本語訳}\begin{itemize}
\item Line Production（この用語は講義スライド上の専門語として扱う）
\item Scheduling problem（この用語は講義スライド上の専門語として扱う）
\item Problem horizon（この用語は講義スライド上の専門語として扱う）
\item Set of machines（この用語は講義スライド上の専門語として扱う）
\item Arriving jobs: duration, deadline, type（この用語は講義スライド上の専門語として扱う）
\item Sequence dependent set-up times（この用語は講義スライド上の専門語として扱う）
\item Exogenous: Arriving jobs（この用語は講義スライド上の専門語として扱う）
\item 状態: Current schedule, new jobs
\item 決定: Sequence, machine
\end{itemize}\NoteSection{試験対策ポイント}\begin{itemize}
\item dynamismは、時間の進行により状態が変わり、意思決定を繰り返す性質である。
\item static planning, dynamic planning, a priori optimization, online reoptimizationを区別する。
\end{itemize}
\end{minipage}
\end{adjustbox}
\end{minipage}\hfill
\begin{minipage}[t]{0.58\textwidth}
\begin{adjustbox}{max totalsize={\textwidth}{0.89\textheight},center}
\begin{minipage}{\textwidth}\scriptsize
\NoteSection{解説}このページでは「Line Production」を理解する。スライド上の表現は短いが、試験では定義、具体例、モデル上の意味、手法の限界まで説明できる必要がある。\par
Dynamism は、意思決定問題が時間とともに進行し、状態が変化することを意味する。静的問題では、すべての入力を最初に与えられ、その情報に基づいて一度だけ計画を作る。動的問題では、新しい注文、交通状況、車両位置、キャンセル、故障などが途中で発生し、計画を更新する必要がある。\par
重要なのは、dynamism は stochasticity と同じではないという点である。stochasticity は未来が不確実であること、dynamism は時間とともに決定と状態が進む構造を持つことである。例えば、すべての未来注文が既知でも、時点ごとに車両位置が変わり決定を繰り返すなら動的である。逆に、将来の需要が確率的でも、最初に一括で決めて変更しないなら静的な確率計画として扱える。\par
a priori optimization は、実行前に全体計画を作り、実行中は大きく変更しない考え方である。計算しやすく、全体最適を考えやすいが、予期しない情報に弱い。online または dynamic planning は、各時点で現在状態を観測し、必要に応じて決定を変える。現実に適応しやすいが、計算時間や一貫性が問題になる。\par
配送の例では、朝に全ルートを決めるのが静的計画である。昼に新規注文が入ったり、車両が遅れたりしたときに、現在の車両位置と残り注文を見て再計画するのが動的計画である。動的計画では、今の選択が将来の柔軟性を奪わないかも考える必要がある。\par
試験答案では、静的/動的の違いを『計画を変更するか』だけでなく、状態、時点、決定、外生情報、遷移の連鎖として説明する。特に動的問題では、状態が更新され、次の決定の入力になるというループを明確に書くと満点答案に近づく。\par\NoteSection{確認問題と解答}\begin{enumerate}\item \textbf{L04-S020: 「Line Production」の概念を、定義・具体例・モデル上の意味の順に説明せよ。}\\定義としては、dynamismは、時間の進行により状態が変わり、意思決定を繰り返す性質である。。具体例では、配送・在庫・フードデリバリーのような現実問題を用い、状態、決定、外生情報、報酬、または情報モデル/意思決定モデルのどこに対応するかを明示する。モデル上の意味として、この概念が目的関数、制約、状態表現、将来価値、または方策選択にどのような影響を与えるかを書く。試験では、用語だけを列挙せず、入力データから意思決定、さらに現実の実装へつながる流れを文章で示すことが重要である。\item \textbf{L04-S020: このスライドの考え方を、講義全体のDDDMプロセスの中に位置づけよ。}\\DDDMプロセスでは、現実世界のデータを集約して情報モデルを作り、意思決定モデルまたは方策で行動を選び、実装後に結果を観測して次の状態へ進む。このスライドの概念は、そのどこか一部だけではなく、データ表現、意思決定、将来不確実性、計算可能性のいずれかに関わる。答案では、まずこの概念がどの段階に属するかを述べ、次に他の段階へどう影響するかを書く。例えば価値関数なら将来価値を現在決定へ戻す役割、FIFOなら旅行時間情報モデルの整合性を保証する役割、CFAなら目的関数を調整して将来に良い行動を誘導する役割である。\end{enumerate}
\end{minipage}
\end{adjustbox}
\end{minipage}


\clearpage
\SlideHeader{Lecture 4 - Dynamism}{021}{Inventory Management}
\begin{minipage}[t]{0.405\textwidth}
\SlideImage{21}{../Materials/2026/3DM_04_Dynamism_SS26.pdf}
\begin{adjustbox}{max totalsize={\textwidth}{0.43\textheight},center}
\begin{minipage}{\textwidth}\scriptsize
\NoteSection{日本語訳}\begin{itemize}
\item Inventory Management（この用語は講義スライド上の専門語として扱う）
\item Inventory problem（この用語は講義スライド上の専門語として扱う）
\item Horizon in terms of days（この用語は講義スライド上の専門語として扱う）
\item Initial inventory and max capacity（この用語は講義スライド上の専門語として扱う）
\item Uncertain demand per day（この用語は講義スライド上の専門語として扱う）
\item Delivery cost（この用語は講義スライド上の専門語として扱う）
\item Holding costs per unit and day（この用語は講義スライド上の専門語として扱う）
\item If demand>inventory, difference is（この用語は講義スライド上の専門語として扱う）
\item served by (expensive) third party（この用語は講義スライド上の専門語として扱う）
\end{itemize}\NoteSection{試験対策ポイント}\begin{itemize}
\item dynamismは、時間の進行により状態が変わり、意思決定を繰り返す性質である。
\item static planning, dynamic planning, a priori optimization, online reoptimizationを区別する。
\end{itemize}
\end{minipage}
\end{adjustbox}
\end{minipage}\hfill
\begin{minipage}[t]{0.58\textwidth}
\begin{adjustbox}{max totalsize={\textwidth}{0.89\textheight},center}
\begin{minipage}{\textwidth}\scriptsize
\NoteSection{解説}このページでは「Inventory Management」を理解する。スライド上の表現は短いが、試験では定義、具体例、モデル上の意味、手法の限界まで説明できる必要がある。\par
Dynamism は、意思決定問題が時間とともに進行し、状態が変化することを意味する。静的問題では、すべての入力を最初に与えられ、その情報に基づいて一度だけ計画を作る。動的問題では、新しい注文、交通状況、車両位置、キャンセル、故障などが途中で発生し、計画を更新する必要がある。\par
重要なのは、dynamism は stochasticity と同じではないという点である。stochasticity は未来が不確実であること、dynamism は時間とともに決定と状態が進む構造を持つことである。例えば、すべての未来注文が既知でも、時点ごとに車両位置が変わり決定を繰り返すなら動的である。逆に、将来の需要が確率的でも、最初に一括で決めて変更しないなら静的な確率計画として扱える。\par
a priori optimization は、実行前に全体計画を作り、実行中は大きく変更しない考え方である。計算しやすく、全体最適を考えやすいが、予期しない情報に弱い。online または dynamic planning は、各時点で現在状態を観測し、必要に応じて決定を変える。現実に適応しやすいが、計算時間や一貫性が問題になる。\par
配送の例では、朝に全ルートを決めるのが静的計画である。昼に新規注文が入ったり、車両が遅れたりしたときに、現在の車両位置と残り注文を見て再計画するのが動的計画である。動的計画では、今の選択が将来の柔軟性を奪わないかも考える必要がある。\par
試験答案では、静的/動的の違いを『計画を変更するか』だけでなく、状態、時点、決定、外生情報、遷移の連鎖として説明する。特に動的問題では、状態が更新され、次の決定の入力になるというループを明確に書くと満点答案に近づく。\par\NoteSection{確認問題と解答}\begin{enumerate}\item \textbf{L04-S021: 「Inventory Management」の概念を、定義・具体例・モデル上の意味の順に説明せよ。}\\定義としては、dynamismは、時間の進行により状態が変わり、意思決定を繰り返す性質である。。具体例では、配送・在庫・フードデリバリーのような現実問題を用い、状態、決定、外生情報、報酬、または情報モデル/意思決定モデルのどこに対応するかを明示する。モデル上の意味として、この概念が目的関数、制約、状態表現、将来価値、または方策選択にどのような影響を与えるかを書く。試験では、用語だけを列挙せず、入力データから意思決定、さらに現実の実装へつながる流れを文章で示すことが重要である。\item \textbf{L04-S021: このスライドの考え方を、講義全体のDDDMプロセスの中に位置づけよ。}\\DDDMプロセスでは、現実世界のデータを集約して情報モデルを作り、意思決定モデルまたは方策で行動を選び、実装後に結果を観測して次の状態へ進む。このスライドの概念は、そのどこか一部だけではなく、データ表現、意思決定、将来不確実性、計算可能性のいずれかに関わる。答案では、まずこの概念がどの段階に属するかを述べ、次に他の段階へどう影響するかを書く。例えば価値関数なら将来価値を現在決定へ戻す役割、FIFOなら旅行時間情報モデルの整合性を保証する役割、CFAなら目的関数を調整して将来に良い行動を誘導する役割である。\end{enumerate}
\end{minipage}
\end{adjustbox}
\end{minipage}


\clearpage
\SlideHeader{Lecture 4 - Dynamism}{022}{Agenda}
\begin{minipage}[t]{0.405\textwidth}
\SlideImage{22}{../Materials/2026/3DM_04_Dynamism_SS26.pdf}
\begin{adjustbox}{max totalsize={\textwidth}{0.43\textheight},center}
\begin{minipage}{\textwidth}\scriptsize
\NoteSection{日本語訳}\begin{itemize}
\item アジェンダ
\item 動的性質
\item ケーススタディ: Urban Delivery and Traffic Management
\item Preparation of Modelling（この用語は講義スライド上の専門語として扱う）
\end{itemize}\NoteSection{試験対策ポイント}\begin{itemize}
\item dynamismは、時間の進行により状態が変わり、意思決定を繰り返す性質である。
\item static planning, dynamic planning, a priori optimization, online reoptimizationを区別する。
\end{itemize}
\end{minipage}
\end{adjustbox}
\end{minipage}\hfill
\begin{minipage}[t]{0.58\textwidth}
\begin{adjustbox}{max totalsize={\textwidth}{0.89\textheight},center}
\begin{minipage}{\textwidth}\scriptsize
\NoteSection{解説}このページでは「Agenda」を理解する。スライド上の表現は短いが、試験では定義、具体例、モデル上の意味、手法の限界まで説明できる必要がある。\par
Dynamism は、意思決定問題が時間とともに進行し、状態が変化することを意味する。静的問題では、すべての入力を最初に与えられ、その情報に基づいて一度だけ計画を作る。動的問題では、新しい注文、交通状況、車両位置、キャンセル、故障などが途中で発生し、計画を更新する必要がある。\par
重要なのは、dynamism は stochasticity と同じではないという点である。stochasticity は未来が不確実であること、dynamism は時間とともに決定と状態が進む構造を持つことである。例えば、すべての未来注文が既知でも、時点ごとに車両位置が変わり決定を繰り返すなら動的である。逆に、将来の需要が確率的でも、最初に一括で決めて変更しないなら静的な確率計画として扱える。\par
a priori optimization は、実行前に全体計画を作り、実行中は大きく変更しない考え方である。計算しやすく、全体最適を考えやすいが、予期しない情報に弱い。online または dynamic planning は、各時点で現在状態を観測し、必要に応じて決定を変える。現実に適応しやすいが、計算時間や一貫性が問題になる。\par
配送の例では、朝に全ルートを決めるのが静的計画である。昼に新規注文が入ったり、車両が遅れたりしたときに、現在の車両位置と残り注文を見て再計画するのが動的計画である。動的計画では、今の選択が将来の柔軟性を奪わないかも考える必要がある。\par
試験答案では、静的/動的の違いを『計画を変更するか』だけでなく、状態、時点、決定、外生情報、遷移の連鎖として説明する。特に動的問題では、状態が更新され、次の決定の入力になるというループを明確に書くと満点答案に近づく。\par\NoteSection{確認問題と解答}\begin{enumerate}\item \textbf{L04-S022: 「Agenda」の概念を、定義・具体例・モデル上の意味の順に説明せよ。}\\定義としては、dynamismは、時間の進行により状態が変わり、意思決定を繰り返す性質である。。具体例では、配送・在庫・フードデリバリーのような現実問題を用い、状態、決定、外生情報、報酬、または情報モデル/意思決定モデルのどこに対応するかを明示する。モデル上の意味として、この概念が目的関数、制約、状態表現、将来価値、または方策選択にどのような影響を与えるかを書く。試験では、用語だけを列挙せず、入力データから意思決定、さらに現実の実装へつながる流れを文章で示すことが重要である。\item \textbf{L04-S022: このスライドの考え方を、講義全体のDDDMプロセスの中に位置づけよ。}\\DDDMプロセスでは、現実世界のデータを集約して情報モデルを作り、意思決定モデルまたは方策で行動を選び、実装後に結果を観測して次の状態へ進む。このスライドの概念は、そのどこか一部だけではなく、データ表現、意思決定、将来不確実性、計算可能性のいずれかに関わる。答案では、まずこの概念がどの段階に属するかを述べ、次に他の段階へどう影響するかを書く。例えば価値関数なら将来価値を現在決定へ戻す役割、FIFOなら旅行時間情報モデルの整合性を保証する役割、CFAなら目的関数を調整して将来に良い行動を誘導する役割である。\end{enumerate}
\end{minipage}
\end{adjustbox}
\end{minipage}


\clearpage
\SlideHeader{Lecture 4 - Dynamism}{023}{Case Study}
\begin{minipage}[t]{0.405\textwidth}
\SlideImage{23}{../Materials/2026/3DM_04_Dynamism_SS26.pdf}
\begin{adjustbox}{max totalsize={\textwidth}{0.43\textheight},center}
\begin{minipage}{\textwidth}\scriptsize
\NoteSection{日本語訳}\begin{itemize}
\item Case Study（この用語は講義スライド上の専門語として扱う）
\item Urban Delivery Felix Köster, Marlin W. Ulmer, Dirk C. Mattfeld, Geir Hasle,Anticipating emission-sensitive traffic（この用語は講義スライド上の専門語として扱う）
\item management strategies for 動的 delivery routing, 交通・輸送 Research Part D: Transport and
\item Environment, Volume 62, 2018, Pages 345-361（この用語は講義スライド上の専門語として扱う）
\item Traffic Management（この用語は講義スライド上の専門語として扱う）
\item Stochastic 旅行時間s
\item Question: Can 動的 replanning save 旅行時間s?
\item berlin-（この用語は講義スライド上の専門語として扱う）
\end{itemize}\NoteSection{試験対策ポイント}\begin{itemize}
\item dynamismは、時間の進行により状態が変わり、意思決定を繰り返す性質である。
\item static planning, dynamic planning, a priori optimization, online reoptimizationを区別する。
\end{itemize}
\end{minipage}
\end{adjustbox}
\end{minipage}\hfill
\begin{minipage}[t]{0.58\textwidth}
\begin{adjustbox}{max totalsize={\textwidth}{0.89\textheight},center}
\begin{minipage}{\textwidth}\scriptsize
\NoteSection{解説}このページでは「Case Study」を理解する。スライド上の表現は短いが、試験では定義、具体例、モデル上の意味、手法の限界まで説明できる必要がある。\par
Dynamism は、意思決定問題が時間とともに進行し、状態が変化することを意味する。静的問題では、すべての入力を最初に与えられ、その情報に基づいて一度だけ計画を作る。動的問題では、新しい注文、交通状況、車両位置、キャンセル、故障などが途中で発生し、計画を更新する必要がある。\par
重要なのは、dynamism は stochasticity と同じではないという点である。stochasticity は未来が不確実であること、dynamism は時間とともに決定と状態が進む構造を持つことである。例えば、すべての未来注文が既知でも、時点ごとに車両位置が変わり決定を繰り返すなら動的である。逆に、将来の需要が確率的でも、最初に一括で決めて変更しないなら静的な確率計画として扱える。\par
a priori optimization は、実行前に全体計画を作り、実行中は大きく変更しない考え方である。計算しやすく、全体最適を考えやすいが、予期しない情報に弱い。online または dynamic planning は、各時点で現在状態を観測し、必要に応じて決定を変える。現実に適応しやすいが、計算時間や一貫性が問題になる。\par
配送の例では、朝に全ルートを決めるのが静的計画である。昼に新規注文が入ったり、車両が遅れたりしたときに、現在の車両位置と残り注文を見て再計画するのが動的計画である。動的計画では、今の選択が将来の柔軟性を奪わないかも考える必要がある。\par
試験答案では、静的/動的の違いを『計画を変更するか』だけでなく、状態、時点、決定、外生情報、遷移の連鎖として説明する。特に動的問題では、状態が更新され、次の決定の入力になるというループを明確に書くと満点答案に近づく。\par\NoteSection{確認問題と解答}\begin{enumerate}\item \textbf{L04-S023: 「Case Study」の概念を、定義・具体例・モデル上の意味の順に説明せよ。}\\定義としては、dynamismは、時間の進行により状態が変わり、意思決定を繰り返す性質である。。具体例では、配送・在庫・フードデリバリーのような現実問題を用い、状態、決定、外生情報、報酬、または情報モデル/意思決定モデルのどこに対応するかを明示する。モデル上の意味として、この概念が目的関数、制約、状態表現、将来価値、または方策選択にどのような影響を与えるかを書く。試験では、用語だけを列挙せず、入力データから意思決定、さらに現実の実装へつながる流れを文章で示すことが重要である。\item \textbf{L04-S023: このスライドの考え方を、講義全体のDDDMプロセスの中に位置づけよ。}\\DDDMプロセスでは、現実世界のデータを集約して情報モデルを作り、意思決定モデルまたは方策で行動を選び、実装後に結果を観測して次の状態へ進む。このスライドの概念は、そのどこか一部だけではなく、データ表現、意思決定、将来不確実性、計算可能性のいずれかに関わる。答案では、まずこの概念がどの段階に属するかを述べ、次に他の段階へどう影響するかを書く。例えば価値関数なら将来価値を現在決定へ戻す役割、FIFOなら旅行時間情報モデルの整合性を保証する役割、CFAなら目的関数を調整して将来に良い行動を誘導する役割である。\end{enumerate}
\end{minipage}
\end{adjustbox}
\end{minipage}


\clearpage
\SlideHeader{Lecture 4 - Dynamism}{024}{Case Study}
\begin{minipage}[t]{0.405\textwidth}
\SlideImage{24}{../Materials/2026/3DM_04_Dynamism_SS26.pdf}
\begin{adjustbox}{max totalsize={\textwidth}{0.43\textheight},center}
\begin{minipage}{\textwidth}\scriptsize
\NoteSection{日本語訳}\begin{itemize}
\item Case Study（この用語は講義スライド上の専門語として扱う）
\item Traffic Management（この用語は講義スライド上の専門語として扱う）
\item Urban Delivery Problem（この用語は講義スライド上の専門語として扱う）
\end{itemize}\NoteSection{試験対策ポイント}\begin{itemize}
\item dynamismは、時間の進行により状態が変わり、意思決定を繰り返す性質である。
\item static planning, dynamic planning, a priori optimization, online reoptimizationを区別する。
\end{itemize}
\end{minipage}
\end{adjustbox}
\end{minipage}\hfill
\begin{minipage}[t]{0.58\textwidth}
\begin{adjustbox}{max totalsize={\textwidth}{0.89\textheight},center}
\begin{minipage}{\textwidth}\scriptsize
\NoteSection{解説}このページでは「Case Study」を理解する。スライド上の表現は短いが、試験では定義、具体例、モデル上の意味、手法の限界まで説明できる必要がある。\par
Dynamism は、意思決定問題が時間とともに進行し、状態が変化することを意味する。静的問題では、すべての入力を最初に与えられ、その情報に基づいて一度だけ計画を作る。動的問題では、新しい注文、交通状況、車両位置、キャンセル、故障などが途中で発生し、計画を更新する必要がある。\par
重要なのは、dynamism は stochasticity と同じではないという点である。stochasticity は未来が不確実であること、dynamism は時間とともに決定と状態が進む構造を持つことである。例えば、すべての未来注文が既知でも、時点ごとに車両位置が変わり決定を繰り返すなら動的である。逆に、将来の需要が確率的でも、最初に一括で決めて変更しないなら静的な確率計画として扱える。\par
a priori optimization は、実行前に全体計画を作り、実行中は大きく変更しない考え方である。計算しやすく、全体最適を考えやすいが、予期しない情報に弱い。online または dynamic planning は、各時点で現在状態を観測し、必要に応じて決定を変える。現実に適応しやすいが、計算時間や一貫性が問題になる。\par
配送の例では、朝に全ルートを決めるのが静的計画である。昼に新規注文が入ったり、車両が遅れたりしたときに、現在の車両位置と残り注文を見て再計画するのが動的計画である。動的計画では、今の選択が将来の柔軟性を奪わないかも考える必要がある。\par
試験答案では、静的/動的の違いを『計画を変更するか』だけでなく、状態、時点、決定、外生情報、遷移の連鎖として説明する。特に動的問題では、状態が更新され、次の決定の入力になるというループを明確に書くと満点答案に近づく。\par\NoteSection{確認問題と解答}\begin{enumerate}\item \textbf{L04-S024: 「Case Study」の概念を、定義・具体例・モデル上の意味の順に説明せよ。}\\定義としては、dynamismは、時間の進行により状態が変わり、意思決定を繰り返す性質である。。具体例では、配送・在庫・フードデリバリーのような現実問題を用い、状態、決定、外生情報、報酬、または情報モデル/意思決定モデルのどこに対応するかを明示する。モデル上の意味として、この概念が目的関数、制約、状態表現、将来価値、または方策選択にどのような影響を与えるかを書く。試験では、用語だけを列挙せず、入力データから意思決定、さらに現実の実装へつながる流れを文章で示すことが重要である。\item \textbf{L04-S024: このスライドの考え方を、講義全体のDDDMプロセスの中に位置づけよ。}\\DDDMプロセスでは、現実世界のデータを集約して情報モデルを作り、意思決定モデルまたは方策で行動を選び、実装後に結果を観測して次の状態へ進む。このスライドの概念は、そのどこか一部だけではなく、データ表現、意思決定、将来不確実性、計算可能性のいずれかに関わる。答案では、まずこの概念がどの段階に属するかを述べ、次に他の段階へどう影響するかを書く。例えば価値関数なら将来価値を現在決定へ戻す役割、FIFOなら旅行時間情報モデルの整合性を保証する役割、CFAなら目的関数を調整して将来に良い行動を誘導する役割である。\end{enumerate}
\end{minipage}
\end{adjustbox}
\end{minipage}


\clearpage
\SlideHeader{Lecture 4 - Dynamism}{025}{Motivation: Delivery service}
\begin{minipage}[t]{0.405\textwidth}
\SlideImage{25}{../Materials/2026/3DM_04_Dynamism_SS26.pdf}
\begin{adjustbox}{max totalsize={\textwidth}{0.43\textheight},center}
\begin{minipage}{\textwidth}\scriptsize
\NoteSection{日本語訳}\begin{itemize}
\item Motivation: Delivery service（この用語は講義スライド上の専門語として扱う）
\item Parcel delivery（この用語は講義スライド上の専門語として扱う）
\item Last-Mile Delivery（この用語は講義スライド上の専門語として扱う）
\item 50\% of total transport costs（この用語は講義スライド上の専門語として扱う）
\item Vehicle fleet（この用語は講義スライド上の専門語として扱う）
\item Vehicles deliver parcels in urban districts（この用語は講義スライド上の専門語として扱う）
\item Initial 分布 of parcels to vehicles
\item Individual routing 決定
\end{itemize}\NoteSection{試験対策ポイント}\begin{itemize}
\item dynamismは、時間の進行により状態が変わり、意思決定を繰り返す性質である。
\item static planning, dynamic planning, a priori optimization, online reoptimizationを区別する。
\end{itemize}
\end{minipage}
\end{adjustbox}
\end{minipage}\hfill
\begin{minipage}[t]{0.58\textwidth}
\begin{adjustbox}{max totalsize={\textwidth}{0.89\textheight},center}
\begin{minipage}{\textwidth}\scriptsize
\NoteSection{解説}このページでは「Motivation: Delivery service」を理解する。スライド上の表現は短いが、試験では定義、具体例、モデル上の意味、手法の限界まで説明できる必要がある。\par
Dynamism は、意思決定問題が時間とともに進行し、状態が変化することを意味する。静的問題では、すべての入力を最初に与えられ、その情報に基づいて一度だけ計画を作る。動的問題では、新しい注文、交通状況、車両位置、キャンセル、故障などが途中で発生し、計画を更新する必要がある。\par
重要なのは、dynamism は stochasticity と同じではないという点である。stochasticity は未来が不確実であること、dynamism は時間とともに決定と状態が進む構造を持つことである。例えば、すべての未来注文が既知でも、時点ごとに車両位置が変わり決定を繰り返すなら動的である。逆に、将来の需要が確率的でも、最初に一括で決めて変更しないなら静的な確率計画として扱える。\par
a priori optimization は、実行前に全体計画を作り、実行中は大きく変更しない考え方である。計算しやすく、全体最適を考えやすいが、予期しない情報に弱い。online または dynamic planning は、各時点で現在状態を観測し、必要に応じて決定を変える。現実に適応しやすいが、計算時間や一貫性が問題になる。\par
配送の例では、朝に全ルートを決めるのが静的計画である。昼に新規注文が入ったり、車両が遅れたりしたときに、現在の車両位置と残り注文を見て再計画するのが動的計画である。動的計画では、今の選択が将来の柔軟性を奪わないかも考える必要がある。\par
試験答案では、静的/動的の違いを『計画を変更するか』だけでなく、状態、時点、決定、外生情報、遷移の連鎖として説明する。特に動的問題では、状態が更新され、次の決定の入力になるというループを明確に書くと満点答案に近づく。\par\NoteSection{確認問題と解答}\begin{enumerate}\item \textbf{L04-S025: 「Motivation: Delivery service」の概念を、定義・具体例・モデル上の意味の順に説明せよ。}\\定義としては、dynamismは、時間の進行により状態が変わり、意思決定を繰り返す性質である。。具体例では、配送・在庫・フードデリバリーのような現実問題を用い、状態、決定、外生情報、報酬、または情報モデル/意思決定モデルのどこに対応するかを明示する。モデル上の意味として、この概念が目的関数、制約、状態表現、将来価値、または方策選択にどのような影響を与えるかを書く。試験では、用語だけを列挙せず、入力データから意思決定、さらに現実の実装へつながる流れを文章で示すことが重要である。\item \textbf{L04-S025: このスライドの考え方を、講義全体のDDDMプロセスの中に位置づけよ。}\\DDDMプロセスでは、現実世界のデータを集約して情報モデルを作り、意思決定モデルまたは方策で行動を選び、実装後に結果を観測して次の状態へ進む。このスライドの概念は、そのどこか一部だけではなく、データ表現、意思決定、将来不確実性、計算可能性のいずれかに関わる。答案では、まずこの概念がどの段階に属するかを述べ、次に他の段階へどう影響するかを書く。例えば価値関数なら将来価値を現在決定へ戻す役割、FIFOなら旅行時間情報モデルの整合性を保証する役割、CFAなら目的関数を調整して将来に良い行動を誘導する役割である。\end{enumerate}
\end{minipage}
\end{adjustbox}
\end{minipage}


\clearpage
\SlideHeader{Lecture 4 - Dynamism}{026}{Motivation: Traffic Management}
\begin{minipage}[t]{0.405\textwidth}
\SlideImage{26}{../Materials/2026/3DM_04_Dynamism_SS26.pdf}
\begin{adjustbox}{max totalsize={\textwidth}{0.43\textheight},center}
\begin{minipage}{\textwidth}\scriptsize
\NoteSection{日本語訳}\begin{itemize}
\item Motivation: Traffic Management（この用語は講義スライド上の専門語として扱う）
\item Controls traffic flows within the city（この用語は講義スライド上の専門語として扱う）
\item 目標:
\item Enable efficient mobility（この用語は講義スライド上の専門語として扱う）
\item Social and political guidelines uvm-bs.de（この用語は講義スライド上の専門語として扱う）
\item Emissions:（この用語は講義スライド上の専門語として扱う）
\item Volume（この用語は講義スライド上の専門語として扱う）
\item Air pollution（この用語は講義スライド上の専門語として扱う）
\item EU regulation on nitrogen oxides（この用語は講義スライド上の専門語として扱う）
\end{itemize}\NoteSection{試験対策ポイント}\begin{itemize}
\item dynamismは、時間の進行により状態が変わり、意思決定を繰り返す性質である。
\item static planning, dynamic planning, a priori optimization, online reoptimizationを区別する。
\end{itemize}
\end{minipage}
\end{adjustbox}
\end{minipage}\hfill
\begin{minipage}[t]{0.58\textwidth}
\begin{adjustbox}{max totalsize={\textwidth}{0.89\textheight},center}
\begin{minipage}{\textwidth}\scriptsize
\NoteSection{解説}このページでは「Motivation: Traffic Management」を理解する。スライド上の表現は短いが、試験では定義、具体例、モデル上の意味、手法の限界まで説明できる必要がある。\par
Dynamism は、意思決定問題が時間とともに進行し、状態が変化することを意味する。静的問題では、すべての入力を最初に与えられ、その情報に基づいて一度だけ計画を作る。動的問題では、新しい注文、交通状況、車両位置、キャンセル、故障などが途中で発生し、計画を更新する必要がある。\par
重要なのは、dynamism は stochasticity と同じではないという点である。stochasticity は未来が不確実であること、dynamism は時間とともに決定と状態が進む構造を持つことである。例えば、すべての未来注文が既知でも、時点ごとに車両位置が変わり決定を繰り返すなら動的である。逆に、将来の需要が確率的でも、最初に一括で決めて変更しないなら静的な確率計画として扱える。\par
a priori optimization は、実行前に全体計画を作り、実行中は大きく変更しない考え方である。計算しやすく、全体最適を考えやすいが、予期しない情報に弱い。online または dynamic planning は、各時点で現在状態を観測し、必要に応じて決定を変える。現実に適応しやすいが、計算時間や一貫性が問題になる。\par
配送の例では、朝に全ルートを決めるのが静的計画である。昼に新規注文が入ったり、車両が遅れたりしたときに、現在の車両位置と残り注文を見て再計画するのが動的計画である。動的計画では、今の選択が将来の柔軟性を奪わないかも考える必要がある。\par
試験答案では、静的/動的の違いを『計画を変更するか』だけでなく、状態、時点、決定、外生情報、遷移の連鎖として説明する。特に動的問題では、状態が更新され、次の決定の入力になるというループを明確に書くと満点答案に近づく。\par\NoteSection{確認問題と解答}\begin{enumerate}\item \textbf{L04-S026: 「Motivation: Traffic Management」の概念を、定義・具体例・モデル上の意味の順に説明せよ。}\\定義としては、dynamismは、時間の進行により状態が変わり、意思決定を繰り返す性質である。。具体例では、配送・在庫・フードデリバリーのような現実問題を用い、状態、決定、外生情報、報酬、または情報モデル/意思決定モデルのどこに対応するかを明示する。モデル上の意味として、この概念が目的関数、制約、状態表現、将来価値、または方策選択にどのような影響を与えるかを書く。試験では、用語だけを列挙せず、入力データから意思決定、さらに現実の実装へつながる流れを文章で示すことが重要である。\item \textbf{L04-S026: このスライドの考え方を、講義全体のDDDMプロセスの中に位置づけよ。}\\DDDMプロセスでは、現実世界のデータを集約して情報モデルを作り、意思決定モデルまたは方策で行動を選び、実装後に結果を観測して次の状態へ進む。このスライドの概念は、そのどこか一部だけではなく、データ表現、意思決定、将来不確実性、計算可能性のいずれかに関わる。答案では、まずこの概念がどの段階に属するかを述べ、次に他の段階へどう影響するかを書く。例えば価値関数なら将来価値を現在決定へ戻す役割、FIFOなら旅行時間情報モデルの整合性を保証する役割、CFAなら目的関数を調整して将来に良い行動を誘導する役割である。\end{enumerate}
\end{minipage}
\end{adjustbox}
\end{minipage}


\clearpage
\SlideHeader{Lecture 4 - Dynamism}{027}{Dynamic Traffic Management}
\begin{minipage}[t]{0.405\textwidth}
\SlideImage{27}{../Materials/2026/3DM_04_Dynamism_SS26.pdf}
\begin{adjustbox}{max totalsize={\textwidth}{0.43\textheight},center}
\begin{minipage}{\textwidth}\scriptsize
\NoteSection{日本語訳}\begin{itemize}
\item Dynamic Traffic Management（この用語は講義スライド上の専門語として扱う）
\item Constant measurement of air pollution（この用語は講義スライド上の専門語として扱う）
\item If NO2 pollution reaches 60g/m3:（この用語は講義スライド上の専門語として扱う）
\item Avoidance of further inflows（この用語は講義スライド上の専門語として扱う）
\item Enabling fast outward flows（この用語は講義スライド上の専門語として扱う）
\item Tools:（この用語は講義スライド上の専門語として扱う）
\item Traffic lights wikipedia.de（この用語は講義スライド上の専門語として扱う）
\item Changes in road capacity（この用語は講義スライド上の専門語として扱う）
\item Change in speed limits（この用語は講義スライド上の専門語として扱う）
\end{itemize}\NoteSection{試験対策ポイント}\begin{itemize}
\item dynamismは、時間の進行により状態が変わり、意思決定を繰り返す性質である。
\item static planning, dynamic planning, a priori optimization, online reoptimizationを区別する。
\end{itemize}
\end{minipage}
\end{adjustbox}
\end{minipage}\hfill
\begin{minipage}[t]{0.58\textwidth}
\begin{adjustbox}{max totalsize={\textwidth}{0.89\textheight},center}
\begin{minipage}{\textwidth}\scriptsize
\NoteSection{解説}このページでは「Dynamic Traffic Management」を理解する。スライド上の表現は短いが、試験では定義、具体例、モデル上の意味、手法の限界まで説明できる必要がある。\par
Dynamism は、意思決定問題が時間とともに進行し、状態が変化することを意味する。静的問題では、すべての入力を最初に与えられ、その情報に基づいて一度だけ計画を作る。動的問題では、新しい注文、交通状況、車両位置、キャンセル、故障などが途中で発生し、計画を更新する必要がある。\par
重要なのは、dynamism は stochasticity と同じではないという点である。stochasticity は未来が不確実であること、dynamism は時間とともに決定と状態が進む構造を持つことである。例えば、すべての未来注文が既知でも、時点ごとに車両位置が変わり決定を繰り返すなら動的である。逆に、将来の需要が確率的でも、最初に一括で決めて変更しないなら静的な確率計画として扱える。\par
a priori optimization は、実行前に全体計画を作り、実行中は大きく変更しない考え方である。計算しやすく、全体最適を考えやすいが、予期しない情報に弱い。online または dynamic planning は、各時点で現在状態を観測し、必要に応じて決定を変える。現実に適応しやすいが、計算時間や一貫性が問題になる。\par
配送の例では、朝に全ルートを決めるのが静的計画である。昼に新規注文が入ったり、車両が遅れたりしたときに、現在の車両位置と残り注文を見て再計画するのが動的計画である。動的計画では、今の選択が将来の柔軟性を奪わないかも考える必要がある。\par
試験答案では、静的/動的の違いを『計画を変更するか』だけでなく、状態、時点、決定、外生情報、遷移の連鎖として説明する。特に動的問題では、状態が更新され、次の決定の入力になるというループを明確に書くと満点答案に近づく。\par\NoteSection{確認問題と解答}\begin{enumerate}\item \textbf{L04-S027: 「Dynamic Traffic Management」の概念を、定義・具体例・モデル上の意味の順に説明せよ。}\\定義としては、dynamismは、時間の進行により状態が変わり、意思決定を繰り返す性質である。。具体例では、配送・在庫・フードデリバリーのような現実問題を用い、状態、決定、外生情報、報酬、または情報モデル/意思決定モデルのどこに対応するかを明示する。モデル上の意味として、この概念が目的関数、制約、状態表現、将来価値、または方策選択にどのような影響を与えるかを書く。試験では、用語だけを列挙せず、入力データから意思決定、さらに現実の実装へつながる流れを文章で示すことが重要である。\item \textbf{L04-S027: このスライドの考え方を、講義全体のDDDMプロセスの中に位置づけよ。}\\DDDMプロセスでは、現実世界のデータを集約して情報モデルを作り、意思決定モデルまたは方策で行動を選び、実装後に結果を観測して次の状態へ進む。このスライドの概念は、そのどこか一部だけではなく、データ表現、意思決定、将来不確実性、計算可能性のいずれかに関わる。答案では、まずこの概念がどの段階に属するかを述べ、次に他の段階へどう影響するかを書く。例えば価値関数なら将来価値を現在決定へ戻す役割、FIFOなら旅行時間情報モデルの整合性を保証する役割、CFAなら目的関数を調整して将来に良い行動を誘導する役割である。\end{enumerate}
\end{minipage}
\end{adjustbox}
\end{minipage}


\clearpage
\SlideHeader{Lecture 4 - Dynamism}{028}{Traffic Management Braunschweig}
\begin{minipage}[t]{0.405\textwidth}
\SlideImage{28}{../Materials/2026/3DM_04_Dynamism_SS26.pdf}
\begin{adjustbox}{max totalsize={\textwidth}{0.43\textheight},center}
\begin{minipage}{\textwidth}\scriptsize
\NoteSection{日本語訳}\begin{itemize}
\item Traffic Management Braunschweig（この用語は講義スライド上の専門語として扱う）
\item Yearly average（この用語は講義スライド上の専門語として扱う）
\end{itemize}\NoteSection{試験対策ポイント}\begin{itemize}
\item dynamismは、時間の進行により状態が変わり、意思決定を繰り返す性質である。
\item static planning, dynamic planning, a priori optimization, online reoptimizationを区別する。
\end{itemize}
\end{minipage}
\end{adjustbox}
\end{minipage}\hfill
\begin{minipage}[t]{0.58\textwidth}
\begin{adjustbox}{max totalsize={\textwidth}{0.89\textheight},center}
\begin{minipage}{\textwidth}\scriptsize
\NoteSection{解説}このページでは「Traffic Management Braunschweig」を理解する。スライド上の表現は短いが、試験では定義、具体例、モデル上の意味、手法の限界まで説明できる必要がある。\par
Dynamism は、意思決定問題が時間とともに進行し、状態が変化することを意味する。静的問題では、すべての入力を最初に与えられ、その情報に基づいて一度だけ計画を作る。動的問題では、新しい注文、交通状況、車両位置、キャンセル、故障などが途中で発生し、計画を更新する必要がある。\par
重要なのは、dynamism は stochasticity と同じではないという点である。stochasticity は未来が不確実であること、dynamism は時間とともに決定と状態が進む構造を持つことである。例えば、すべての未来注文が既知でも、時点ごとに車両位置が変わり決定を繰り返すなら動的である。逆に、将来の需要が確率的でも、最初に一括で決めて変更しないなら静的な確率計画として扱える。\par
a priori optimization は、実行前に全体計画を作り、実行中は大きく変更しない考え方である。計算しやすく、全体最適を考えやすいが、予期しない情報に弱い。online または dynamic planning は、各時点で現在状態を観測し、必要に応じて決定を変える。現実に適応しやすいが、計算時間や一貫性が問題になる。\par
配送の例では、朝に全ルートを決めるのが静的計画である。昼に新規注文が入ったり、車両が遅れたりしたときに、現在の車両位置と残り注文を見て再計画するのが動的計画である。動的計画では、今の選択が将来の柔軟性を奪わないかも考える必要がある。\par
試験答案では、静的/動的の違いを『計画を変更するか』だけでなく、状態、時点、決定、外生情報、遷移の連鎖として説明する。特に動的問題では、状態が更新され、次の決定の入力になるというループを明確に書くと満点答案に近づく。\par\NoteSection{確認問題と解答}\begin{enumerate}\item \textbf{L04-S028: 「Traffic Management Braunschweig」の概念を、定義・具体例・モデル上の意味の順に説明せよ。}\\定義としては、dynamismは、時間の進行により状態が変わり、意思決定を繰り返す性質である。。具体例では、配送・在庫・フードデリバリーのような現実問題を用い、状態、決定、外生情報、報酬、または情報モデル/意思決定モデルのどこに対応するかを明示する。モデル上の意味として、この概念が目的関数、制約、状態表現、将来価値、または方策選択にどのような影響を与えるかを書く。試験では、用語だけを列挙せず、入力データから意思決定、さらに現実の実装へつながる流れを文章で示すことが重要である。\item \textbf{L04-S028: このスライドの考え方を、講義全体のDDDMプロセスの中に位置づけよ。}\\DDDMプロセスでは、現実世界のデータを集約して情報モデルを作り、意思決定モデルまたは方策で行動を選び、実装後に結果を観測して次の状態へ進む。このスライドの概念は、そのどこか一部だけではなく、データ表現、意思決定、将来不確実性、計算可能性のいずれかに関わる。答案では、まずこの概念がどの段階に属するかを述べ、次に他の段階へどう影響するかを書く。例えば価値関数なら将来価値を現在決定へ戻す役割、FIFOなら旅行時間情報モデルの整合性を保証する役割、CFAなら目的関数を調整して将来に良い行動を誘導する役割である。\end{enumerate}
\end{minipage}
\end{adjustbox}
\end{minipage}


\clearpage
\SlideHeader{Lecture 4 - Dynamism}{029}{Hot-Spot 1: Hildesheimer Str.}
\begin{minipage}[t]{0.405\textwidth}
\SlideImage{29}{../Materials/2026/3DM_04_Dynamism_SS26.pdf}
\begin{adjustbox}{max totalsize={\textwidth}{0.43\textheight},center}
\begin{minipage}{\textwidth}\scriptsize
\NoteSection{日本語訳}\begin{itemize}
\item Hot-Spot 1: Hildesheimer Str.（この用語は講義スライド上の専門語として扱う）
\item Limit traffic from A391（この用語は講義スライド上の専門語として扱う）
\item Suggest exits at Celler, Münchner, and（この用語は講義スライド上の専門語として扱う）
\item Hamburger Straße.（この用語は講義スライド上の専門語として扱う）
\item Reduction（この用語は講義スライド上の専門語として扱う）
\item Increase（この用語は講義スライド上の専門語として扱う）
\end{itemize}\NoteSection{試験対策ポイント}\begin{itemize}
\item dynamismは、時間の進行により状態が変わり、意思決定を繰り返す性質である。
\item static planning, dynamic planning, a priori optimization, online reoptimizationを区別する。
\end{itemize}
\end{minipage}
\end{adjustbox}
\end{minipage}\hfill
\begin{minipage}[t]{0.58\textwidth}
\begin{adjustbox}{max totalsize={\textwidth}{0.89\textheight},center}
\begin{minipage}{\textwidth}\scriptsize
\NoteSection{解説}このページでは「Hot-Spot 1: Hildesheimer Str.」を理解する。スライド上の表現は短いが、試験では定義、具体例、モデル上の意味、手法の限界まで説明できる必要がある。\par
Dynamism は、意思決定問題が時間とともに進行し、状態が変化することを意味する。静的問題では、すべての入力を最初に与えられ、その情報に基づいて一度だけ計画を作る。動的問題では、新しい注文、交通状況、車両位置、キャンセル、故障などが途中で発生し、計画を更新する必要がある。\par
重要なのは、dynamism は stochasticity と同じではないという点である。stochasticity は未来が不確実であること、dynamism は時間とともに決定と状態が進む構造を持つことである。例えば、すべての未来注文が既知でも、時点ごとに車両位置が変わり決定を繰り返すなら動的である。逆に、将来の需要が確率的でも、最初に一括で決めて変更しないなら静的な確率計画として扱える。\par
a priori optimization は、実行前に全体計画を作り、実行中は大きく変更しない考え方である。計算しやすく、全体最適を考えやすいが、予期しない情報に弱い。online または dynamic planning は、各時点で現在状態を観測し、必要に応じて決定を変える。現実に適応しやすいが、計算時間や一貫性が問題になる。\par
配送の例では、朝に全ルートを決めるのが静的計画である。昼に新規注文が入ったり、車両が遅れたりしたときに、現在の車両位置と残り注文を見て再計画するのが動的計画である。動的計画では、今の選択が将来の柔軟性を奪わないかも考える必要がある。\par
試験答案では、静的/動的の違いを『計画を変更するか』だけでなく、状態、時点、決定、外生情報、遷移の連鎖として説明する。特に動的問題では、状態が更新され、次の決定の入力になるというループを明確に書くと満点答案に近づく。\par\NoteSection{確認問題と解答}\begin{enumerate}\item \textbf{L04-S029: 「Hot-Spot 1: Hildesheimer Str.」の概念を、定義・具体例・モデル上の意味の順に説明せよ。}\\定義としては、dynamismは、時間の進行により状態が変わり、意思決定を繰り返す性質である。。具体例では、配送・在庫・フードデリバリーのような現実問題を用い、状態、決定、外生情報、報酬、または情報モデル/意思決定モデルのどこに対応するかを明示する。モデル上の意味として、この概念が目的関数、制約、状態表現、将来価値、または方策選択にどのような影響を与えるかを書く。試験では、用語だけを列挙せず、入力データから意思決定、さらに現実の実装へつながる流れを文章で示すことが重要である。\item \textbf{L04-S029: このスライドの考え方を、講義全体のDDDMプロセスの中に位置づけよ。}\\DDDMプロセスでは、現実世界のデータを集約して情報モデルを作り、意思決定モデルまたは方策で行動を選び、実装後に結果を観測して次の状態へ進む。このスライドの概念は、そのどこか一部だけではなく、データ表現、意思決定、将来不確実性、計算可能性のいずれかに関わる。答案では、まずこの概念がどの段階に属するかを述べ、次に他の段階へどう影響するかを書く。例えば価値関数なら将来価値を現在決定へ戻す役割、FIFOなら旅行時間情報モデルの整合性を保証する役割、CFAなら目的関数を調整して将来に良い行動を誘導する役割である。\end{enumerate}
\end{minipage}
\end{adjustbox}
\end{minipage}


\clearpage
\SlideHeader{Lecture 4 - Dynamism}{030}{Hot-Spot: Downtown}
\begin{minipage}[t]{0.405\textwidth}
\SlideImage{30}{../Materials/2026/3DM_04_Dynamism_SS26.pdf}
\begin{adjustbox}{max totalsize={\textwidth}{0.43\textheight},center}
\begin{minipage}{\textwidth}\scriptsize
\NoteSection{日本語訳}\begin{itemize}
\item Hot-Spot: Downtown（この用語は講義スライド上の専門語として扱う）
\item Closure Fallersleber Straße（この用語は講義スライド上の専門語として扱う）
\item Reduction（この用語は講義スライド上の専門語として扱う）
\item Increase（この用語は講義スライド上の専門語として扱う）
\end{itemize}\NoteSection{試験対策ポイント}\begin{itemize}
\item dynamismは、時間の進行により状態が変わり、意思決定を繰り返す性質である。
\item static planning, dynamic planning, a priori optimization, online reoptimizationを区別する。
\end{itemize}
\end{minipage}
\end{adjustbox}
\end{minipage}\hfill
\begin{minipage}[t]{0.58\textwidth}
\begin{adjustbox}{max totalsize={\textwidth}{0.89\textheight},center}
\begin{minipage}{\textwidth}\scriptsize
\NoteSection{解説}このページでは「Hot-Spot: Downtown」を理解する。スライド上の表現は短いが、試験では定義、具体例、モデル上の意味、手法の限界まで説明できる必要がある。\par
Dynamism は、意思決定問題が時間とともに進行し、状態が変化することを意味する。静的問題では、すべての入力を最初に与えられ、その情報に基づいて一度だけ計画を作る。動的問題では、新しい注文、交通状況、車両位置、キャンセル、故障などが途中で発生し、計画を更新する必要がある。\par
重要なのは、dynamism は stochasticity と同じではないという点である。stochasticity は未来が不確実であること、dynamism は時間とともに決定と状態が進む構造を持つことである。例えば、すべての未来注文が既知でも、時点ごとに車両位置が変わり決定を繰り返すなら動的である。逆に、将来の需要が確率的でも、最初に一括で決めて変更しないなら静的な確率計画として扱える。\par
a priori optimization は、実行前に全体計画を作り、実行中は大きく変更しない考え方である。計算しやすく、全体最適を考えやすいが、予期しない情報に弱い。online または dynamic planning は、各時点で現在状態を観測し、必要に応じて決定を変える。現実に適応しやすいが、計算時間や一貫性が問題になる。\par
配送の例では、朝に全ルートを決めるのが静的計画である。昼に新規注文が入ったり、車両が遅れたりしたときに、現在の車両位置と残り注文を見て再計画するのが動的計画である。動的計画では、今の選択が将来の柔軟性を奪わないかも考える必要がある。\par
試験答案では、静的/動的の違いを『計画を変更するか』だけでなく、状態、時点、決定、外生情報、遷移の連鎖として説明する。特に動的問題では、状態が更新され、次の決定の入力になるというループを明確に書くと満点答案に近づく。\par\NoteSection{確認問題と解答}\begin{enumerate}\item \textbf{L04-S030: 「Hot-Spot: Downtown」の概念を、定義・具体例・モデル上の意味の順に説明せよ。}\\定義としては、dynamismは、時間の進行により状態が変わり、意思決定を繰り返す性質である。。具体例では、配送・在庫・フードデリバリーのような現実問題を用い、状態、決定、外生情報、報酬、または情報モデル/意思決定モデルのどこに対応するかを明示する。モデル上の意味として、この概念が目的関数、制約、状態表現、将来価値、または方策選択にどのような影響を与えるかを書く。試験では、用語だけを列挙せず、入力データから意思決定、さらに現実の実装へつながる流れを文章で示すことが重要である。\item \textbf{L04-S030: このスライドの考え方を、講義全体のDDDMプロセスの中に位置づけよ。}\\DDDMプロセスでは、現実世界のデータを集約して情報モデルを作り、意思決定モデルまたは方策で行動を選び、実装後に結果を観測して次の状態へ進む。このスライドの概念は、そのどこか一部だけではなく、データ表現、意思決定、将来不確実性、計算可能性のいずれかに関わる。答案では、まずこの概念がどの段階に属するかを述べ、次に他の段階へどう影響するかを書く。例えば価値関数なら将来価値を現在決定へ戻す役割、FIFOなら旅行時間情報モデルの整合性を保証する役割、CFAなら目的関数を調整して将来に良い行動を誘導する役割である。\end{enumerate}
\end{minipage}
\end{adjustbox}
\end{minipage}


\clearpage
\SlideHeader{Lecture 4 - Dynamism}{031}{Hot-Spot 4: Altewiekring}
\begin{minipage}[t]{0.405\textwidth}
\SlideImage{31}{../Materials/2026/3DM_04_Dynamism_SS26.pdf}
\begin{adjustbox}{max totalsize={\textwidth}{0.43\textheight},center}
\begin{minipage}{\textwidth}\scriptsize
\NoteSection{日本語訳}\begin{itemize}
\item Hot-Spot 4: Altewiekring（この用語は講義スライド上の専門語として扱う）
\item Reduction（この用語は講義スライド上の専門語として扱う）
\item Limit left turn at Schillstraße Increase（この用語は講義スライド上の専門語として扱う）
\end{itemize}\NoteSection{試験対策ポイント}\begin{itemize}
\item dynamismは、時間の進行により状態が変わり、意思決定を繰り返す性質である。
\item static planning, dynamic planning, a priori optimization, online reoptimizationを区別する。
\end{itemize}
\end{minipage}
\end{adjustbox}
\end{minipage}\hfill
\begin{minipage}[t]{0.58\textwidth}
\begin{adjustbox}{max totalsize={\textwidth}{0.89\textheight},center}
\begin{minipage}{\textwidth}\scriptsize
\NoteSection{解説}このページでは「Hot-Spot 4: Altewiekring」を理解する。スライド上の表現は短いが、試験では定義、具体例、モデル上の意味、手法の限界まで説明できる必要がある。\par
Dynamism は、意思決定問題が時間とともに進行し、状態が変化することを意味する。静的問題では、すべての入力を最初に与えられ、その情報に基づいて一度だけ計画を作る。動的問題では、新しい注文、交通状況、車両位置、キャンセル、故障などが途中で発生し、計画を更新する必要がある。\par
重要なのは、dynamism は stochasticity と同じではないという点である。stochasticity は未来が不確実であること、dynamism は時間とともに決定と状態が進む構造を持つことである。例えば、すべての未来注文が既知でも、時点ごとに車両位置が変わり決定を繰り返すなら動的である。逆に、将来の需要が確率的でも、最初に一括で決めて変更しないなら静的な確率計画として扱える。\par
a priori optimization は、実行前に全体計画を作り、実行中は大きく変更しない考え方である。計算しやすく、全体最適を考えやすいが、予期しない情報に弱い。online または dynamic planning は、各時点で現在状態を観測し、必要に応じて決定を変える。現実に適応しやすいが、計算時間や一貫性が問題になる。\par
配送の例では、朝に全ルートを決めるのが静的計画である。昼に新規注文が入ったり、車両が遅れたりしたときに、現在の車両位置と残り注文を見て再計画するのが動的計画である。動的計画では、今の選択が将来の柔軟性を奪わないかも考える必要がある。\par
試験答案では、静的/動的の違いを『計画を変更するか』だけでなく、状態、時点、決定、外生情報、遷移の連鎖として説明する。特に動的問題では、状態が更新され、次の決定の入力になるというループを明確に書くと満点答案に近づく。\par\NoteSection{確認問題と解答}\begin{enumerate}\item \textbf{L04-S031: 「Hot-Spot 4: Altewiekring」の概念を、定義・具体例・モデル上の意味の順に説明せよ。}\\定義としては、dynamismは、時間の進行により状態が変わり、意思決定を繰り返す性質である。。具体例では、配送・在庫・フードデリバリーのような現実問題を用い、状態、決定、外生情報、報酬、または情報モデル/意思決定モデルのどこに対応するかを明示する。モデル上の意味として、この概念が目的関数、制約、状態表現、将来価値、または方策選択にどのような影響を与えるかを書く。試験では、用語だけを列挙せず、入力データから意思決定、さらに現実の実装へつながる流れを文章で示すことが重要である。\item \textbf{L04-S031: このスライドの考え方を、講義全体のDDDMプロセスの中に位置づけよ。}\\DDDMプロセスでは、現実世界のデータを集約して情報モデルを作り、意思決定モデルまたは方策で行動を選び、実装後に結果を観測して次の状態へ進む。このスライドの概念は、そのどこか一部だけではなく、データ表現、意思決定、将来不確実性、計算可能性のいずれかに関わる。答案では、まずこの概念がどの段階に属するかを述べ、次に他の段階へどう影響するかを書く。例えば価値関数なら将来価値を現在決定へ戻す役割、FIFOなら旅行時間情報モデルの整合性を保証する役割、CFAなら目的関数を調整して将来に良い行動を誘導する役割である。\end{enumerate}
\end{minipage}
\end{adjustbox}
\end{minipage}


\clearpage
\SlideHeader{Lecture 4 - Dynamism}{032}{Delivery Service}
\begin{minipage}[t]{0.405\textwidth}
\SlideImage{32}{../Materials/2026/3DM_04_Dynamism_SS26.pdf}
\begin{adjustbox}{max totalsize={\textwidth}{0.43\textheight},center}
\begin{minipage}{\textwidth}\scriptsize
\NoteSection{日本語訳}\begin{itemize}
\item Delivery Service（この用語は講義スライド上の専門語として扱う）
\item Vehicle fleet（この用語は講義スライド上の専門語として扱う）
\item Vehicles deliver packages in urban districts（この用語は講義スライド上の専門語として扱う）
\item Initial 分布 to vehicles
\item Individual routing 決定
\item 情報モデル: Stochastic 旅行時間s due to
\item traffic management 決定s
\end{itemize}\NoteSection{試験対策ポイント}\begin{itemize}
\item dynamismは、時間の進行により状態が変わり、意思決定を繰り返す性質である。
\item static planning, dynamic planning, a priori optimization, online reoptimizationを区別する。
\end{itemize}
\end{minipage}
\end{adjustbox}
\end{minipage}\hfill
\begin{minipage}[t]{0.58\textwidth}
\begin{adjustbox}{max totalsize={\textwidth}{0.89\textheight},center}
\begin{minipage}{\textwidth}\scriptsize
\NoteSection{解説}このページでは「Delivery Service」を理解する。スライド上の表現は短いが、試験では定義、具体例、モデル上の意味、手法の限界まで説明できる必要がある。\par
Dynamism は、意思決定問題が時間とともに進行し、状態が変化することを意味する。静的問題では、すべての入力を最初に与えられ、その情報に基づいて一度だけ計画を作る。動的問題では、新しい注文、交通状況、車両位置、キャンセル、故障などが途中で発生し、計画を更新する必要がある。\par
重要なのは、dynamism は stochasticity と同じではないという点である。stochasticity は未来が不確実であること、dynamism は時間とともに決定と状態が進む構造を持つことである。例えば、すべての未来注文が既知でも、時点ごとに車両位置が変わり決定を繰り返すなら動的である。逆に、将来の需要が確率的でも、最初に一括で決めて変更しないなら静的な確率計画として扱える。\par
a priori optimization は、実行前に全体計画を作り、実行中は大きく変更しない考え方である。計算しやすく、全体最適を考えやすいが、予期しない情報に弱い。online または dynamic planning は、各時点で現在状態を観測し、必要に応じて決定を変える。現実に適応しやすいが、計算時間や一貫性が問題になる。\par
配送の例では、朝に全ルートを決めるのが静的計画である。昼に新規注文が入ったり、車両が遅れたりしたときに、現在の車両位置と残り注文を見て再計画するのが動的計画である。動的計画では、今の選択が将来の柔軟性を奪わないかも考える必要がある。\par
試験答案では、静的/動的の違いを『計画を変更するか』だけでなく、状態、時点、決定、外生情報、遷移の連鎖として説明する。特に動的問題では、状態が更新され、次の決定の入力になるというループを明確に書くと満点答案に近づく。\par\NoteSection{確認問題と解答}\begin{enumerate}\item \textbf{L04-S032: 「Delivery Service」の概念を、定義・具体例・モデル上の意味の順に説明せよ。}\\定義としては、dynamismは、時間の進行により状態が変わり、意思決定を繰り返す性質である。。具体例では、配送・在庫・フードデリバリーのような現実問題を用い、状態、決定、外生情報、報酬、または情報モデル/意思決定モデルのどこに対応するかを明示する。モデル上の意味として、この概念が目的関数、制約、状態表現、将来価値、または方策選択にどのような影響を与えるかを書く。試験では、用語だけを列挙せず、入力データから意思決定、さらに現実の実装へつながる流れを文章で示すことが重要である。\item \textbf{L04-S032: このスライドの考え方を、講義全体のDDDMプロセスの中に位置づけよ。}\\DDDMプロセスでは、現実世界のデータを集約して情報モデルを作り、意思決定モデルまたは方策で行動を選び、実装後に結果を観測して次の状態へ進む。このスライドの概念は、そのどこか一部だけではなく、データ表現、意思決定、将来不確実性、計算可能性のいずれかに関わる。答案では、まずこの概念がどの段階に属するかを述べ、次に他の段階へどう影響するかを書く。例えば価値関数なら将来価値を現在決定へ戻す役割、FIFOなら旅行時間情報モデルの整合性を保証する役割、CFAなら目的関数を調整して将来に良い行動を誘導する役割である。\end{enumerate}
\end{minipage}
\end{adjustbox}
\end{minipage}


\clearpage
\SlideHeader{Lecture 4 - Dynamism}{033}{Traffic Management and Travel Times}
\begin{minipage}[t]{0.405\textwidth}
\SlideImage{33}{../Materials/2026/3DM_04_Dynamism_SS26.pdf}
\begin{adjustbox}{max totalsize={\textwidth}{0.43\textheight},center}
\begin{minipage}{\textwidth}\scriptsize
\NoteSection{日本語訳}\begin{itemize}
\item Traffic Management and 旅行時間
\item Traffic Management controls street capacities（この用語は講義スライド上の専門語として扱う）
\item Capacities impact traffic flows（この用語は講義スライド上の専門語として扱う）
\item Traffic flows define speed per street segment（この用語は講義スライド上の専門語として扱う）
\item Shortest paths and 旅行時間 matrices
\end{itemize}\NoteSection{試験対策ポイント}\begin{itemize}
\item dynamismは、時間の進行により状態が変わり、意思決定を繰り返す性質である。
\item static planning, dynamic planning, a priori optimization, online reoptimizationを区別する。
\end{itemize}
\end{minipage}
\end{adjustbox}
\end{minipage}\hfill
\begin{minipage}[t]{0.58\textwidth}
\begin{adjustbox}{max totalsize={\textwidth}{0.89\textheight},center}
\begin{minipage}{\textwidth}\scriptsize
\NoteSection{解説}このページでは「Traffic Management and Travel Times」を理解する。スライド上の表現は短いが、試験では定義、具体例、モデル上の意味、手法の限界まで説明できる必要がある。\par
Dynamism は、意思決定問題が時間とともに進行し、状態が変化することを意味する。静的問題では、すべての入力を最初に与えられ、その情報に基づいて一度だけ計画を作る。動的問題では、新しい注文、交通状況、車両位置、キャンセル、故障などが途中で発生し、計画を更新する必要がある。\par
重要なのは、dynamism は stochasticity と同じではないという点である。stochasticity は未来が不確実であること、dynamism は時間とともに決定と状態が進む構造を持つことである。例えば、すべての未来注文が既知でも、時点ごとに車両位置が変わり決定を繰り返すなら動的である。逆に、将来の需要が確率的でも、最初に一括で決めて変更しないなら静的な確率計画として扱える。\par
a priori optimization は、実行前に全体計画を作り、実行中は大きく変更しない考え方である。計算しやすく、全体最適を考えやすいが、予期しない情報に弱い。online または dynamic planning は、各時点で現在状態を観測し、必要に応じて決定を変える。現実に適応しやすいが、計算時間や一貫性が問題になる。\par
配送の例では、朝に全ルートを決めるのが静的計画である。昼に新規注文が入ったり、車両が遅れたりしたときに、現在の車両位置と残り注文を見て再計画するのが動的計画である。動的計画では、今の選択が将来の柔軟性を奪わないかも考える必要がある。\par
試験答案では、静的/動的の違いを『計画を変更するか』だけでなく、状態、時点、決定、外生情報、遷移の連鎖として説明する。特に動的問題では、状態が更新され、次の決定の入力になるというループを明確に書くと満点答案に近づく。\par\NoteSection{確認問題と解答}\begin{enumerate}\item \textbf{L04-S033: 「Traffic Management and Travel Times」の概念を、定義・具体例・モデル上の意味の順に説明せよ。}\\定義としては、dynamismは、時間の進行により状態が変わり、意思決定を繰り返す性質である。。具体例では、配送・在庫・フードデリバリーのような現実問題を用い、状態、決定、外生情報、報酬、または情報モデル/意思決定モデルのどこに対応するかを明示する。モデル上の意味として、この概念が目的関数、制約、状態表現、将来価値、または方策選択にどのような影響を与えるかを書く。試験では、用語だけを列挙せず、入力データから意思決定、さらに現実の実装へつながる流れを文章で示すことが重要である。\item \textbf{L04-S033: このスライドの考え方を、講義全体のDDDMプロセスの中に位置づけよ。}\\DDDMプロセスでは、現実世界のデータを集約して情報モデルを作り、意思決定モデルまたは方策で行動を選び、実装後に結果を観測して次の状態へ進む。このスライドの概念は、そのどこか一部だけではなく、データ表現、意思決定、将来不確実性、計算可能性のいずれかに関わる。答案では、まずこの概念がどの段階に属するかを述べ、次に他の段階へどう影響するかを書く。例えば価値関数なら将来価値を現在決定へ戻す役割、FIFOなら旅行時間情報モデルの整合性を保証する役割、CFAなら目的関数を調整して将来に良い行動を誘導する役割である。\end{enumerate}
\end{minipage}
\end{adjustbox}
\end{minipage}


\clearpage
\SlideHeader{Lecture 4 - Dynamism}{034}{Delivery service}
\begin{minipage}[t]{0.405\textwidth}
\SlideImage{34}{../Materials/2026/3DM_04_Dynamism_SS26.pdf}
\begin{adjustbox}{max totalsize={\textwidth}{0.43\textheight},center}
\begin{minipage}{\textwidth}\scriptsize
\NoteSection{日本語訳}\begin{itemize}
\item Delivery service（この用語は講義スライド上の専門語として扱う）
\item Vehicle fleet（この用語は講義スライド上の専門語として扱う）
\item Vehicles deliver packages in urban districts（この用語は講義スライド上の専門語として扱う）
\item Initial 分布 to vehicles
\item Individual routing 決定
\item 情報モデル: Stochastic 旅行時間s due to
\item traffic management 決定s
\item 意思決定モデル: Where to go next after serving a
\item district (how to adapt the tentative route)（この用語は講義スライド上の専門語として扱う）
\end{itemize}\NoteSection{試験対策ポイント}\begin{itemize}
\item dynamismは、時間の進行により状態が変わり、意思決定を繰り返す性質である。
\item static planning, dynamic planning, a priori optimization, online reoptimizationを区別する。
\end{itemize}
\end{minipage}
\end{adjustbox}
\end{minipage}\hfill
\begin{minipage}[t]{0.58\textwidth}
\begin{adjustbox}{max totalsize={\textwidth}{0.89\textheight},center}
\begin{minipage}{\textwidth}\scriptsize
\NoteSection{解説}このページでは「Delivery service」を理解する。スライド上の表現は短いが、試験では定義、具体例、モデル上の意味、手法の限界まで説明できる必要がある。\par
Dynamism は、意思決定問題が時間とともに進行し、状態が変化することを意味する。静的問題では、すべての入力を最初に与えられ、その情報に基づいて一度だけ計画を作る。動的問題では、新しい注文、交通状況、車両位置、キャンセル、故障などが途中で発生し、計画を更新する必要がある。\par
重要なのは、dynamism は stochasticity と同じではないという点である。stochasticity は未来が不確実であること、dynamism は時間とともに決定と状態が進む構造を持つことである。例えば、すべての未来注文が既知でも、時点ごとに車両位置が変わり決定を繰り返すなら動的である。逆に、将来の需要が確率的でも、最初に一括で決めて変更しないなら静的な確率計画として扱える。\par
a priori optimization は、実行前に全体計画を作り、実行中は大きく変更しない考え方である。計算しやすく、全体最適を考えやすいが、予期しない情報に弱い。online または dynamic planning は、各時点で現在状態を観測し、必要に応じて決定を変える。現実に適応しやすいが、計算時間や一貫性が問題になる。\par
配送の例では、朝に全ルートを決めるのが静的計画である。昼に新規注文が入ったり、車両が遅れたりしたときに、現在の車両位置と残り注文を見て再計画するのが動的計画である。動的計画では、今の選択が将来の柔軟性を奪わないかも考える必要がある。\par
試験答案では、静的/動的の違いを『計画を変更するか』だけでなく、状態、時点、決定、外生情報、遷移の連鎖として説明する。特に動的問題では、状態が更新され、次の決定の入力になるというループを明確に書くと満点答案に近づく。\par\NoteSection{確認問題と解答}\begin{enumerate}\item \textbf{L04-S034: 「Delivery service」の概念を、定義・具体例・モデル上の意味の順に説明せよ。}\\定義としては、dynamismは、時間の進行により状態が変わり、意思決定を繰り返す性質である。。具体例では、配送・在庫・フードデリバリーのような現実問題を用い、状態、決定、外生情報、報酬、または情報モデル/意思決定モデルのどこに対応するかを明示する。モデル上の意味として、この概念が目的関数、制約、状態表現、将来価値、または方策選択にどのような影響を与えるかを書く。試験では、用語だけを列挙せず、入力データから意思決定、さらに現実の実装へつながる流れを文章で示すことが重要である。\item \textbf{L04-S034: このスライドの考え方を、講義全体のDDDMプロセスの中に位置づけよ。}\\DDDMプロセスでは、現実世界のデータを集約して情報モデルを作り、意思決定モデルまたは方策で行動を選び、実装後に結果を観測して次の状態へ進む。このスライドの概念は、そのどこか一部だけではなく、データ表現、意思決定、将来不確実性、計算可能性のいずれかに関わる。答案では、まずこの概念がどの段階に属するかを述べ、次に他の段階へどう影響するかを書く。例えば価値関数なら将来価値を現在決定へ戻す役割、FIFOなら旅行時間情報モデルの整合性を保証する役割、CFAなら目的関数を調整して将来に良い行動を誘導する役割である。\end{enumerate}
\end{minipage}
\end{adjustbox}
\end{minipage}


\clearpage
\SlideHeader{Lecture 4 - Dynamism}{035}{Case Study Overview}
\begin{minipage}[t]{0.405\textwidth}
\SlideImage{35}{../Materials/2026/3DM_04_Dynamism_SS26.pdf}
\begin{adjustbox}{max totalsize={\textwidth}{0.43\textheight},center}
\begin{minipage}{\textwidth}\scriptsize
\NoteSection{日本語訳}\begin{itemize}
\item Case Study Overview（この用語は講義スライド上の専門語として扱う）
\item Instance Generation（この用語は講義スライド上の専門語として扱う）
\item Solution Method（この用語は講義スライド上の専門語として扱う）
\item Results（この用語は講義スライド上の専門語として扱う）
\end{itemize}\NoteSection{試験対策ポイント}\begin{itemize}
\item dynamismは、時間の進行により状態が変わり、意思決定を繰り返す性質である。
\item static planning, dynamic planning, a priori optimization, online reoptimizationを区別する。
\end{itemize}
\end{minipage}
\end{adjustbox}
\end{minipage}\hfill
\begin{minipage}[t]{0.58\textwidth}
\begin{adjustbox}{max totalsize={\textwidth}{0.89\textheight},center}
\begin{minipage}{\textwidth}\scriptsize
\NoteSection{解説}このページでは「Case Study Overview」を理解する。スライド上の表現は短いが、試験では定義、具体例、モデル上の意味、手法の限界まで説明できる必要がある。\par
Dynamism は、意思決定問題が時間とともに進行し、状態が変化することを意味する。静的問題では、すべての入力を最初に与えられ、その情報に基づいて一度だけ計画を作る。動的問題では、新しい注文、交通状況、車両位置、キャンセル、故障などが途中で発生し、計画を更新する必要がある。\par
重要なのは、dynamism は stochasticity と同じではないという点である。stochasticity は未来が不確実であること、dynamism は時間とともに決定と状態が進む構造を持つことである。例えば、すべての未来注文が既知でも、時点ごとに車両位置が変わり決定を繰り返すなら動的である。逆に、将来の需要が確率的でも、最初に一括で決めて変更しないなら静的な確率計画として扱える。\par
a priori optimization は、実行前に全体計画を作り、実行中は大きく変更しない考え方である。計算しやすく、全体最適を考えやすいが、予期しない情報に弱い。online または dynamic planning は、各時点で現在状態を観測し、必要に応じて決定を変える。現実に適応しやすいが、計算時間や一貫性が問題になる。\par
配送の例では、朝に全ルートを決めるのが静的計画である。昼に新規注文が入ったり、車両が遅れたりしたときに、現在の車両位置と残り注文を見て再計画するのが動的計画である。動的計画では、今の選択が将来の柔軟性を奪わないかも考える必要がある。\par
試験答案では、静的/動的の違いを『計画を変更するか』だけでなく、状態、時点、決定、外生情報、遷移の連鎖として説明する。特に動的問題では、状態が更新され、次の決定の入力になるというループを明確に書くと満点答案に近づく。\par\NoteSection{確認問題と解答}\begin{enumerate}\item \textbf{L04-S035: 「Case Study Overview」の概念を、定義・具体例・モデル上の意味の順に説明せよ。}\\定義としては、dynamismは、時間の進行により状態が変わり、意思決定を繰り返す性質である。。具体例では、配送・在庫・フードデリバリーのような現実問題を用い、状態、決定、外生情報、報酬、または情報モデル/意思決定モデルのどこに対応するかを明示する。モデル上の意味として、この概念が目的関数、制約、状態表現、将来価値、または方策選択にどのような影響を与えるかを書く。試験では、用語だけを列挙せず、入力データから意思決定、さらに現実の実装へつながる流れを文章で示すことが重要である。\item \textbf{L04-S035: このスライドの考え方を、講義全体のDDDMプロセスの中に位置づけよ。}\\DDDMプロセスでは、現実世界のデータを集約して情報モデルを作り、意思決定モデルまたは方策で行動を選び、実装後に結果を観測して次の状態へ進む。このスライドの概念は、そのどこか一部だけではなく、データ表現、意思決定、将来不確実性、計算可能性のいずれかに関わる。答案では、まずこの概念がどの段階に属するかを述べ、次に他の段階へどう影響するかを書く。例えば価値関数なら将来価値を現在決定へ戻す役割、FIFOなら旅行時間情報モデルの整合性を保証する役割、CFAなら目的関数を調整して将来に良い行動を誘導する役割である。\end{enumerate}
\end{minipage}
\end{adjustbox}
\end{minipage}


\clearpage
\SlideHeader{Lecture 4 - Dynamism}{036}{Instances}
\begin{minipage}[t]{0.405\textwidth}
\SlideImage{36}{../Materials/2026/3DM_04_Dynamism_SS26.pdf}
\begin{adjustbox}{max totalsize={\textwidth}{0.43\textheight},center}
\begin{minipage}{\textwidth}\scriptsize
\NoteSection{日本語訳}\begin{itemize}
\item Instances（この用語は講義スライド上の専門語として扱う）
\item 3 Hotspots in Braunschweig（この用語は講義スライド上の専門語として扱う）
\item Change of hotspot transport strategy（この用語は講義スライド上の専門語として扱う）
\item Min. 1 hour active（この用語は講義スライド上の専門語として扱う）
\item Each hotspot is independent（この用語は講義スライド上の専門語として扱う）
\item 8 Traffic strategies（この用語は講義スライド上の専門語として扱う）
\item Vehicle speed:（この用語は講義スライド上の専門語として扱う）
\item City roads: 25 km/h（この用語は講義スライド上の専門語として扱う）
\item Reduced: 10 km/h（この用語は講義スライド上の専門語として扱う）
\end{itemize}\NoteSection{試験対策ポイント}\begin{itemize}
\item dynamismは、時間の進行により状態が変わり、意思決定を繰り返す性質である。
\item static planning, dynamic planning, a priori optimization, online reoptimizationを区別する。
\end{itemize}
\end{minipage}
\end{adjustbox}
\end{minipage}\hfill
\begin{minipage}[t]{0.58\textwidth}
\begin{adjustbox}{max totalsize={\textwidth}{0.89\textheight},center}
\begin{minipage}{\textwidth}\scriptsize
\NoteSection{解説}このページでは「Instances」を理解する。スライド上の表現は短いが、試験では定義、具体例、モデル上の意味、手法の限界まで説明できる必要がある。\par
Dynamism は、意思決定問題が時間とともに進行し、状態が変化することを意味する。静的問題では、すべての入力を最初に与えられ、その情報に基づいて一度だけ計画を作る。動的問題では、新しい注文、交通状況、車両位置、キャンセル、故障などが途中で発生し、計画を更新する必要がある。\par
重要なのは、dynamism は stochasticity と同じではないという点である。stochasticity は未来が不確実であること、dynamism は時間とともに決定と状態が進む構造を持つことである。例えば、すべての未来注文が既知でも、時点ごとに車両位置が変わり決定を繰り返すなら動的である。逆に、将来の需要が確率的でも、最初に一括で決めて変更しないなら静的な確率計画として扱える。\par
a priori optimization は、実行前に全体計画を作り、実行中は大きく変更しない考え方である。計算しやすく、全体最適を考えやすいが、予期しない情報に弱い。online または dynamic planning は、各時点で現在状態を観測し、必要に応じて決定を変える。現実に適応しやすいが、計算時間や一貫性が問題になる。\par
配送の例では、朝に全ルートを決めるのが静的計画である。昼に新規注文が入ったり、車両が遅れたりしたときに、現在の車両位置と残り注文を見て再計画するのが動的計画である。動的計画では、今の選択が将来の柔軟性を奪わないかも考える必要がある。\par
試験答案では、静的/動的の違いを『計画を変更するか』だけでなく、状態、時点、決定、外生情報、遷移の連鎖として説明する。特に動的問題では、状態が更新され、次の決定の入力になるというループを明確に書くと満点答案に近づく。\par\NoteSection{確認問題と解答}\begin{enumerate}\item \textbf{L04-S036: 「Instances」の概念を、定義・具体例・モデル上の意味の順に説明せよ。}\\定義としては、dynamismは、時間の進行により状態が変わり、意思決定を繰り返す性質である。。具体例では、配送・在庫・フードデリバリーのような現実問題を用い、状態、決定、外生情報、報酬、または情報モデル/意思決定モデルのどこに対応するかを明示する。モデル上の意味として、この概念が目的関数、制約、状態表現、将来価値、または方策選択にどのような影響を与えるかを書く。試験では、用語だけを列挙せず、入力データから意思決定、さらに現実の実装へつながる流れを文章で示すことが重要である。\item \textbf{L04-S036: このスライドの考え方を、講義全体のDDDMプロセスの中に位置づけよ。}\\DDDMプロセスでは、現実世界のデータを集約して情報モデルを作り、意思決定モデルまたは方策で行動を選び、実装後に結果を観測して次の状態へ進む。このスライドの概念は、そのどこか一部だけではなく、データ表現、意思決定、将来不確実性、計算可能性のいずれかに関わる。答案では、まずこの概念がどの段階に属するかを述べ、次に他の段階へどう影響するかを書く。例えば価値関数なら将来価値を現在決定へ戻す役割、FIFOなら旅行時間情報モデルの整合性を保証する役割、CFAなら目的関数を調整して将来に良い行動を誘導する役割である。\end{enumerate}
\end{minipage}
\end{adjustbox}
\end{minipage}


\clearpage
\SlideHeader{Lecture 4 - Dynamism}{037}{Emission Data}
\begin{minipage}[t]{0.405\textwidth}
\SlideImage{37}{../Materials/2026/3DM_04_Dynamism_SS26.pdf}
\begin{adjustbox}{max totalsize={\textwidth}{0.43\textheight},center}
\begin{minipage}{\textwidth}\scriptsize
\NoteSection{日本語訳}\begin{itemize}
\item Emission Data（この用語は講義スライド上の専門語として扱う）
\item Data set of the 状態 of Lower Saxony:
\item 3 months hourly values for NO2（この用語は講義スライド上の専門語として扱う）
\item 92 days（この用語は講義スライド上の専門語として扱う）
\item Standard deviation 3-11 micrograms/m3（この用語は講義スライド上の専門語として扱う）
\item 24\% of time over 60 micrograms/m3（この用語は講義スライド上の専門語として扱う）
\item Hotspots: Emissions data from Braunschweig, Hannover and（この用語は講義スライド上の専門語として扱う）
\item Osnabrück（この用語は講義スライド上の専門語として扱う）
\item If the pollution at a hotspot is greater than 60 micrograms/m3,（この用語は講義スライド上の専門語として扱う）
\end{itemize}\NoteSection{試験対策ポイント}\begin{itemize}
\item dynamismは、時間の進行により状態が変わり、意思決定を繰り返す性質である。
\item static planning, dynamic planning, a priori optimization, online reoptimizationを区別する。
\end{itemize}
\end{minipage}
\end{adjustbox}
\end{minipage}\hfill
\begin{minipage}[t]{0.58\textwidth}
\begin{adjustbox}{max totalsize={\textwidth}{0.89\textheight},center}
\begin{minipage}{\textwidth}\scriptsize
\NoteSection{解説}このページでは「Emission Data」を理解する。スライド上の表現は短いが、試験では定義、具体例、モデル上の意味、手法の限界まで説明できる必要がある。\par
Dynamism は、意思決定問題が時間とともに進行し、状態が変化することを意味する。静的問題では、すべての入力を最初に与えられ、その情報に基づいて一度だけ計画を作る。動的問題では、新しい注文、交通状況、車両位置、キャンセル、故障などが途中で発生し、計画を更新する必要がある。\par
重要なのは、dynamism は stochasticity と同じではないという点である。stochasticity は未来が不確実であること、dynamism は時間とともに決定と状態が進む構造を持つことである。例えば、すべての未来注文が既知でも、時点ごとに車両位置が変わり決定を繰り返すなら動的である。逆に、将来の需要が確率的でも、最初に一括で決めて変更しないなら静的な確率計画として扱える。\par
a priori optimization は、実行前に全体計画を作り、実行中は大きく変更しない考え方である。計算しやすく、全体最適を考えやすいが、予期しない情報に弱い。online または dynamic planning は、各時点で現在状態を観測し、必要に応じて決定を変える。現実に適応しやすいが、計算時間や一貫性が問題になる。\par
配送の例では、朝に全ルートを決めるのが静的計画である。昼に新規注文が入ったり、車両が遅れたりしたときに、現在の車両位置と残り注文を見て再計画するのが動的計画である。動的計画では、今の選択が将来の柔軟性を奪わないかも考える必要がある。\par
試験答案では、静的/動的の違いを『計画を変更するか』だけでなく、状態、時点、決定、外生情報、遷移の連鎖として説明する。特に動的問題では、状態が更新され、次の決定の入力になるというループを明確に書くと満点答案に近づく。\par\NoteSection{確認問題と解答}\begin{enumerate}\item \textbf{L04-S037: 「Emission Data」の概念を、定義・具体例・モデル上の意味の順に説明せよ。}\\定義としては、dynamismは、時間の進行により状態が変わり、意思決定を繰り返す性質である。。具体例では、配送・在庫・フードデリバリーのような現実問題を用い、状態、決定、外生情報、報酬、または情報モデル/意思決定モデルのどこに対応するかを明示する。モデル上の意味として、この概念が目的関数、制約、状態表現、将来価値、または方策選択にどのような影響を与えるかを書く。試験では、用語だけを列挙せず、入力データから意思決定、さらに現実の実装へつながる流れを文章で示すことが重要である。\item \textbf{L04-S037: このスライドの考え方を、講義全体のDDDMプロセスの中に位置づけよ。}\\DDDMプロセスでは、現実世界のデータを集約して情報モデルを作り、意思決定モデルまたは方策で行動を選び、実装後に結果を観測して次の状態へ進む。このスライドの概念は、そのどこか一部だけではなく、データ表現、意思決定、将来不確実性、計算可能性のいずれかに関わる。答案では、まずこの概念がどの段階に属するかを述べ、次に他の段階へどう影響するかを書く。例えば価値関数なら将来価値を現在決定へ戻す役割、FIFOなら旅行時間情報モデルの整合性を保証する役割、CFAなら目的関数を調整して将来に良い行動を誘導する役割である。\end{enumerate}
\end{minipage}
\end{adjustbox}
\end{minipage}


\clearpage
\SlideHeader{Lecture 4 - Dynamism}{038}{Case Study Overview}
\begin{minipage}[t]{0.405\textwidth}
\SlideImage{38}{../Materials/2026/3DM_04_Dynamism_SS26.pdf}
\begin{adjustbox}{max totalsize={\textwidth}{0.43\textheight},center}
\begin{minipage}{\textwidth}\scriptsize
\NoteSection{日本語訳}\begin{itemize}
\item Case Study Overview（この用語は講義スライド上の専門語として扱う）
\item Instance Generation（この用語は講義スライド上の専門語として扱う）
\item Solution Method（この用語は講義スライド上の専門語として扱う）
\item Results（この用語は講義スライド上の専門語として扱う）
\end{itemize}\NoteSection{試験対策ポイント}\begin{itemize}
\item dynamismは、時間の進行により状態が変わり、意思決定を繰り返す性質である。
\item static planning, dynamic planning, a priori optimization, online reoptimizationを区別する。
\end{itemize}
\end{minipage}
\end{adjustbox}
\end{minipage}\hfill
\begin{minipage}[t]{0.58\textwidth}
\begin{adjustbox}{max totalsize={\textwidth}{0.89\textheight},center}
\begin{minipage}{\textwidth}\scriptsize
\NoteSection{解説}このページでは「Case Study Overview」を理解する。スライド上の表現は短いが、試験では定義、具体例、モデル上の意味、手法の限界まで説明できる必要がある。\par
Dynamism は、意思決定問題が時間とともに進行し、状態が変化することを意味する。静的問題では、すべての入力を最初に与えられ、その情報に基づいて一度だけ計画を作る。動的問題では、新しい注文、交通状況、車両位置、キャンセル、故障などが途中で発生し、計画を更新する必要がある。\par
重要なのは、dynamism は stochasticity と同じではないという点である。stochasticity は未来が不確実であること、dynamism は時間とともに決定と状態が進む構造を持つことである。例えば、すべての未来注文が既知でも、時点ごとに車両位置が変わり決定を繰り返すなら動的である。逆に、将来の需要が確率的でも、最初に一括で決めて変更しないなら静的な確率計画として扱える。\par
a priori optimization は、実行前に全体計画を作り、実行中は大きく変更しない考え方である。計算しやすく、全体最適を考えやすいが、予期しない情報に弱い。online または dynamic planning は、各時点で現在状態を観測し、必要に応じて決定を変える。現実に適応しやすいが、計算時間や一貫性が問題になる。\par
配送の例では、朝に全ルートを決めるのが静的計画である。昼に新規注文が入ったり、車両が遅れたりしたときに、現在の車両位置と残り注文を見て再計画するのが動的計画である。動的計画では、今の選択が将来の柔軟性を奪わないかも考える必要がある。\par
試験答案では、静的/動的の違いを『計画を変更するか』だけでなく、状態、時点、決定、外生情報、遷移の連鎖として説明する。特に動的問題では、状態が更新され、次の決定の入力になるというループを明確に書くと満点答案に近づく。\par\NoteSection{確認問題と解答}\begin{enumerate}\item \textbf{L04-S038: 「Case Study Overview」の概念を、定義・具体例・モデル上の意味の順に説明せよ。}\\定義としては、dynamismは、時間の進行により状態が変わり、意思決定を繰り返す性質である。。具体例では、配送・在庫・フードデリバリーのような現実問題を用い、状態、決定、外生情報、報酬、または情報モデル/意思決定モデルのどこに対応するかを明示する。モデル上の意味として、この概念が目的関数、制約、状態表現、将来価値、または方策選択にどのような影響を与えるかを書く。試験では、用語だけを列挙せず、入力データから意思決定、さらに現実の実装へつながる流れを文章で示すことが重要である。\item \textbf{L04-S038: このスライドの考え方を、講義全体のDDDMプロセスの中に位置づけよ。}\\DDDMプロセスでは、現実世界のデータを集約して情報モデルを作り、意思決定モデルまたは方策で行動を選び、実装後に結果を観測して次の状態へ進む。このスライドの概念は、そのどこか一部だけではなく、データ表現、意思決定、将来不確実性、計算可能性のいずれかに関わる。答案では、まずこの概念がどの段階に属するかを述べ、次に他の段階へどう影響するかを書く。例えば価値関数なら将来価値を現在決定へ戻す役割、FIFOなら旅行時間情報モデルの整合性を保証する役割、CFAなら目的関数を調整して将来に良い行動を誘導する役割である。\end{enumerate}
\end{minipage}
\end{adjustbox}
\end{minipage}


\clearpage
\SlideHeader{Lecture 4 - Dynamism}{039}{Method}
\begin{minipage}[t]{0.405\textwidth}
\SlideImage{39}{../Materials/2026/3DM_04_Dynamism_SS26.pdf}
\begin{adjustbox}{max totalsize={\textwidth}{0.43\textheight},center}
\begin{minipage}{\textwidth}\scriptsize
\NoteSection{日本語訳}\begin{itemize}
\item Method（この用語は講義スライド上の専門語として扱う）
\item Initial 分布: 車両ルーティング問題 with
\item time limit（この用語は講義スライド上の専門語として扱う）
\item For every vehicle in every 決定 状態:
\item Calculate 旅行時間s: Dijkstra-Algorithm
\item Solve the open traveling salesperson problem（この用語は講義スライド上の専門語として扱う）
\item Source: OSM（この用語は講義スライド上の専門語として扱う）
\end{itemize}\NoteSection{試験対策ポイント}\begin{itemize}
\item dynamismは、時間の進行により状態が変わり、意思決定を繰り返す性質である。
\item static planning, dynamic planning, a priori optimization, online reoptimizationを区別する。
\end{itemize}
\end{minipage}
\end{adjustbox}
\end{minipage}\hfill
\begin{minipage}[t]{0.58\textwidth}
\begin{adjustbox}{max totalsize={\textwidth}{0.89\textheight},center}
\begin{minipage}{\textwidth}\scriptsize
\NoteSection{解説}このページでは「Method」を理解する。スライド上の表現は短いが、試験では定義、具体例、モデル上の意味、手法の限界まで説明できる必要がある。\par
Dynamism は、意思決定問題が時間とともに進行し、状態が変化することを意味する。静的問題では、すべての入力を最初に与えられ、その情報に基づいて一度だけ計画を作る。動的問題では、新しい注文、交通状況、車両位置、キャンセル、故障などが途中で発生し、計画を更新する必要がある。\par
重要なのは、dynamism は stochasticity と同じではないという点である。stochasticity は未来が不確実であること、dynamism は時間とともに決定と状態が進む構造を持つことである。例えば、すべての未来注文が既知でも、時点ごとに車両位置が変わり決定を繰り返すなら動的である。逆に、将来の需要が確率的でも、最初に一括で決めて変更しないなら静的な確率計画として扱える。\par
a priori optimization は、実行前に全体計画を作り、実行中は大きく変更しない考え方である。計算しやすく、全体最適を考えやすいが、予期しない情報に弱い。online または dynamic planning は、各時点で現在状態を観測し、必要に応じて決定を変える。現実に適応しやすいが、計算時間や一貫性が問題になる。\par
配送の例では、朝に全ルートを決めるのが静的計画である。昼に新規注文が入ったり、車両が遅れたりしたときに、現在の車両位置と残り注文を見て再計画するのが動的計画である。動的計画では、今の選択が将来の柔軟性を奪わないかも考える必要がある。\par
試験答案では、静的/動的の違いを『計画を変更するか』だけでなく、状態、時点、決定、外生情報、遷移の連鎖として説明する。特に動的問題では、状態が更新され、次の決定の入力になるというループを明確に書くと満点答案に近づく。\par\NoteSection{確認問題と解答}\begin{enumerate}\item \textbf{L04-S039: 「Method」の概念を、定義・具体例・モデル上の意味の順に説明せよ。}\\定義としては、dynamismは、時間の進行により状態が変わり、意思決定を繰り返す性質である。。具体例では、配送・在庫・フードデリバリーのような現実問題を用い、状態、決定、外生情報、報酬、または情報モデル/意思決定モデルのどこに対応するかを明示する。モデル上の意味として、この概念が目的関数、制約、状態表現、将来価値、または方策選択にどのような影響を与えるかを書く。試験では、用語だけを列挙せず、入力データから意思決定、さらに現実の実装へつながる流れを文章で示すことが重要である。\item \textbf{L04-S039: このスライドの考え方を、講義全体のDDDMプロセスの中に位置づけよ。}\\DDDMプロセスでは、現実世界のデータを集約して情報モデルを作り、意思決定モデルまたは方策で行動を選び、実装後に結果を観測して次の状態へ進む。このスライドの概念は、そのどこか一部だけではなく、データ表現、意思決定、将来不確実性、計算可能性のいずれかに関わる。答案では、まずこの概念がどの段階に属するかを述べ、次に他の段階へどう影響するかを書く。例えば価値関数なら将来価値を現在決定へ戻す役割、FIFOなら旅行時間情報モデルの整合性を保証する役割、CFAなら目的関数を調整して将来に良い行動を誘導する役割である。\end{enumerate}
\end{minipage}
\end{adjustbox}
\end{minipage}


\clearpage
\SlideHeader{Lecture 4 - Dynamism}{040}{Results}
\begin{minipage}[t]{0.405\textwidth}
\SlideImage{40}{../Materials/2026/3DM_04_Dynamism_SS26.pdf}
\begin{adjustbox}{max totalsize={\textwidth}{0.43\textheight},center}
\begin{minipage}{\textwidth}\scriptsize
\NoteSection{日本語訳}\begin{itemize}
\item Results（この用語は講義スライド上の専門語として扱う）
\item 1000 Days（この用語は講義スライド上の専門語として扱う）
\item Average sum of 旅行時間s
\item Reduction by 3.5\%（この用語は講義スライド上の専門語として扱う）
\item Travel time Strategy（この用語は講義スライド上の専門語として扱う）
\end{itemize}\NoteSection{試験対策ポイント}\begin{itemize}
\item dynamismは、時間の進行により状態が変わり、意思決定を繰り返す性質である。
\item static planning, dynamic planning, a priori optimization, online reoptimizationを区別する。
\end{itemize}
\end{minipage}
\end{adjustbox}
\end{minipage}\hfill
\begin{minipage}[t]{0.58\textwidth}
\begin{adjustbox}{max totalsize={\textwidth}{0.89\textheight},center}
\begin{minipage}{\textwidth}\scriptsize
\NoteSection{解説}このページでは「Results」を理解する。スライド上の表現は短いが、試験では定義、具体例、モデル上の意味、手法の限界まで説明できる必要がある。\par
Dynamism は、意思決定問題が時間とともに進行し、状態が変化することを意味する。静的問題では、すべての入力を最初に与えられ、その情報に基づいて一度だけ計画を作る。動的問題では、新しい注文、交通状況、車両位置、キャンセル、故障などが途中で発生し、計画を更新する必要がある。\par
重要なのは、dynamism は stochasticity と同じではないという点である。stochasticity は未来が不確実であること、dynamism は時間とともに決定と状態が進む構造を持つことである。例えば、すべての未来注文が既知でも、時点ごとに車両位置が変わり決定を繰り返すなら動的である。逆に、将来の需要が確率的でも、最初に一括で決めて変更しないなら静的な確率計画として扱える。\par
a priori optimization は、実行前に全体計画を作り、実行中は大きく変更しない考え方である。計算しやすく、全体最適を考えやすいが、予期しない情報に弱い。online または dynamic planning は、各時点で現在状態を観測し、必要に応じて決定を変える。現実に適応しやすいが、計算時間や一貫性が問題になる。\par
配送の例では、朝に全ルートを決めるのが静的計画である。昼に新規注文が入ったり、車両が遅れたりしたときに、現在の車両位置と残り注文を見て再計画するのが動的計画である。動的計画では、今の選択が将来の柔軟性を奪わないかも考える必要がある。\par
試験答案では、静的/動的の違いを『計画を変更するか』だけでなく、状態、時点、決定、外生情報、遷移の連鎖として説明する。特に動的問題では、状態が更新され、次の決定の入力になるというループを明確に書くと満点答案に近づく。\par\NoteSection{確認問題と解答}\begin{enumerate}\item \textbf{L04-S040: 「Results」の概念を、定義・具体例・モデル上の意味の順に説明せよ。}\\定義としては、dynamismは、時間の進行により状態が変わり、意思決定を繰り返す性質である。。具体例では、配送・在庫・フードデリバリーのような現実問題を用い、状態、決定、外生情報、報酬、または情報モデル/意思決定モデルのどこに対応するかを明示する。モデル上の意味として、この概念が目的関数、制約、状態表現、将来価値、または方策選択にどのような影響を与えるかを書く。試験では、用語だけを列挙せず、入力データから意思決定、さらに現実の実装へつながる流れを文章で示すことが重要である。\item \textbf{L04-S040: このスライドの考え方を、講義全体のDDDMプロセスの中に位置づけよ。}\\DDDMプロセスでは、現実世界のデータを集約して情報モデルを作り、意思決定モデルまたは方策で行動を選び、実装後に結果を観測して次の状態へ進む。このスライドの概念は、そのどこか一部だけではなく、データ表現、意思決定、将来不確実性、計算可能性のいずれかに関わる。答案では、まずこの概念がどの段階に属するかを述べ、次に他の段階へどう影響するかを書く。例えば価値関数なら将来価値を現在決定へ戻す役割、FIFOなら旅行時間情報モデルの整合性を保証する役割、CFAなら目的関数を調整して将来に良い行動を誘導する役割である。\end{enumerate}
\end{minipage}
\end{adjustbox}
\end{minipage}


\clearpage
\SlideHeader{Lecture 4 - Dynamism}{041}{Street Usage}
\begin{minipage}[t]{0.405\textwidth}
\SlideImage{41}{../Materials/2026/3DM_04_Dynamism_SS26.pdf}
\begin{adjustbox}{max totalsize={\textwidth}{0.43\textheight},center}
\begin{minipage}{\textwidth}\scriptsize
\NoteSection{日本語訳}\begin{itemize}
\item Street Usage（この用語は講義スライド上の専門語として扱う）
\item Static Dynamic（この用語は講義スライド上の専門語として扱う）
\end{itemize}\NoteSection{試験対策ポイント}\begin{itemize}
\item dynamismは、時間の進行により状態が変わり、意思決定を繰り返す性質である。
\item static planning, dynamic planning, a priori optimization, online reoptimizationを区別する。
\end{itemize}
\end{minipage}
\end{adjustbox}
\end{minipage}\hfill
\begin{minipage}[t]{0.58\textwidth}
\begin{adjustbox}{max totalsize={\textwidth}{0.89\textheight},center}
\begin{minipage}{\textwidth}\scriptsize
\NoteSection{解説}このページでは「Street Usage」を理解する。スライド上の表現は短いが、試験では定義、具体例、モデル上の意味、手法の限界まで説明できる必要がある。\par
Dynamism は、意思決定問題が時間とともに進行し、状態が変化することを意味する。静的問題では、すべての入力を最初に与えられ、その情報に基づいて一度だけ計画を作る。動的問題では、新しい注文、交通状況、車両位置、キャンセル、故障などが途中で発生し、計画を更新する必要がある。\par
重要なのは、dynamism は stochasticity と同じではないという点である。stochasticity は未来が不確実であること、dynamism は時間とともに決定と状態が進む構造を持つことである。例えば、すべての未来注文が既知でも、時点ごとに車両位置が変わり決定を繰り返すなら動的である。逆に、将来の需要が確率的でも、最初に一括で決めて変更しないなら静的な確率計画として扱える。\par
a priori optimization は、実行前に全体計画を作り、実行中は大きく変更しない考え方である。計算しやすく、全体最適を考えやすいが、予期しない情報に弱い。online または dynamic planning は、各時点で現在状態を観測し、必要に応じて決定を変える。現実に適応しやすいが、計算時間や一貫性が問題になる。\par
配送の例では、朝に全ルートを決めるのが静的計画である。昼に新規注文が入ったり、車両が遅れたりしたときに、現在の車両位置と残り注文を見て再計画するのが動的計画である。動的計画では、今の選択が将来の柔軟性を奪わないかも考える必要がある。\par
試験答案では、静的/動的の違いを『計画を変更するか』だけでなく、状態、時点、決定、外生情報、遷移の連鎖として説明する。特に動的問題では、状態が更新され、次の決定の入力になるというループを明確に書くと満点答案に近づく。\par\NoteSection{確認問題と解答}\begin{enumerate}\item \textbf{L04-S041: 「Street Usage」の概念を、定義・具体例・モデル上の意味の順に説明せよ。}\\定義としては、dynamismは、時間の進行により状態が変わり、意思決定を繰り返す性質である。。具体例では、配送・在庫・フードデリバリーのような現実問題を用い、状態、決定、外生情報、報酬、または情報モデル/意思決定モデルのどこに対応するかを明示する。モデル上の意味として、この概念が目的関数、制約、状態表現、将来価値、または方策選択にどのような影響を与えるかを書く。試験では、用語だけを列挙せず、入力データから意思決定、さらに現実の実装へつながる流れを文章で示すことが重要である。\item \textbf{L04-S041: このスライドの考え方を、講義全体のDDDMプロセスの中に位置づけよ。}\\DDDMプロセスでは、現実世界のデータを集約して情報モデルを作り、意思決定モデルまたは方策で行動を選び、実装後に結果を観測して次の状態へ進む。このスライドの概念は、そのどこか一部だけではなく、データ表現、意思決定、将来不確実性、計算可能性のいずれかに関わる。答案では、まずこの概念がどの段階に属するかを述べ、次に他の段階へどう影響するかを書く。例えば価値関数なら将来価値を現在決定へ戻す役割、FIFOなら旅行時間情報モデルの整合性を保証する役割、CFAなら目的関数を調整して将来に良い行動を誘導する役割である。\end{enumerate}
\end{minipage}
\end{adjustbox}
\end{minipage}


\clearpage
\SlideHeader{Lecture 4 - Dynamism}{042}{Agenda}
\begin{minipage}[t]{0.405\textwidth}
\SlideImage{42}{../Materials/2026/3DM_04_Dynamism_SS26.pdf}
\begin{adjustbox}{max totalsize={\textwidth}{0.43\textheight},center}
\begin{minipage}{\textwidth}\scriptsize
\NoteSection{日本語訳}\begin{itemize}
\item アジェンダ
\item 動的性質
\item ケーススタディ: Urban Delivery and Traffic Management
\item Preparation of Modelling（この用語は講義スライド上の専門語として扱う）
\end{itemize}\NoteSection{試験対策ポイント}\begin{itemize}
\item dynamismは、時間の進行により状態が変わり、意思決定を繰り返す性質である。
\item static planning, dynamic planning, a priori optimization, online reoptimizationを区別する。
\end{itemize}
\end{minipage}
\end{adjustbox}
\end{minipage}\hfill
\begin{minipage}[t]{0.58\textwidth}
\begin{adjustbox}{max totalsize={\textwidth}{0.89\textheight},center}
\begin{minipage}{\textwidth}\scriptsize
\NoteSection{解説}このページでは「Agenda」を理解する。スライド上の表現は短いが、試験では定義、具体例、モデル上の意味、手法の限界まで説明できる必要がある。\par
Dynamism は、意思決定問題が時間とともに進行し、状態が変化することを意味する。静的問題では、すべての入力を最初に与えられ、その情報に基づいて一度だけ計画を作る。動的問題では、新しい注文、交通状況、車両位置、キャンセル、故障などが途中で発生し、計画を更新する必要がある。\par
重要なのは、dynamism は stochasticity と同じではないという点である。stochasticity は未来が不確実であること、dynamism は時間とともに決定と状態が進む構造を持つことである。例えば、すべての未来注文が既知でも、時点ごとに車両位置が変わり決定を繰り返すなら動的である。逆に、将来の需要が確率的でも、最初に一括で決めて変更しないなら静的な確率計画として扱える。\par
a priori optimization は、実行前に全体計画を作り、実行中は大きく変更しない考え方である。計算しやすく、全体最適を考えやすいが、予期しない情報に弱い。online または dynamic planning は、各時点で現在状態を観測し、必要に応じて決定を変える。現実に適応しやすいが、計算時間や一貫性が問題になる。\par
配送の例では、朝に全ルートを決めるのが静的計画である。昼に新規注文が入ったり、車両が遅れたりしたときに、現在の車両位置と残り注文を見て再計画するのが動的計画である。動的計画では、今の選択が将来の柔軟性を奪わないかも考える必要がある。\par
試験答案では、静的/動的の違いを『計画を変更するか』だけでなく、状態、時点、決定、外生情報、遷移の連鎖として説明する。特に動的問題では、状態が更新され、次の決定の入力になるというループを明確に書くと満点答案に近づく。\par\NoteSection{確認問題と解答}\begin{enumerate}\item \textbf{L04-S042: 「Agenda」の概念を、定義・具体例・モデル上の意味の順に説明せよ。}\\定義としては、dynamismは、時間の進行により状態が変わり、意思決定を繰り返す性質である。。具体例では、配送・在庫・フードデリバリーのような現実問題を用い、状態、決定、外生情報、報酬、または情報モデル/意思決定モデルのどこに対応するかを明示する。モデル上の意味として、この概念が目的関数、制約、状態表現、将来価値、または方策選択にどのような影響を与えるかを書く。試験では、用語だけを列挙せず、入力データから意思決定、さらに現実の実装へつながる流れを文章で示すことが重要である。\item \textbf{L04-S042: このスライドの考え方を、講義全体のDDDMプロセスの中に位置づけよ。}\\DDDMプロセスでは、現実世界のデータを集約して情報モデルを作り、意思決定モデルまたは方策で行動を選び、実装後に結果を観測して次の状態へ進む。このスライドの概念は、そのどこか一部だけではなく、データ表現、意思決定、将来不確実性、計算可能性のいずれかに関わる。答案では、まずこの概念がどの段階に属するかを述べ、次に他の段階へどう影響するかを書く。例えば価値関数なら将来価値を現在決定へ戻す役割、FIFOなら旅行時間情報モデルの整合性を保証する役割、CFAなら目的関数を調整して将来に良い行動を誘導する役割である。\end{enumerate}
\end{minipage}
\end{adjustbox}
\end{minipage}


\clearpage
\SlideHeader{Lecture 4 - Dynamism}{043}{Dynamic Decision Process}
\begin{minipage}[t]{0.405\textwidth}
\SlideImage{43}{../Materials/2026/3DM_04_Dynamism_SS26.pdf}
\begin{adjustbox}{max totalsize={\textwidth}{0.43\textheight},center}
\begin{minipage}{\textwidth}\scriptsize
\NoteSection{日本語訳}\begin{itemize}
\item 動的意思決定プロセス
\item Two processes:（この用語は講義スライド上の専門語として扱う）
\item Exogeneous: Changing information (確率的ity)
\item Endogenous: Adaption of plans (dynamism)（この用語は講義スライド上の専門語として扱う）
\item Meisel (2011)（この用語は講義スライド上の専門語として扱う）
\item → A model should capture 確率的ity and dynamism (over time)
\end{itemize}\NoteSection{試験対策ポイント}\begin{itemize}
\item dynamismは、時間の進行により状態が変わり、意思決定を繰り返す性質である。
\item static planning, dynamic planning, a priori optimization, online reoptimizationを区別する。
\end{itemize}
\end{minipage}
\end{adjustbox}
\end{minipage}\hfill
\begin{minipage}[t]{0.58\textwidth}
\begin{adjustbox}{max totalsize={\textwidth}{0.89\textheight},center}
\begin{minipage}{\textwidth}\scriptsize
\NoteSection{解説}このページでは「Dynamic Decision Process」を理解する。スライド上の表現は短いが、試験では定義、具体例、モデル上の意味、手法の限界まで説明できる必要がある。\par
Dynamism は、意思決定問題が時間とともに進行し、状態が変化することを意味する。静的問題では、すべての入力を最初に与えられ、その情報に基づいて一度だけ計画を作る。動的問題では、新しい注文、交通状況、車両位置、キャンセル、故障などが途中で発生し、計画を更新する必要がある。\par
重要なのは、dynamism は stochasticity と同じではないという点である。stochasticity は未来が不確実であること、dynamism は時間とともに決定と状態が進む構造を持つことである。例えば、すべての未来注文が既知でも、時点ごとに車両位置が変わり決定を繰り返すなら動的である。逆に、将来の需要が確率的でも、最初に一括で決めて変更しないなら静的な確率計画として扱える。\par
a priori optimization は、実行前に全体計画を作り、実行中は大きく変更しない考え方である。計算しやすく、全体最適を考えやすいが、予期しない情報に弱い。online または dynamic planning は、各時点で現在状態を観測し、必要に応じて決定を変える。現実に適応しやすいが、計算時間や一貫性が問題になる。\par
配送の例では、朝に全ルートを決めるのが静的計画である。昼に新規注文が入ったり、車両が遅れたりしたときに、現在の車両位置と残り注文を見て再計画するのが動的計画である。動的計画では、今の選択が将来の柔軟性を奪わないかも考える必要がある。\par
試験答案では、静的/動的の違いを『計画を変更するか』だけでなく、状態、時点、決定、外生情報、遷移の連鎖として説明する。特に動的問題では、状態が更新され、次の決定の入力になるというループを明確に書くと満点答案に近づく。\par\NoteSection{確認問題と解答}\begin{enumerate}\item \textbf{L04-S043: 「Dynamic Decision Process」の概念を、定義・具体例・モデル上の意味の順に説明せよ。}\\定義としては、dynamismは、時間の進行により状態が変わり、意思決定を繰り返す性質である。。具体例では、配送・在庫・フードデリバリーのような現実問題を用い、状態、決定、外生情報、報酬、または情報モデル/意思決定モデルのどこに対応するかを明示する。モデル上の意味として、この概念が目的関数、制約、状態表現、将来価値、または方策選択にどのような影響を与えるかを書く。試験では、用語だけを列挙せず、入力データから意思決定、さらに現実の実装へつながる流れを文章で示すことが重要である。\item \textbf{L04-S043: このスライドの考え方を、講義全体のDDDMプロセスの中に位置づけよ。}\\DDDMプロセスでは、現実世界のデータを集約して情報モデルを作り、意思決定モデルまたは方策で行動を選び、実装後に結果を観測して次の状態へ進む。このスライドの概念は、そのどこか一部だけではなく、データ表現、意思決定、将来不確実性、計算可能性のいずれかに関わる。答案では、まずこの概念がどの段階に属するかを述べ、次に他の段階へどう影響するかを書く。例えば価値関数なら将来価値を現在決定へ戻す役割、FIFOなら旅行時間情報モデルの整合性を保証する役割、CFAなら目的関数を調整して将来に良い行動を誘導する役割である。\end{enumerate}
\end{minipage}
\end{adjustbox}
\end{minipage}


\clearpage
\SlideHeader{Lecture 4 - Dynamism}{044}{“Standard” Static Vehicle Routing Problem}
\begin{minipage}[t]{0.405\textwidth}
\SlideImage{44}{../Materials/2026/3DM_04_Dynamism_SS26.pdf}
\begin{adjustbox}{max totalsize={\textwidth}{0.43\textheight},center}
\begin{minipage}{\textwidth}\scriptsize
\NoteSection{日本語訳}\begin{itemize}
\item “Standard” Static 車両ルーティング問題
\item Depot, customers, vehicles（この用語は講義スライド上の専門語として扱う）
\item Constraint: capacity（この用語は講義スライド上の専門語として扱う）
\item 目的: 最小化する 旅行時間
\item 決定s
\item Assignment of customers to vehicles（この用語は講義スライド上の専門語として扱う）
\item Sequence of customers per vehicle（この用語は講義スライド上の専門語として扱う）
\item Modeling: mixed-integer programming（この用語は講義スライド上の専門語として扱う）
\item Solution: set of routes（この用語は講義スライド上の専門語として扱う）
\end{itemize}\NoteSection{試験対策ポイント}\begin{itemize}
\item dynamismは、時間の進行により状態が変わり、意思決定を繰り返す性質である。
\item static planning, dynamic planning, a priori optimization, online reoptimizationを区別する。
\end{itemize}
\end{minipage}
\end{adjustbox}
\end{minipage}\hfill
\begin{minipage}[t]{0.58\textwidth}
\begin{adjustbox}{max totalsize={\textwidth}{0.89\textheight},center}
\begin{minipage}{\textwidth}\scriptsize
\NoteSection{解説}このページでは「“Standard” Static Vehicle Routing Problem」を理解する。スライド上の表現は短いが、試験では定義、具体例、モデル上の意味、手法の限界まで説明できる必要がある。\par
Dynamism は、意思決定問題が時間とともに進行し、状態が変化することを意味する。静的問題では、すべての入力を最初に与えられ、その情報に基づいて一度だけ計画を作る。動的問題では、新しい注文、交通状況、車両位置、キャンセル、故障などが途中で発生し、計画を更新する必要がある。\par
重要なのは、dynamism は stochasticity と同じではないという点である。stochasticity は未来が不確実であること、dynamism は時間とともに決定と状態が進む構造を持つことである。例えば、すべての未来注文が既知でも、時点ごとに車両位置が変わり決定を繰り返すなら動的である。逆に、将来の需要が確率的でも、最初に一括で決めて変更しないなら静的な確率計画として扱える。\par
a priori optimization は、実行前に全体計画を作り、実行中は大きく変更しない考え方である。計算しやすく、全体最適を考えやすいが、予期しない情報に弱い。online または dynamic planning は、各時点で現在状態を観測し、必要に応じて決定を変える。現実に適応しやすいが、計算時間や一貫性が問題になる。\par
配送の例では、朝に全ルートを決めるのが静的計画である。昼に新規注文が入ったり、車両が遅れたりしたときに、現在の車両位置と残り注文を見て再計画するのが動的計画である。動的計画では、今の選択が将来の柔軟性を奪わないかも考える必要がある。\par
試験答案では、静的/動的の違いを『計画を変更するか』だけでなく、状態、時点、決定、外生情報、遷移の連鎖として説明する。特に動的問題では、状態が更新され、次の決定の入力になるというループを明確に書くと満点答案に近づく。\par\NoteSection{確認問題と解答}\begin{enumerate}\item \textbf{L04-S044: 「“Standard” Static Vehicle Routing Problem」の概念を、定義・具体例・モデル上の意味の順に説明せよ。}\\定義としては、dynamismは、時間の進行により状態が変わり、意思決定を繰り返す性質である。。具体例では、配送・在庫・フードデリバリーのような現実問題を用い、状態、決定、外生情報、報酬、または情報モデル/意思決定モデルのどこに対応するかを明示する。モデル上の意味として、この概念が目的関数、制約、状態表現、将来価値、または方策選択にどのような影響を与えるかを書く。試験では、用語だけを列挙せず、入力データから意思決定、さらに現実の実装へつながる流れを文章で示すことが重要である。\item \textbf{L04-S044: このスライドの考え方を、講義全体のDDDMプロセスの中に位置づけよ。}\\DDDMプロセスでは、現実世界のデータを集約して情報モデルを作り、意思決定モデルまたは方策で行動を選び、実装後に結果を観測して次の状態へ進む。このスライドの概念は、そのどこか一部だけではなく、データ表現、意思決定、将来不確実性、計算可能性のいずれかに関わる。答案では、まずこの概念がどの段階に属するかを述べ、次に他の段階へどう影響するかを書く。例えば価値関数なら将来価値を現在決定へ戻す役割、FIFOなら旅行時間情報モデルの整合性を保証する役割、CFAなら目的関数を調整して将来に良い行動を誘導する役割である。\end{enumerate}
\end{minipage}
\end{adjustbox}
\end{minipage}


\clearpage
\SlideHeader{Lecture 4 - Dynamism}{045}{Markov Decision Process}
\begin{minipage}[t]{0.405\textwidth}
\SlideImage{45}{../Materials/2026/3DM_04_Dynamism_SS26.pdf}
\begin{adjustbox}{max totalsize={\textwidth}{0.43\textheight},center}
\begin{minipage}{\textwidth}\scriptsize
\NoteSection{日本語訳}\begin{itemize}
\item マルコフ意思決定過程
\item 𝑆𝑘 𝑥 𝑆𝑘𝑥 𝜔𝑘 𝑆𝑘+1
\item 𝑅 𝑆𝑘 , 𝑥
\item Components:（この用語は講義スライド上の専門語として扱う）
\item 決定 point 𝑘 (either periodic or event triggered)
\item 決定 状態 𝑆𝑘 ∈ 𝐒
\item 決定s 𝑥 ∈ 𝑋 (e.g., with Mixed-Integer Programming)
\item 報酬 𝑅 𝑆𝑘 , 𝑥 ∈ ℝ (or costs as negative 報酬)
\item Post-決定 状態 𝑆𝑘𝑥 ∈ 𝐒 × 𝑋
\end{itemize}\NoteSection{試験対策ポイント}\begin{itemize}
\item dynamismは、時間の進行により状態が変わり、意思決定を繰り返す性質である。
\item static planning, dynamic planning, a priori optimization, online reoptimizationを区別する。
\end{itemize}
\end{minipage}
\end{adjustbox}
\end{minipage}\hfill
\begin{minipage}[t]{0.58\textwidth}
\begin{adjustbox}{max totalsize={\textwidth}{0.89\textheight},center}
\begin{minipage}{\textwidth}\scriptsize
\NoteSection{解説}このページでは「Markov Decision Process」を理解する。スライド上の表現は短いが、試験では定義、具体例、モデル上の意味、手法の限界まで説明できる必要がある。\par
Dynamism は、意思決定問題が時間とともに進行し、状態が変化することを意味する。静的問題では、すべての入力を最初に与えられ、その情報に基づいて一度だけ計画を作る。動的問題では、新しい注文、交通状況、車両位置、キャンセル、故障などが途中で発生し、計画を更新する必要がある。\par
重要なのは、dynamism は stochasticity と同じではないという点である。stochasticity は未来が不確実であること、dynamism は時間とともに決定と状態が進む構造を持つことである。例えば、すべての未来注文が既知でも、時点ごとに車両位置が変わり決定を繰り返すなら動的である。逆に、将来の需要が確率的でも、最初に一括で決めて変更しないなら静的な確率計画として扱える。\par
a priori optimization は、実行前に全体計画を作り、実行中は大きく変更しない考え方である。計算しやすく、全体最適を考えやすいが、予期しない情報に弱い。online または dynamic planning は、各時点で現在状態を観測し、必要に応じて決定を変える。現実に適応しやすいが、計算時間や一貫性が問題になる。\par
配送の例では、朝に全ルートを決めるのが静的計画である。昼に新規注文が入ったり、車両が遅れたりしたときに、現在の車両位置と残り注文を見て再計画するのが動的計画である。動的計画では、今の選択が将来の柔軟性を奪わないかも考える必要がある。\par
試験答案では、静的/動的の違いを『計画を変更するか』だけでなく、状態、時点、決定、外生情報、遷移の連鎖として説明する。特に動的問題では、状態が更新され、次の決定の入力になるというループを明確に書くと満点答案に近づく。\par\NoteSection{確認問題と解答}\begin{enumerate}\item \textbf{L04-S045: 「Markov Decision Process」の概念を、定義・具体例・モデル上の意味の順に説明せよ。}\\定義としては、dynamismは、時間の進行により状態が変わり、意思決定を繰り返す性質である。。具体例では、配送・在庫・フードデリバリーのような現実問題を用い、状態、決定、外生情報、報酬、または情報モデル/意思決定モデルのどこに対応するかを明示する。モデル上の意味として、この概念が目的関数、制約、状態表現、将来価値、または方策選択にどのような影響を与えるかを書く。試験では、用語だけを列挙せず、入力データから意思決定、さらに現実の実装へつながる流れを文章で示すことが重要である。\item \textbf{L04-S045: このスライドの考え方を、講義全体のDDDMプロセスの中に位置づけよ。}\\DDDMプロセスでは、現実世界のデータを集約して情報モデルを作り、意思決定モデルまたは方策で行動を選び、実装後に結果を観測して次の状態へ進む。このスライドの概念は、そのどこか一部だけではなく、データ表現、意思決定、将来不確実性、計算可能性のいずれかに関わる。答案では、まずこの概念がどの段階に属するかを述べ、次に他の段階へどう影響するかを書く。例えば価値関数なら将来価値を現在決定へ戻す役割、FIFOなら旅行時間情報モデルの整合性を保証する役割、CFAなら目的関数を調整して将来に良い行動を誘導する役割である。\end{enumerate}
\end{minipage}
\end{adjustbox}
\end{minipage}


\clearpage
\SlideHeader{Lecture 4 - Dynamism}{046}{Example: Parcel Service}
\begin{minipage}[t]{0.405\textwidth}
\SlideImage{46}{../Materials/2026/3DM_04_Dynamism_SS26.pdf}
\begin{adjustbox}{max totalsize={\textwidth}{0.43\textheight},center}
\begin{minipage}{\textwidth}\scriptsize
\NoteSection{日本語訳}\begin{itemize}
\item 例: Parcel Service
\item 𝑥 𝜔𝑘
\item 𝑆𝑘 𝑆𝑘𝑥 𝑆𝑘+1
\item 𝑅 𝑆𝑘 , 𝑥 = 30
\item 𝑡 = 60 D 𝑡 = 60 D 𝑡 = 90 D
\item 3 3 3
\item 2 30 2 2
\item 1 1
\end{itemize}\NoteSection{試験対策ポイント}\begin{itemize}
\item dynamismは、時間の進行により状態が変わり、意思決定を繰り返す性質である。
\item static planning, dynamic planning, a priori optimization, online reoptimizationを区別する。
\end{itemize}
\end{minipage}
\end{adjustbox}
\end{minipage}\hfill
\begin{minipage}[t]{0.58\textwidth}
\begin{adjustbox}{max totalsize={\textwidth}{0.89\textheight},center}
\begin{minipage}{\textwidth}\scriptsize
\NoteSection{解説}このページでは「Example: Parcel Service」を理解する。スライド上の表現は短いが、試験では定義、具体例、モデル上の意味、手法の限界まで説明できる必要がある。\par
Dynamism は、意思決定問題が時間とともに進行し、状態が変化することを意味する。静的問題では、すべての入力を最初に与えられ、その情報に基づいて一度だけ計画を作る。動的問題では、新しい注文、交通状況、車両位置、キャンセル、故障などが途中で発生し、計画を更新する必要がある。\par
重要なのは、dynamism は stochasticity と同じではないという点である。stochasticity は未来が不確実であること、dynamism は時間とともに決定と状態が進む構造を持つことである。例えば、すべての未来注文が既知でも、時点ごとに車両位置が変わり決定を繰り返すなら動的である。逆に、将来の需要が確率的でも、最初に一括で決めて変更しないなら静的な確率計画として扱える。\par
a priori optimization は、実行前に全体計画を作り、実行中は大きく変更しない考え方である。計算しやすく、全体最適を考えやすいが、予期しない情報に弱い。online または dynamic planning は、各時点で現在状態を観測し、必要に応じて決定を変える。現実に適応しやすいが、計算時間や一貫性が問題になる。\par
配送の例では、朝に全ルートを決めるのが静的計画である。昼に新規注文が入ったり、車両が遅れたりしたときに、現在の車両位置と残り注文を見て再計画するのが動的計画である。動的計画では、今の選択が将来の柔軟性を奪わないかも考える必要がある。\par
試験答案では、静的/動的の違いを『計画を変更するか』だけでなく、状態、時点、決定、外生情報、遷移の連鎖として説明する。特に動的問題では、状態が更新され、次の決定の入力になるというループを明確に書くと満点答案に近づく。\par\NoteSection{確認問題と解答}\begin{enumerate}\item \textbf{L04-S046: 「Example: Parcel Service」の概念を、定義・具体例・モデル上の意味の順に説明せよ。}\\定義としては、dynamismは、時間の進行により状態が変わり、意思決定を繰り返す性質である。。具体例では、配送・在庫・フードデリバリーのような現実問題を用い、状態、決定、外生情報、報酬、または情報モデル/意思決定モデルのどこに対応するかを明示する。モデル上の意味として、この概念が目的関数、制約、状態表現、将来価値、または方策選択にどのような影響を与えるかを書く。試験では、用語だけを列挙せず、入力データから意思決定、さらに現実の実装へつながる流れを文章で示すことが重要である。\item \textbf{L04-S046: このスライドの考え方を、講義全体のDDDMプロセスの中に位置づけよ。}\\DDDMプロセスでは、現実世界のデータを集約して情報モデルを作り、意思決定モデルまたは方策で行動を選び、実装後に結果を観測して次の状態へ進む。このスライドの概念は、そのどこか一部だけではなく、データ表現、意思決定、将来不確実性、計算可能性のいずれかに関わる。答案では、まずこの概念がどの段階に属するかを述べ、次に他の段階へどう影響するかを書く。例えば価値関数なら将来価値を現在決定へ戻す役割、FIFOなら旅行時間情報モデルの整合性を保証する役割、CFAなら目的関数を調整して将来に良い行動を誘導する役割である。\end{enumerate}
\end{minipage}
\end{adjustbox}
\end{minipage}


\clearpage
\SlideHeader{Lecture 4 - Dynamism}{047}{Problem Modeling: Markov Decision Process}
\begin{minipage}[t]{0.405\textwidth}
\SlideImage{47}{../Materials/2026/3DM_04_Dynamism_SS26.pdf}
\begin{adjustbox}{max totalsize={\textwidth}{0.43\textheight},center}
\begin{minipage}{\textwidth}\scriptsize
\NoteSection{日本語訳}\begin{itemize}
\item Problem Modeling: マルコフ意思決定過程
\item 𝑥 𝜔𝑘
\item 𝑆𝑘 𝑆𝑘𝑥 𝑆𝑘+1
\item 𝑅 𝑆𝑘 , 𝑥
\item S0 SK
\item SK
\item Initial 状態 𝑆0
\item Final 状態 𝑆𝐾
\end{itemize}\NoteSection{試験対策ポイント}\begin{itemize}
\item dynamismは、時間の進行により状態が変わり、意思決定を繰り返す性質である。
\item static planning, dynamic planning, a priori optimization, online reoptimizationを区別する。
\end{itemize}
\end{minipage}
\end{adjustbox}
\end{minipage}\hfill
\begin{minipage}[t]{0.58\textwidth}
\begin{adjustbox}{max totalsize={\textwidth}{0.89\textheight},center}
\begin{minipage}{\textwidth}\scriptsize
\NoteSection{解説}このページでは「Problem Modeling: Markov Decision Process」を理解する。スライド上の表現は短いが、試験では定義、具体例、モデル上の意味、手法の限界まで説明できる必要がある。\par
Dynamism は、意思決定問題が時間とともに進行し、状態が変化することを意味する。静的問題では、すべての入力を最初に与えられ、その情報に基づいて一度だけ計画を作る。動的問題では、新しい注文、交通状況、車両位置、キャンセル、故障などが途中で発生し、計画を更新する必要がある。\par
重要なのは、dynamism は stochasticity と同じではないという点である。stochasticity は未来が不確実であること、dynamism は時間とともに決定と状態が進む構造を持つことである。例えば、すべての未来注文が既知でも、時点ごとに車両位置が変わり決定を繰り返すなら動的である。逆に、将来の需要が確率的でも、最初に一括で決めて変更しないなら静的な確率計画として扱える。\par
a priori optimization は、実行前に全体計画を作り、実行中は大きく変更しない考え方である。計算しやすく、全体最適を考えやすいが、予期しない情報に弱い。online または dynamic planning は、各時点で現在状態を観測し、必要に応じて決定を変える。現実に適応しやすいが、計算時間や一貫性が問題になる。\par
配送の例では、朝に全ルートを決めるのが静的計画である。昼に新規注文が入ったり、車両が遅れたりしたときに、現在の車両位置と残り注文を見て再計画するのが動的計画である。動的計画では、今の選択が将来の柔軟性を奪わないかも考える必要がある。\par
試験答案では、静的/動的の違いを『計画を変更するか』だけでなく、状態、時点、決定、外生情報、遷移の連鎖として説明する。特に動的問題では、状態が更新され、次の決定の入力になるというループを明確に書くと満点答案に近づく。\par\NoteSection{確認問題と解答}\begin{enumerate}\item \textbf{L04-S047: 「Problem Modeling: Markov Decision Process」の概念を、定義・具体例・モデル上の意味の順に説明せよ。}\\定義としては、dynamismは、時間の進行により状態が変わり、意思決定を繰り返す性質である。。具体例では、配送・在庫・フードデリバリーのような現実問題を用い、状態、決定、外生情報、報酬、または情報モデル/意思決定モデルのどこに対応するかを明示する。モデル上の意味として、この概念が目的関数、制約、状態表現、将来価値、または方策選択にどのような影響を与えるかを書く。試験では、用語だけを列挙せず、入力データから意思決定、さらに現実の実装へつながる流れを文章で示すことが重要である。\item \textbf{L04-S047: このスライドの考え方を、講義全体のDDDMプロセスの中に位置づけよ。}\\DDDMプロセスでは、現実世界のデータを集約して情報モデルを作り、意思決定モデルまたは方策で行動を選び、実装後に結果を観測して次の状態へ進む。このスライドの概念は、そのどこか一部だけではなく、データ表現、意思決定、将来不確実性、計算可能性のいずれかに関わる。答案では、まずこの概念がどの段階に属するかを述べ、次に他の段階へどう影響するかを書く。例えば価値関数なら将来価値を現在決定へ戻す役割、FIFOなら旅行時間情報モデルの整合性を保証する役割、CFAなら目的関数を調整して将来に良い行動を誘導する役割である。\end{enumerate}
\end{minipage}
\end{adjustbox}
\end{minipage}


\clearpage
\SlideHeader{Lecture 4 - Dynamism}{048}{Solution: Decision Policy}
\begin{minipage}[t]{0.405\textwidth}
\SlideImage{48}{../Materials/2026/3DM_04_Dynamism_SS26.pdf}
\begin{adjustbox}{max totalsize={\textwidth}{0.43\textheight},center}
\begin{minipage}{\textwidth}\scriptsize
\NoteSection{日本語訳}\begin{itemize}
\item Solution: 決定 方策
\item A 決定 方策 assigns a 決定 to every 状態
\item S0 SK
\item SK
\item 目的: find optimal 方策 𝜋 ∗ maximizing 期待される overall 報酬s
\end{itemize}\NoteSection{試験対策ポイント}\begin{itemize}
\item dynamismは、時間の進行により状態が変わり、意思決定を繰り返す性質である。
\item static planning, dynamic planning, a priori optimization, online reoptimizationを区別する。
\end{itemize}
\end{minipage}
\end{adjustbox}
\end{minipage}\hfill
\begin{minipage}[t]{0.58\textwidth}
\begin{adjustbox}{max totalsize={\textwidth}{0.89\textheight},center}
\begin{minipage}{\textwidth}\scriptsize
\NoteSection{解説}このページでは「Solution: Decision Policy」を理解する。スライド上の表現は短いが、試験では定義、具体例、モデル上の意味、手法の限界まで説明できる必要がある。\par
Dynamism は、意思決定問題が時間とともに進行し、状態が変化することを意味する。静的問題では、すべての入力を最初に与えられ、その情報に基づいて一度だけ計画を作る。動的問題では、新しい注文、交通状況、車両位置、キャンセル、故障などが途中で発生し、計画を更新する必要がある。\par
重要なのは、dynamism は stochasticity と同じではないという点である。stochasticity は未来が不確実であること、dynamism は時間とともに決定と状態が進む構造を持つことである。例えば、すべての未来注文が既知でも、時点ごとに車両位置が変わり決定を繰り返すなら動的である。逆に、将来の需要が確率的でも、最初に一括で決めて変更しないなら静的な確率計画として扱える。\par
a priori optimization は、実行前に全体計画を作り、実行中は大きく変更しない考え方である。計算しやすく、全体最適を考えやすいが、予期しない情報に弱い。online または dynamic planning は、各時点で現在状態を観測し、必要に応じて決定を変える。現実に適応しやすいが、計算時間や一貫性が問題になる。\par
配送の例では、朝に全ルートを決めるのが静的計画である。昼に新規注文が入ったり、車両が遅れたりしたときに、現在の車両位置と残り注文を見て再計画するのが動的計画である。動的計画では、今の選択が将来の柔軟性を奪わないかも考える必要がある。\par
試験答案では、静的/動的の違いを『計画を変更するか』だけでなく、状態、時点、決定、外生情報、遷移の連鎖として説明する。特に動的問題では、状態が更新され、次の決定の入力になるというループを明確に書くと満点答案に近づく。\par\NoteSection{確認問題と解答}\begin{enumerate}\item \textbf{L04-S048: 「Solution: Decision Policy」の概念を、定義・具体例・モデル上の意味の順に説明せよ。}\\定義としては、dynamismは、時間の進行により状態が変わり、意思決定を繰り返す性質である。。具体例では、配送・在庫・フードデリバリーのような現実問題を用い、状態、決定、外生情報、報酬、または情報モデル/意思決定モデルのどこに対応するかを明示する。モデル上の意味として、この概念が目的関数、制約、状態表現、将来価値、または方策選択にどのような影響を与えるかを書く。試験では、用語だけを列挙せず、入力データから意思決定、さらに現実の実装へつながる流れを文章で示すことが重要である。\item \textbf{L04-S048: このスライドの考え方を、講義全体のDDDMプロセスの中に位置づけよ。}\\DDDMプロセスでは、現実世界のデータを集約して情報モデルを作り、意思決定モデルまたは方策で行動を選び、実装後に結果を観測して次の状態へ進む。このスライドの概念は、そのどこか一部だけではなく、データ表現、意思決定、将来不確実性、計算可能性のいずれかに関わる。答案では、まずこの概念がどの段階に属するかを述べ、次に他の段階へどう影響するかを書く。例えば価値関数なら将来価値を現在決定へ戻す役割、FIFOなら旅行時間情報モデルの整合性を保証する役割、CFAなら目的関数を調整して将来に良い行動を誘導する役割である。\end{enumerate}
\end{minipage}
\end{adjustbox}
\end{minipage}


\clearpage
\SlideHeader{Lecture 4 - Dynamism}{049}{Optimal Solution}
\begin{minipage}[t]{0.405\textwidth}
\SlideImage{49}{../Materials/2026/3DM_04_Dynamism_SS26.pdf}
\begin{adjustbox}{max totalsize={\textwidth}{0.43\textheight},center}
\begin{minipage}{\textwidth}\scriptsize
\NoteSection{日本語訳}\begin{itemize}
\item 最適解
\item An optimal solution 𝜋 ∗ satisfies:（この用語は講義スライド上の専門語として扱う）
\item 𝐾
\item ∗
\item 𝜋 ∗ = arg max𝜋∈Π 𝔼 ෍ 𝑅(𝑆𝑖 , 𝑋𝑖𝜋 (𝑆𝑖 ))|𝑆0
\item 𝑖=0
\item 方策 𝜋 ∗ represents the 方策 𝜋 ∈ Π, maximizing the 期待される 報酬
\item Given initial 状態 𝑆0
\item 𝑋𝑘𝜋 𝑆𝑘 indicates the 決定 made by 方策 𝜋 in 状態 𝑆𝑘
\end{itemize}\NoteSection{試験対策ポイント}\begin{itemize}
\item dynamismは、時間の進行により状態が変わり、意思決定を繰り返す性質である。
\item static planning, dynamic planning, a priori optimization, online reoptimizationを区別する。
\end{itemize}
\end{minipage}
\end{adjustbox}
\end{minipage}\hfill
\begin{minipage}[t]{0.58\textwidth}
\begin{adjustbox}{max totalsize={\textwidth}{0.89\textheight},center}
\begin{minipage}{\textwidth}\scriptsize
\NoteSection{解説}このページでは「Optimal Solution」を理解する。スライド上の表現は短いが、試験では定義、具体例、モデル上の意味、手法の限界まで説明できる必要がある。\par
Dynamism は、意思決定問題が時間とともに進行し、状態が変化することを意味する。静的問題では、すべての入力を最初に与えられ、その情報に基づいて一度だけ計画を作る。動的問題では、新しい注文、交通状況、車両位置、キャンセル、故障などが途中で発生し、計画を更新する必要がある。\par
重要なのは、dynamism は stochasticity と同じではないという点である。stochasticity は未来が不確実であること、dynamism は時間とともに決定と状態が進む構造を持つことである。例えば、すべての未来注文が既知でも、時点ごとに車両位置が変わり決定を繰り返すなら動的である。逆に、将来の需要が確率的でも、最初に一括で決めて変更しないなら静的な確率計画として扱える。\par
a priori optimization は、実行前に全体計画を作り、実行中は大きく変更しない考え方である。計算しやすく、全体最適を考えやすいが、予期しない情報に弱い。online または dynamic planning は、各時点で現在状態を観測し、必要に応じて決定を変える。現実に適応しやすいが、計算時間や一貫性が問題になる。\par
配送の例では、朝に全ルートを決めるのが静的計画である。昼に新規注文が入ったり、車両が遅れたりしたときに、現在の車両位置と残り注文を見て再計画するのが動的計画である。動的計画では、今の選択が将来の柔軟性を奪わないかも考える必要がある。\par
試験答案では、静的/動的の違いを『計画を変更するか』だけでなく、状態、時点、決定、外生情報、遷移の連鎖として説明する。特に動的問題では、状態が更新され、次の決定の入力になるというループを明確に書くと満点答案に近づく。\par\NoteSection{確認問題と解答}\begin{enumerate}\item \textbf{L04-S049: 「Optimal Solution」の概念を、定義・具体例・モデル上の意味の順に説明せよ。}\\定義としては、dynamismは、時間の進行により状態が変わり、意思決定を繰り返す性質である。。具体例では、配送・在庫・フードデリバリーのような現実問題を用い、状態、決定、外生情報、報酬、または情報モデル/意思決定モデルのどこに対応するかを明示する。モデル上の意味として、この概念が目的関数、制約、状態表現、将来価値、または方策選択にどのような影響を与えるかを書く。試験では、用語だけを列挙せず、入力データから意思決定、さらに現実の実装へつながる流れを文章で示すことが重要である。\item \textbf{L04-S049: このスライドの考え方を、講義全体のDDDMプロセスの中に位置づけよ。}\\DDDMプロセスでは、現実世界のデータを集約して情報モデルを作り、意思決定モデルまたは方策で行動を選び、実装後に結果を観測して次の状態へ進む。このスライドの概念は、そのどこか一部だけではなく、データ表現、意思決定、将来不確実性、計算可能性のいずれかに関わる。答案では、まずこの概念がどの段階に属するかを述べ、次に他の段階へどう影響するかを書く。例えば価値関数なら将来価値を現在決定へ戻す役割、FIFOなら旅行時間情報モデルの整合性を保証する役割、CFAなら目的関数を調整して将来に良い行動を誘導する役割である。\end{enumerate}
\end{minipage}
\end{adjustbox}
\end{minipage}


\clearpage
\SlideHeader{Lecture 4 - Dynamism}{050}{Summary}
\begin{minipage}[t]{0.405\textwidth}
\SlideImage{50}{../Materials/2026/3DM_04_Dynamism_SS26.pdf}
\begin{adjustbox}{max totalsize={\textwidth}{0.43\textheight},center}
\begin{minipage}{\textwidth}\scriptsize
\NoteSection{日本語訳}\begin{itemize}
\item まとめ
\item Motivation und Definition: 動的性質
\item Case Study: Dynamic re-planning reduces cost（この用語は講義スライド上の専門語として扱う）
\item Modelling: マルコフ意思決定過程
\end{itemize}\NoteSection{試験対策ポイント}\begin{itemize}
\item dynamismは、時間の進行により状態が変わり、意思決定を繰り返す性質である。
\item static planning, dynamic planning, a priori optimization, online reoptimizationを区別する。
\end{itemize}
\end{minipage}
\end{adjustbox}
\end{minipage}\hfill
\begin{minipage}[t]{0.58\textwidth}
\begin{adjustbox}{max totalsize={\textwidth}{0.89\textheight},center}
\begin{minipage}{\textwidth}\scriptsize
\NoteSection{解説}このページでは「Summary」を理解する。スライド上の表現は短いが、試験では定義、具体例、モデル上の意味、手法の限界まで説明できる必要がある。\par
Dynamism は、意思決定問題が時間とともに進行し、状態が変化することを意味する。静的問題では、すべての入力を最初に与えられ、その情報に基づいて一度だけ計画を作る。動的問題では、新しい注文、交通状況、車両位置、キャンセル、故障などが途中で発生し、計画を更新する必要がある。\par
重要なのは、dynamism は stochasticity と同じではないという点である。stochasticity は未来が不確実であること、dynamism は時間とともに決定と状態が進む構造を持つことである。例えば、すべての未来注文が既知でも、時点ごとに車両位置が変わり決定を繰り返すなら動的である。逆に、将来の需要が確率的でも、最初に一括で決めて変更しないなら静的な確率計画として扱える。\par
a priori optimization は、実行前に全体計画を作り、実行中は大きく変更しない考え方である。計算しやすく、全体最適を考えやすいが、予期しない情報に弱い。online または dynamic planning は、各時点で現在状態を観測し、必要に応じて決定を変える。現実に適応しやすいが、計算時間や一貫性が問題になる。\par
配送の例では、朝に全ルートを決めるのが静的計画である。昼に新規注文が入ったり、車両が遅れたりしたときに、現在の車両位置と残り注文を見て再計画するのが動的計画である。動的計画では、今の選択が将来の柔軟性を奪わないかも考える必要がある。\par
試験答案では、静的/動的の違いを『計画を変更するか』だけでなく、状態、時点、決定、外生情報、遷移の連鎖として説明する。特に動的問題では、状態が更新され、次の決定の入力になるというループを明確に書くと満点答案に近づく。\par\NoteSection{確認問題と解答}\begin{enumerate}\item \textbf{L04-S050: 「Summary」の概念を、定義・具体例・モデル上の意味の順に説明せよ。}\\定義としては、dynamismは、時間の進行により状態が変わり、意思決定を繰り返す性質である。。具体例では、配送・在庫・フードデリバリーのような現実問題を用い、状態、決定、外生情報、報酬、または情報モデル/意思決定モデルのどこに対応するかを明示する。モデル上の意味として、この概念が目的関数、制約、状態表現、将来価値、または方策選択にどのような影響を与えるかを書く。試験では、用語だけを列挙せず、入力データから意思決定、さらに現実の実装へつながる流れを文章で示すことが重要である。\item \textbf{L04-S050: このスライドの考え方を、講義全体のDDDMプロセスの中に位置づけよ。}\\DDDMプロセスでは、現実世界のデータを集約して情報モデルを作り、意思決定モデルまたは方策で行動を選び、実装後に結果を観測して次の状態へ進む。このスライドの概念は、そのどこか一部だけではなく、データ表現、意思決定、将来不確実性、計算可能性のいずれかに関わる。答案では、まずこの概念がどの段階に属するかを述べ、次に他の段階へどう影響するかを書く。例えば価値関数なら将来価値を現在決定へ戻す役割、FIFOなら旅行時間情報モデルの整合性を保証する役割、CFAなら目的関数を調整して将来に良い行動を誘導する役割である。\end{enumerate}
\end{minipage}
\end{adjustbox}
\end{minipage}


\clearpage
\SlideHeader{Lecture 4 - Dynamism}{051}{Outlook}
\begin{minipage}[t]{0.405\textwidth}
\SlideImage{51}{../Materials/2026/3DM_04_Dynamism_SS26.pdf}
\begin{adjustbox}{max totalsize={\textwidth}{0.43\textheight},center}
\begin{minipage}{\textwidth}\scriptsize
\NoteSection{日本語訳}\begin{itemize}
\item Outlook（この用語は講義スライド上の専門語として扱う）
\item Anticipation: Integration of future developments in 決定 making
\item Next lecture:（この用語は講義スライド上の専門語として扱う）
\item マルコフ意思決定過程
\item Optimal solutions（この用語は講義スライド上の専門語として扱う）
\item S0 SK
\item SK
\end{itemize}\NoteSection{試験対策ポイント}\begin{itemize}
\item dynamismは、時間の進行により状態が変わり、意思決定を繰り返す性質である。
\item static planning, dynamic planning, a priori optimization, online reoptimizationを区別する。
\end{itemize}
\end{minipage}
\end{adjustbox}
\end{minipage}\hfill
\begin{minipage}[t]{0.58\textwidth}
\begin{adjustbox}{max totalsize={\textwidth}{0.89\textheight},center}
\begin{minipage}{\textwidth}\scriptsize
\NoteSection{解説}このページでは「Outlook」を理解する。スライド上の表現は短いが、試験では定義、具体例、モデル上の意味、手法の限界まで説明できる必要がある。\par
Dynamism は、意思決定問題が時間とともに進行し、状態が変化することを意味する。静的問題では、すべての入力を最初に与えられ、その情報に基づいて一度だけ計画を作る。動的問題では、新しい注文、交通状況、車両位置、キャンセル、故障などが途中で発生し、計画を更新する必要がある。\par
重要なのは、dynamism は stochasticity と同じではないという点である。stochasticity は未来が不確実であること、dynamism は時間とともに決定と状態が進む構造を持つことである。例えば、すべての未来注文が既知でも、時点ごとに車両位置が変わり決定を繰り返すなら動的である。逆に、将来の需要が確率的でも、最初に一括で決めて変更しないなら静的な確率計画として扱える。\par
a priori optimization は、実行前に全体計画を作り、実行中は大きく変更しない考え方である。計算しやすく、全体最適を考えやすいが、予期しない情報に弱い。online または dynamic planning は、各時点で現在状態を観測し、必要に応じて決定を変える。現実に適応しやすいが、計算時間や一貫性が問題になる。\par
配送の例では、朝に全ルートを決めるのが静的計画である。昼に新規注文が入ったり、車両が遅れたりしたときに、現在の車両位置と残り注文を見て再計画するのが動的計画である。動的計画では、今の選択が将来の柔軟性を奪わないかも考える必要がある。\par
試験答案では、静的/動的の違いを『計画を変更するか』だけでなく、状態、時点、決定、外生情報、遷移の連鎖として説明する。特に動的問題では、状態が更新され、次の決定の入力になるというループを明確に書くと満点答案に近づく。\par\NoteSection{確認問題と解答}\begin{enumerate}\item \textbf{L04-S051: 「Outlook」の概念を、定義・具体例・モデル上の意味の順に説明せよ。}\\定義としては、dynamismは、時間の進行により状態が変わり、意思決定を繰り返す性質である。。具体例では、配送・在庫・フードデリバリーのような現実問題を用い、状態、決定、外生情報、報酬、または情報モデル/意思決定モデルのどこに対応するかを明示する。モデル上の意味として、この概念が目的関数、制約、状態表現、将来価値、または方策選択にどのような影響を与えるかを書く。試験では、用語だけを列挙せず、入力データから意思決定、さらに現実の実装へつながる流れを文章で示すことが重要である。\item \textbf{L04-S051: このスライドの考え方を、講義全体のDDDMプロセスの中に位置づけよ。}\\DDDMプロセスでは、現実世界のデータを集約して情報モデルを作り、意思決定モデルまたは方策で行動を選び、実装後に結果を観測して次の状態へ進む。このスライドの概念は、そのどこか一部だけではなく、データ表現、意思決定、将来不確実性、計算可能性のいずれかに関わる。答案では、まずこの概念がどの段階に属するかを述べ、次に他の段階へどう影響するかを書く。例えば価値関数なら将来価値を現在決定へ戻す役割、FIFOなら旅行時間情報モデルの整合性を保証する役割、CFAなら目的関数を調整して将来に良い行動を誘導する役割である。\end{enumerate}
\end{minipage}
\end{adjustbox}
\end{minipage}

\clearpage\ExamHeader{Lecture 4 - Dynamism}
\ExamQuestion{WS22/23 Aufgabe 6 (4 点)}
\ExamSubsection{問題内容}
動的意思決定プロセスを記述する量。設問では、状態・決定・外生情報・遷移・報酬/費用・時点をセットで説明すること。 ことが求められる。\par
\ExamSubsection{満点答案}
満点答案では、Real Worldから観測データを集約してInformation Modelを作り、そこからDecision Modelを作り、解をImplementation/Disaggregationで現実の行動へ戻す流れを説明する。Information Modelは現実の圧縮表現であり、Decision Modelは決定変数、目的関数、制約により行動選択を評価する構造である。\par
\ExamSubsection{具体例による解説}
具体例として、配送問題なら、顧客注文・車両位置・旅行時間を情報モデルにし、訪問順や割当を意思決定モデルで決める。在庫問題なら、在庫水準を状態、注文量を決定、需要を外生情報、在庫更新式を遷移、欠品費・保管費を費用として整理する。問題文の文脈に合わせて、抽象用語を必ず具体的な変数や行動に置き換える。\par
\ExamSubsection{採点で落とさない要素}
\begin{itemize}
\item 設問が求める概念の定義を最初に明確に書く。
\item 講義の専門語を列挙するだけでなく、現実例とモデル要素へ対応付ける。
\item 利点だけでなく限界・注意点も最低一つ述べる。
\item 式が出る問題では、記号の意味を文章で説明する。
\item 新しい事例問題では、状態・決定・外生情報・遷移・報酬/費用の順に整理する。
\end{itemize}
\vspace{2mm}\hrule\vspace{2mm}
\ExamQuestion{WS23/24 Aufgabe 8 (7 点)}
\ExamSubsection{問題内容}
動的意思決定問題の構成要素と時点tの状態。設問では、状態が静的情報・動的情報・履歴/リソースなどから構成されることを説明すること。 ことが求められる。\par
\ExamSubsection{満点答案}
満点答案では、Real Worldから観測データを集約してInformation Modelを作り、そこからDecision Modelを作り、解をImplementation/Disaggregationで現実の行動へ戻す流れを説明する。Information Modelは現実の圧縮表現であり、Decision Modelは決定変数、目的関数、制約により行動選択を評価する構造である。\par
\ExamSubsection{具体例による解説}
具体例として、配送問題なら、顧客注文・車両位置・旅行時間を情報モデルにし、訪問順や割当を意思決定モデルで決める。在庫問題なら、在庫水準を状態、注文量を決定、需要を外生情報、在庫更新式を遷移、欠品費・保管費を費用として整理する。問題文の文脈に合わせて、抽象用語を必ず具体的な変数や行動に置き換える。\par
\ExamSubsection{採点で落とさない要素}
\begin{itemize}
\item 設問が求める概念の定義を最初に明確に書く。
\item 講義の専門語を列挙するだけでなく、現実例とモデル要素へ対応付ける。
\item 利点だけでなく限界・注意点も最低一つ述べる。
\item 式が出る問題では、記号の意味を文章で説明する。
\item 新しい事例問題では、状態・決定・外生情報・遷移・報酬/費用の順に整理する。
\end{itemize}
\vspace{2mm}\hrule\vspace{2mm}
\ExamQuestion{SS25 Exercise 1 (2 点)}
\ExamSubsection{問題内容}
a priori optimizationの動的意思決定に対する利点。設問では、静的最適化の単純性・計算容易性などを説明すること。 ことが求められる。\par
\ExamSubsection{満点答案}
満点答案では、Real Worldから観測データを集約してInformation Modelを作り、そこからDecision Modelを作り、解をImplementation/Disaggregationで現実の行動へ戻す流れを説明する。Information Modelは現実の圧縮表現であり、Decision Modelは決定変数、目的関数、制約により行動選択を評価する構造である。\par
\ExamSubsection{具体例による解説}
具体例として、配送問題なら、顧客注文・車両位置・旅行時間を情報モデルにし、訪問順や割当を意思決定モデルで決める。在庫問題なら、在庫水準を状態、注文量を決定、需要を外生情報、在庫更新式を遷移、欠品費・保管費を費用として整理する。問題文の文脈に合わせて、抽象用語を必ず具体的な変数や行動に置き換える。\par
\ExamSubsection{採点で落とさない要素}
\begin{itemize}
\item 設問が求める概念の定義を最初に明確に書く。
\item 講義の専門語を列挙するだけでなく、現実例とモデル要素へ対応付ける。
\item 利点だけでなく限界・注意点も最低一つ述べる。
\item 式が出る問題では、記号の意味を文章で説明する。
\item 新しい事例問題では、状態・決定・外生情報・遷移・報酬/費用の順に整理する。
\end{itemize}
\vspace{2mm}\hrule\vspace{2mm}